\documentclass[aps,pra,twocolumn,superscriptaddress,longbibliography]{revtex4-2}

\usepackage{amsmath,amssymb,graphicx,xcolor,bm,braket,dcolumn,booktabs}
\usepackage{hyperref}

% dcolumn shorthand: enable `d{N.M}` as alias for `D{.}{.}{N.M}`
\newcolumntype{d}[1]{D{.}{.}{#1}}


% ---------------------------------------------------------------
% \gvq{registry-key}{value} -- a load-bearing quantity, carrying a
% foreign key into debug/qa/numeric_registry.py alongside its literal
% value.  Renders as the value alone, so this file and the archived
% PDF stay self-contained; the key gives the QA gate (C20) referential
% integrity, and `--index <key>` names every locus citing the
% quantity when it moves.  See CLAUDE.md SS15.
\newcommand{\gvq}[2]{#2}

\begin{document}

\title{Structurally Sparse Qubit Hamiltonians\\from Spectral Graph Theory}
\thanks{Paper 14 in the Geometric Vacuum series}

\author{J.~Loutey}
\affiliation{Independent Researcher, Kent, Washington}

\date{March 28, 2026}

\begin{abstract}
We show that qubit Hamiltonians derived from spectral graph theory lattices
exhibit structurally favorable Pauli decompositions for quantum simulation.
After Jordan--Wigner transformation, the GeoVac lattice Hamiltonian yields
$\sim Q^{3.8}$ Pauli string terms in the number of qubits~$Q$ (corrected
2026-08-29 to the exact Coulomb selection rule; the retired pair-diagonal
encoding gave $Q^{3.15}$), compared
to $O(Q^{4+})$ for conventional Gaussian-basis Hamiltonians
(our own log-log fit of the published Pauli counts of
Trenev~\textit{et al.}~\cite{trenev2025} gives
$Q^{4.25}$ for LiH and $Q^{3.92}$ for H$_2$O, spanning $Q^{3.9}$--$Q^{4.3}$).  The advantage originates in angular momentum
selection rules---enforced by Gaunt integrals at the two-electron integral
level---which cause the electron repulsion integral (ERI) tensor density
to decay as $\approx M^{-0.49}$ with
the number of spatial orbitals~$M$, while
Gaussian ERI tensors remain near-dense at $O(1)$.
Beyond term counts, we quantify two operationally relevant cost metrics:
the number of qubitwise-commuting (QWC) measurement groups scales as
$O(Q^{4.01})$, and the Pauli 1-norm
$\lambda = \sum_i |c_i|$---which controls both Trotter step counts and
qubitization query complexity---scales as $O(Q^{1.77})$, $R^2 = 0.9997$
(re-measured 2026-08-30 under the exact ERI rule; the retired
pair-diagonal fits gave $O(Q^{3.36})$ and $O(Q^{1.69})$).  The 1-norm
remains strongly sub-quadratic and its equal-qubit advantage survives at
$3.8\times$ and $7.1\times$; the QWC exponent, however, now exceeds the
Pauli term exponent and sits inside the Gaussian $O(Q^{4\text{--}5})$
band, so no measurement-group scaling advantage is claimed.  The sub-quadratic 1-norm scaling is the key result
for fault-tolerant simulation: it implies that Trotter circuit depth
grows far more slowly with system size than term counts alone would
suggest.
All Gaussian baselines are validated against published integral values
(Szabo and Ostlund for H$_2$; analytical STO for He).

For composed multi-center systems, the
pseudopotential approximation of the composed natural geometry
framework~\cite{loutey_paper17} yields block-diagonal electron
repulsion integrals whose Pauli count factorizes exactly into a
basis factor times the qubit count, $N_{\mathrm{Pauli}} =
c(n_{\max})\,Q$, with rational coefficients $c(1) = 3/2$,
$c(2) = 279/10$, $c(3) = 7089/14$ identical across every molecule
tested (eight systems, three periodic-table rows).  Along the
within-molecule basis axis the corresponding log--log exponent is
$3.17$ (fit over $n_{\max} = 1$--$3$; molecule-independent to all
reported digits), with the local slope rising toward $\sim 3.8$ by
$n_{\max} = 4$---comparable to the atomic path and approaching the
lower edge of the Gaussian span.  The decisive advantage is therefore
the \emph{system-size} axis: at fixed basis the count is exactly
linear in~$Q$, against Gaussian baselines that scale as
$Q^{3.9}$--$Q^{4.3}$ (published counts of
Trenev~\textit{et al.}~\cite{trenev2025}, span fitted by us); at equal
qubit counts this yields $54\times$--$317\times$ fewer Pauli terms for
H$_2$O across $n_{\max} = 1$--$3$.  These comparisons are at matched
qubit counts, not matched accuracy; the composed basis carries
classical equilibrium-geometry errors of 5--26\% (LiH/BeH$_2$/H$_2$O;
Sec.~\ref{sec:accuracy}), making the framework optimal for fixed-geometry
quantum simulation rather than geometry optimisation.
Across the composed molecular library (35 molecules spanning
$Z = 1$--56, H through Ba, including the first transition series ScH--ZnH), the composed
Pauli count is exactly linear in the qubit count:
$N_{\mathrm{Pauli}} = \gvq{composed_coeff}{27.90} \times Q$ at $n_{\max} = 2$ (exact, non-identity
terms; $d$-orbital blocks give 30.03).  All counts in this paper are computed
under the exact global-$M_L$ Coulomb selection rule
($m_a + m_b = m_c + m_d$).  Until 2026-08-29 the production encoders realized
a \emph{pair-diagonal} restriction ($m_a=m_c$, $m_b=m_d$) that was documented
as a deliberate sparsifying approximation but traced to a sign error in the
Gaunt coefficient (the 3j magnetic argument $q = m_c - m_a$ instead of
$m_a - m_c$); the corresponding retired coefficients were $11.10$
(main-group) and $9.23$ ($d$-block).  Under the exact rule these become
$\gvq{composed_coeff}{27.90}$ and $\gvq{composed_coeff_d}{30.03}$: the $d$-block coefficient is the \emph{larger} of the
two, so the $d$-block is \emph{denser} than main-group and the apparent
$d$-block sparsity advantage does not survive (Sec.~\ref{sec:eri_rule}).
\end{abstract}

\maketitle

%% ========================================================================
\section{Introduction}
\label{sec:intro}
%% ========================================================================

Practical quantum chemistry on near-term quantum computers is fundamentally
constrained by gate count.  For variational quantum eigensolvers (VQE) and
Hamiltonian simulation via Trotterization, the number of distinct Pauli
string terms in the qubit Hamiltonian directly determines the circuit depth
and measurement overhead~\cite{mcardle2020,cao2019}.  Reducing the Pauli
term count is therefore a central objective for quantum computational
chemistry.

Prior approaches to Pauli count reduction operate as post-hoc optimizations
on an already-constructed qubit Hamiltonian: qubit tapering exploits
symmetries to eliminate qubits~\cite{bravyi2017}; Pauli grouping strategies
reduce the number of distinct measurement circuits by identifying
simultaneously measurable sets~\cite{verteletskyi2020,yen2020}; and
low-rank factorizations of the two-electron integral tensor can compress
the operator representation~\cite{motta2021,lee2021}.  These techniques
are effective but treat the sparsity structure of the underlying
second-quantized Hamiltonian as given.

In this work we show that a different choice of \emph{basis}---the
spectral graph theory lattice of the GeoVac framework---produces
Hamiltonians that are structurally sparse \emph{before} any compilation
optimization.  The GeoVac lattice encodes electronic structure on a
discrete graph whose nodes carry hydrogenic quantum numbers
$(n, l, m)$~\cite{loutey_paper7,loutey_paper13}.  The key structural
consequence is that the two-electron repulsion integrals (ERIs) inherit
angular momentum selection rules from the underlying Gaunt
integrals~\cite{gaunt1929}, which zero out most of the $M^4$ possible
ERI tensor entries.  As the basis grows, the ERI density decays as
$\approx M^{-0.49}$ (corrected 2026-08-29), in stark
contrast to Gaussian bases where the ERI tensor is nearly dense.

The relationship to numerical low-rank factorizations~\cite{motta2021,
lee2021} can be made precise.  The two-electron operator
$V_{ee} = \sum_P L_P \otimes L_P$ that double factorization and tensor
hypercontraction produce is an element of $\mathcal{A} \otimes \mathcal{A}$
in the tensor-product spectral triple, and in the GeoVac basis the same
decomposition is supplied analytically by the multipole expansion
$1/r_{12} = \sum_{L,M} (4\pi/(2L{+}1))\,(r_<^L/r_>^{L{+}1})\,Y_{LM}(1)\,
Y_{LM}^{*}(2)$, with $L_P$ proportional to $Y_{LM}$ and rank fixed by
the count of distinct radial-density channels at each angular momentum.
Numerically, the double factorization of a GeoVac ERI tensor lives
entirely within this multipole subspace: every singular vector of the
reshaped two-electron tensor has machine-precision unit overlap with the
analytic multipole-channel subspace, verified bit-exactly across
LiH/BeH$_2$/H$_2$O sub-blocks at $n_{\max} \in \{2, 3, 4\}$.  At the
production basis ($n_{\max} = 2$, valence set $\{1s, 2s, 2p\}$) the
double-factorization rank equals the analytic count of distinct radial
densities (seven) exactly.  Double factorization on a GeoVac Hamiltonian
is therefore an SVD-sorted basis for the same algebraic subspace the
multipole expansion supplies by construction, not a complementary
compression of an independent structure (see
\texttt{debug/sprint\_df\_multipole\_lift\_memo.md}).

Under the Jordan--Wigner transformation~\cite{jordan1928,bravyi2002},
each nonzero two-body operator
$a^\dagger_p a^\dagger_q a_s a_r$ maps to a bounded number of Pauli
string terms---typically 4 for density-type integrals ($p{=}s$, $q{=}r$)
and 8 for exchange-type integrals (four distinct indices), independent
of~$Q$.  (The JW $Z$-chains increase the \emph{weight} of each string
but not the \emph{number} of distinct strings per integral.)  Similarly,
each one-body term $a^\dagger_p a_q$ generates $O(1)$ Pauli strings.
The total Pauli count therefore scales as
\begin{equation}
  N_{\mathrm{Pauli}} \sim C_1 \cdot n_{\mathrm{h1}}
  + C_2 \cdot n_{\mathrm{ERI}},
  \label{eq:pauli_scaling}
\end{equation}
where $C_1, C_2$ are $O(1)$ constants and $n_{\mathrm{h1}}$,
$n_{\mathrm{ERI}}$ are the numbers of nonzero one- and two-electron
integrals.  For Gaussian bases with
$n_{\mathrm{ERI}} \sim M^4$ and $Q = 2M$, the two-body term dominates at
$O(Q^{4})$.  For the GeoVac lattice with
$n_{\mathrm{ERI}} \sim M^2$, the leading-order prediction is
$O(Q^2)$ (this $n_{\mathrm{ERI}} \sim M^2$ count corresponds to an ERI
density $\sim 1/M^2$, though empirically the density decays more slowly,
as $\sim 1/M$; see Sec.~\ref{sec:eri_density}).  The empirical exponent
of $\alpha \approx 3.8$ (corrected 2026-08-29; retired rule: 3.15) over the
range $Q = 10$--$110$ exceeds this asymptotic estimate, reflecting
spin-channel multiplicity (each spatial ERI generates up to four
spin-orbital ERIs) and subleading combinatorial terms at moderate~$Q$.
The scaling \emph{gap} between GeoVac and Gaussian encodings is
$0.1$--$0.5$ in exponent ($3.8$ against $3.9$--$4.3$).  It is
\emph{not} robust to the per-integral counting: the retired
pair-diagonal rule put the GeoVac exponent at $3.15$ and hence the gap at
$1$--$1.5$, and the correction closed most of it.  What survives as a
structural consequence of selection-rule sparsity is the exact linearity
in~$Q$ across molecules at fixed basis, not the within-molecule exponent
gap.

We validate this scaling with systematic benchmarks on atoms (H, He) and
molecules (H$_2$) using the GeoVac lattice at increasing basis sizes
($n_{\max} = 2$--$5$), comparing against standard Gaussian basis sets
(STO-3G, 6-31G, cc-pVDZ).  All Jordan--Wigner transformations are
performed with OpenFermion~\cite{mcclean2020}.


%% ========================================================================
\section{Methods}
\label{sec:methods}
%% ========================================================================

\subsection{GeoVac lattice Hamiltonians}
\label{sec:geovac}

The GeoVac framework constructs electronic Hamiltonians on discrete graphs
whose nodes correspond to hydrogenic spin-orbitals $(n, l, m, \sigma)$
with $1 \le n \le n_{\max}$, $0 \le l < n$, $-l \le m \le l$, and
$\sigma \in \{\uparrow, \downarrow\}$.  The number of spatial orbitals is
$M = n_{\max}^2$; the number of qubits (spin-orbitals) is $Q = 2M$.

The one-electron Hamiltonian is a sparse graph Laplacian on this lattice,
with edges connecting angular neighbors ($m \leftrightarrow m \pm 1$) and
radial neighbors ($n \leftrightarrow n \pm 1$) within each $(l, m)$
shell~\cite{loutey_paper7}.  In the ``hybrid'' prescription, the diagonal
carries exact hydrogenic eigenvalues $-Z^2/(2n^2)$ while off-diagonal
elements come from the graph Laplacian scaled by the kinetic scale
factor $K = -1/16$.

The two-electron integrals use the full Slater--Condon framework with
Slater $F^k$ (direct) and $G^k$ (exchange) integrals, where the angular
coupling coefficients are Gaunt integrals computed via Wigner $3j$
symbols~\cite{loutey_paper13}:
\begin{equation}
  c^k(l_1, l_2, l_3) = (-1)^{m_1}
  \begin{pmatrix} l_1 & k & l_3 \\ -m_1 & q & m_3 \end{pmatrix}
  \begin{pmatrix} l_1 & k & l_3 \\ 0 & 0 & 0 \end{pmatrix}.
  \label{eq:gaunt}
\end{equation}
The triangle inequality $|l_1 - l_3| \le k \le l_1 + l_3$ and parity
selection rule $l_1 + k + l_3 = \text{even}$ enforce strict sparsity on
the $c^k$ coefficients, which propagates directly into the ERI tensor.

\subsection{Jordan--Wigner transformation}
\label{sec:jw}

Given $Q$ spin-orbitals, the Jordan--Wigner mapping encodes fermionic
creation and annihilation operators as~\cite{jordan1928}
\begin{equation}
  a^\dagger_p = \frac{1}{2}(X_p - iY_p) \prod_{j<p} Z_j,
  \label{eq:jw}
\end{equation}
where $X_p$, $Y_p$, $Z_p$ are Pauli operators on qubit~$p$.  A one-body
term $h_{pq}\, a^\dagger_p a_q$ generates a Pauli string of weight up to
$|p - q| + 1$ (due to the $Z$-chain), while a two-body term
$v_{pqrs}\, a^\dagger_p a^\dagger_q a_s a_r$ generates strings of weight
up to $\max(|p-r|, |q-s|) + 2$.  The total number of distinct Pauli
strings after collecting terms is our primary metric.

\subsection{Comparison methodology}
\label{sec:comparison}

For the Gaussian baseline, we use published integral values from two
sources.  H$_2$ STO-3G integrals ($M = 2$) are from Szabo and
Ostlund~\cite{szabo1996} (Table~3.15, MO basis after restricted
Hartree--Fock); the FCI energy at $R = 1.4$~bohr is
$E_{\mathrm{FCI}} = -1.1373$~Ha.  He STO-3G integrals ($M = 1$) use
the analytical Slater-type orbital formulas with exponent
$\zeta = 1.6875$ from Hehre, Stewart, and Pople~\cite{hehre1969}:
$h_{11} = \zeta^2/2 - Z\zeta$ and $(11|11) = 5\zeta/8$, giving
$E = -2.848$~Ha.  These integrals are exact (not Gaussian fits) and
serve as genuine baselines for the Pauli count, measurement grouping,
and 1-norm comparisons.

For larger Gaussian bases (6-31G, cc-pVDZ), the Pauli term counts
reported in Table~\ref{tab:pauli} use synthetically generated integrals
with representative sparsity patterns, as PySCF was unavailable in our
build environment.  These are retained for scaling-fit purposes only;
the structural conclusion---Gaussian ERIs are near-dense while lattice
ERIs are selection-rule-sparse---is well established~\cite{helgaker2000}.

For GeoVac, we construct \texttt{LatticeIndex} objects with
\texttt{vee\_method='slater\_full'} and \texttt{h1\_method='hybrid'} at
$n_{\max} = 2, 3, 4, 5$, yielding $M = 5, 14, 30, 55$ spatial orbitals
and $Q = 10, 28, 60, 110$ qubits.  The helium atom ($Z = 2$, 2~electrons)
serves as the primary benchmark for scaling analysis.  For the molecular
comparison, we encode H$_2$ at $R = 1.4$~bohr using
\texttt{MolecularLatticeIndex} with $n_{\max} = 2$ per atom.

All Jordan--Wigner transformations use OpenFermion
1.7.1~\cite{mcclean2020}.  Pauli terms with coefficients below $10^{-12}$
are discarded.


%% ========================================================================
\section{Results}
\label{sec:results}
%% ========================================================================

\subsection{Pauli term counts}
\label{sec:pauli_counts}

Table~\ref{tab:pauli} reports the Pauli term counts, qubit counts, and
ERI densities for all systems studied.

\begin{table}[t]
\caption{Pauli term counts after Jordan--Wigner transformation.
$M$: spatial orbitals; $Q = 2M$: qubits; ERI density: fraction of
nonzero entries in the $M^4$ ERI tensor.  Gaussian entries marked
$\dagger$ use synthetic integrals with representative sparsity;
unmarked Gaussian entries use validated computed integrals.}
\label{tab:pauli}
\begin{ruledtabular}
\begin{tabular}{llrrrd{2.1}}
Method   & Basis     & $M$ & $Q$ & Pauli terms & \multicolumn{1}{c}{ERI (\%)} \\
\hline
Gaussian & He STO-3G    &   1 &   2 &           4 & 100.0 \\
Gaussian & H$_2$ STO-3G &   2 &   4 &          15 & 50.0 \\
Gaussian & 6-31G$^\dagger$ &   4 &   8 &         201 & 50.0 \\
Gaussian & He cc-pVDZ &   5 &  10 &         \gvq{he_ccpvdz_pauli}{156} & 100.0 \\
Gaussian & cc-pVDZ$^\dagger$ &  10 &  20 &      23{,}189 & 100.0 \\
Gaussian & He cc-pVTZ &  14 &  28 &      \gvq{he_ccpvtz_pauli}{21{,}607} & 100.0 \\
\hline
GeoVac   & $n_{\max}{=}2$ &   5 &  10 &         \gvq{he_n2_pauli}{288} & 17.1 \\
GeoVac   & $n_{\max}{=}3$ &  14 &  28 &      \gvq{he_n3_pauli}{14{,}079} & 9.9 \\
GeoVac   & $n_{\max}{=}4$ &  30 &  60 &     \gvq{he_n4_pauli}{250{,}403} & 7.1 \\
GeoVac   & $n_{\max}{=}5$ &  55 & 110 & \gvq{he_n5_pauli}{2{,}434{,}442}\textsuperscript{b} & 5.7 \\
\hline
GeoVac   & H$_2$ $n_{\max}{=}2$\textsuperscript{c} & 10 & 20 &   391 & 1.9 \\
\end{tabular}
\end{ruledtabular}
{\footnotesize \textsuperscript{a}\emph{Convention:} the Pauli-term
column includes the identity term throughout, on both the Gaussian and
GeoVac sides.  The non-identity counts are one lower
($287$ / $14{,}078$ / $250{,}402$ / $2{,}434{,}441$), and it is those that
the scaling fits and 1-norm comparisons elsewhere in this paper use.
\textsuperscript{b}Measured 2026-08-30:\ the exact-rule Jordan--Wigner
transform over the corrected $519{,}585$ ERIs is computationally heavy
($\sim\!28$ minutes) but completed; the retired pair-diagonal count was
$227{,}338$.  This cell read ``---'' until then.
\textsuperscript{c}Retired-rule vintage (its generator predates the
2026-08-29 correction).}
\end{table}

The most striking feature is the ERI density column: Gaussian bases
maintain 50--100\% density regardless of basis size, while GeoVac ERI
density drops monotonically from 17.1\% at $M = 5$ to 5.7\% at $M = 55$
(corrected 2026-08-29 to the exact global-$M_L$ rule; the retired
pair-diagonal counts read 10.4\% and 1.3\%).

The genuine computed Gaussian baselines enable head-to-head comparisons
at identical qubit counts, and under the exact rule the small-basis
comparison is \textbf{unfavorable}: at $Q = 10$, GeoVac He
($n_{\max} = 2$) produces \gvq{he_n2_pauli}{288}~Pauli terms versus \gvq{he_ccpvdz_pauli}{156} for Gaussian He
cc-pVDZ---$1.8\times$ \emph{more} (the retired pair-diagonal count, 120,
had claimed a $1.3\times$ reduction).  At $Q = 28$, GeoVac He
($n_{\max} = 3$) produces \gvq{he_n3_pauli}{14{,}079}~terms versus \gvq{he_ccpvtz_pauli}{21{,}607} for Gaussian
He cc-pVTZ---a $1.5\times$ reduction (retired: $8.1\times$).  The
crossover reflects the density gap: Gaussian ERI density stays at
100\% while GeoVac's decays as $\approx M^{-0.49}$ (measured over
$M = 5$--$55$; the retired rule's $M^{-0.85}$--$M^{-1}$ overstated the
decay), so the atomic advantage appears only at larger bases and is
modest there.

At matched qubit count ($Q = 20$), the molecular comparison is
also instructive: GeoVac H$_2$ with $n_{\max} = 2$ per atom
produces 391 Pauli terms at 1.9\% ERI density, while synthetic Gaussian
cc-pVDZ produces 23{,}189 terms at 100\% ERI density---a
\textbf{59$\times$} reduction in Pauli terms.

\subsection{Scaling exponents}
\label{sec:scaling}

\begin{figure}[t]
  \centering
  \includegraphics[width=0.95\columnwidth]{paper_14_figures/pauli_scaling.png}
  \caption{Pauli term count after Jordan--Wigner transformation versus
  number of qubits~$Q$.  Red squares: Gaussian basis sets (STO-3G, 6-31G,
  cc-pVDZ).  Blue circles: GeoVac lattice ($n_{\max} = 2$--$5$).  Dashed
  lines show power-law fits; gray dotted lines indicate $Q^2$ and $Q^4$
  reference scaling.  The GeoVac exponent ($\alpha = 3.78$, corrected
  2026-08-29) sits just below the $Q^{3.9}$--$Q^{4.3}$ Gaussian molecular
  range (specific exponents $\alpha = 3.92$--$4.25$, our log-log fit
  of the published Pauli counts of Trenev~\textit{et al.}~\cite{trenev2025}).}
  \label{fig:pauli_scaling}
\end{figure}

Power-law fits $N_{\mathrm{Pauli}} = a \cdot Q^\alpha$ in log--log space
yield:
\begin{align}
  \alpha_{\mathrm{GeoVac}}   &= 3.78 \quad \text{(3-point log--log fit, } n_{\max} = 2\text{--}4\text{)}. \label{eq:geovac_exp}
\end{align}
For the Gaussian baseline, we take the published molecular Pauli counts
(Jordan--Wigner with 2-qubit reduction) of
Trenev~\textit{et al.}~\cite{trenev2025} (Table~5, Appendix~B) and fit
them in log--log space, giving system-specific exponents: $\alpha = 4.25$
for LiH ($R^2 = 0.999$, 3~basis sets) and $\alpha = 3.92$ for H$_2$O
($R^2 > 0.999$, 3~basis sets), both within the $Q^{3.9}$--$Q^{4.3}$
span of those fits.  These supersede our earlier estimate
of $\alpha_{\mathrm{Gaussian}} \approx 4.6$ from synthetic
integrals~\cite{helgaker2000}.

The GeoVac exponent reflects the $\sim 1/M$ ERI density decay
(measured $\approx M^{-0.49}$, corrected 2026-08-29), which suppresses the
effective two-body contribution well below the $O(Q^4)$ Gaussian
scaling, with additional suppression from the
band-structured one-body graph Laplacian.

The scaling gap is $\Delta\alpha \approx 0.1$--$0.5$ ($3.78$ against
$3.9$--$4.3$, depending on molecule and Gaussian basis hierarchy).
Corrected 2026-08-30:\ the retired pair-diagonal exponent $3.15$ made this
gap $1.0$--$1.5$, and the inference drawn from it---that the GeoVac
advantage grows with system size along the \emph{basis} axis---does not
survive.  What does survive is the exact linearity in~$Q$ across molecules
at fixed basis (Sec.~\ref{sec:composed}), which is a system-size statement
rather than a basis-size one.

\subsection{ERI density decay}
\label{sec:eri_density}

\begin{figure}[t]
  \centering
  \includegraphics[width=0.95\columnwidth]{paper_14_figures/eri_density.png}
  \caption{ERI tensor density versus number of spatial orbitals~$M$.
  Red squares: Gaussian bases (50--100\% dense).  Blue circles: GeoVac
  lattice, showing monotonic decay from 17.1\% ($M = 5$) to 5.7\%
  ($M = 55$) (corrected 2026-08-29; the retired pair-diagonal figures
  read 10.4\% to 1.3\%).  Gray dashed curve: $\sim M^{-0.49}$ fit, normalized to
  the $n_{\max} = 2$ data point.  The GeoVac ERI density decays as
  $\sim M^{-0.49}$ under the exact rule, shallower than
  the $1/M^2$ reference, the decay being imposed by angular momentum
  selection rules.  (The retired pair-diagonal rule gave
  $\approx M^{-0.85}$--$M^{-1}$, which overstated the decay.)}
  \label{fig:eri_density}
\end{figure}

Table~\ref{tab:eri_density} quantifies the ERI density decay for the
GeoVac lattice.

\begin{table}[t]
\caption{ERI tensor statistics for GeoVac He ($Z = 2$, 2~electrons).
$n_{\mathrm{ERI}}$: nonzero integrals; density $= n_{\mathrm{ERI}}/M^4$.}
\label{tab:eri_density}
\begin{ruledtabular}
\begin{tabular}{rrrrd{2.1}}
$n_{\max}$ & $M$ & $n_{\mathrm{ERI}}$ & $M^4$ & \multicolumn{1}{c}{Density (\%)} \\
\hline
2 &  5 &     107 &        625 & 17.1 \\
3 & 14 &   3{,}800 &     38{,}416 & 9.9 \\
4 & 30 &  57{,}700 &    810{,}000 & 7.1 \\
5 & 55 & 519{,}585 &  9{,}150{,}625 & 5.7 \\
\end{tabular}
\end{ruledtabular}
\end{table}

The density scales approximately as $\rho_{\mathrm{ERI}} \propto M^{-1/2}$
(the four tabulated points give an empirical exponent $\approx M^{-0.46}$;
the frequently-quoted $M^{-0.49}$ is instead the two-point
$M = 5 \to 30$ endpoint slope; the retired pair-diagonal tensor gave
$M^{-0.85}$---the deleted $m$-swap couplings were responsible for most of
the apparent decay), which we can partly understand analytically.
The number of nonzero Gaunt
coefficients $c^k(l_1, l_2, l_3)$ is constrained by the triangle
inequality: for each pair $(l_1, l_3)$, only
$\min(2l_1, 2l_3) + 1 \le 2l_{\min} + 1$ values of $k$ contribute.
Since $l$ ranges up to $n_{\max} - 1$, the total number of allowed
$(l_1, l_2, l_3, k)$ quartets grows as $O(n_{\max}^3)$.  With
$M = n_{\max}^2$ spatial orbitals, this gives
$n_{\mathrm{ERI}} \sim O(n_{\max}^3 \cdot n_{\max}^2) = O(M^{5/2})$,
so $\rho_{\mathrm{ERI}} = n_{\mathrm{ERI}} / M^4 \sim M^{-3/2}$.
The empirical decay ($\approx M^{-0.49}$ under the exact rule;
retired: $M^{-0.85}$--$M^{-1}$) is in fact
\emph{shallower} than this $M^{-3/2}$ order-of-magnitude estimate---and
shallower still than the $M^{-2}$ one might naively read off the
endpoints---because spin-channel multiplicity and the growing number of
surviving $(l_1, l_2, l_3, k)$ combinations at larger $n_{\max}$ partly
offset the $M^4$ denominator; the realized density also depends on which
angular selection rule the production code applies, discussed next.

\subsection{The retired pair-diagonal ERI rule: a sign error, not an approximation}
\label{sec:eri_rule}

All Pauli counts, ERI densities, and multipliers in this paper are
realized under the \emph{exact} Coulomb angular selection rule: a two-electron
integral is retained iff $m_a + m_b = m_c + m_d$ (global $M_L$) with nonzero
Gaunt factors, including the \emph{$m$-swap} integrals (e.g.\
$\langle p_{+1}p_{-1}|p_{-1}p_{+1}\rangle$, per-vertex angular factor
$|c^2(1,{\pm}1;1,{\mp}1)| = \sqrt{6}/5 \approx 0.49$).

\emph{Correction (2026-08-29).}  Earlier versions of this paper reported
counts realized under a \emph{pair-diagonal} restriction ($m_a = m_c$ and
$m_b = m_d$), described as a deliberate sparsifying approximation.  An
independent re-derivation traced that restriction to a sign error: the
production Gaunt coefficient set the 3j magnetic argument to
$q = m_c - m_a$, where the Wigner-3j bottom row requires
$q = m_a - m_c$ (so the row sums to zero); with the wrong sign the
coefficient vanishes identically unless $m_a = m_c$, silently deleting
every $m$-changing multipole (88 of the 126 nonzero $c^k$ at
$l \le 2$).  The error was arbitrated against direct four-dimensional
quadrature of the angular factor from its definition, and the corrected
evaluators agree with the framework's independent precision-physics path
to machine precision.  Three quantitative claims of the earlier disclosure
also fell.  (i)~The re-pricing is \emph{not} a constant factor: it is
$2.51\times$ uniformly across main-group molecules at $n_{\max} = 2$, but
grows along the basis axis ($1.00\times / 2.51\times / 5.40\times$ at
$n_{\max} = 1/2/3$), so the within-molecule exponent moves
$2.54 \to 3.17$ (fit over $n_{\max} \le 3$), with the local slope rising
toward $\sim 3.8$ by $n_{\max} = 4$.  (ii)~The retired rule's low
$d$-block coefficient ($9.23$ vs.\ $11.10$) reverses, as the earlier
disclosure anticipated: measured ScH/TiH give $30.03$ vs.\ the main-group
$27.90$---$d$-blocks are \emph{denser} under the exact rule.  (iii)~The
apparent H$_2$ scaling anomaly (exponent $3.13$ vs.\ the molecules'
$2.50$--$2.52$) dissolves: under the exact rule every system tested gives
the \emph{same} exponent, $3.17$, because the count factorizes exactly as
$N_{\mathrm{Pauli}} = c(n_{\max})\,Q$ with molecule-independent rational
coefficients $c(1) = 3/2$, $c(2) = 279/10$, $c(3) = 7089/14$.  What
survives unchanged is the exact linearity in~$Q$ at fixed basis---the
system-size axis---which is the load-bearing comparison against the
Gaussian $Q^{3.9\text{--}4.3}$ baselines.  The realized rule is pinned by
\texttt{tests/test\_paper14\_eri\_rule.py}; the full defect record is in
the project changelog (v5.1.12).

\subsection{Recoupling sparsity in H-set coupled bases}
\label{sec:recoupling_sparsity}

When $N$-electron angular momenta are organized into pairwise-coupled
(H-set) hyperspherical harmonics---where electrons are first grouped into
pairs $(1,2)$ and $(3,4)$ with definite pair angular momenta $l_{12}$ and
$l_{34}$, then coupled to total $L$---the $6j$ recoupling coefficients
\begin{equation}
  \begin{Bmatrix} l_1 & l_2 & l_{12} \\ l_{12}' & l_1' & k \end{Bmatrix}
\end{equation}
produce additional vanishing ERI matrix elements beyond those enforced by
bare Gaunt selection rules.  The mechanism is structural: the intra-pair
electron repulsion $V_{ee}(1,2)$ preserves $l_{34}$ as a good quantum
number, making the ERI tensor block-diagonal in the partner pair's angular
momentum.  Cross-pair terms $V_{ee}(1,3)$ require a $6j$ recoupling from
the H-set to the K-set (where electrons $(1,3)$ are paired), and the
triangle inequalities on the $6j$ symbol impose selection rules that are
strictly more restrictive than the Gaunt constraints alone.

Table~\ref{tab:recoupling_eri} quantifies the additional sparsity.

\begin{table}[t]
\caption{ERI density in uncoupled vs.\ H-set coupled bases for 4-electron
systems.  The H-set coupling produces 18--26\% fewer nonzero ERI elements
from $6j$ recoupling zeros.}
\label{tab:recoupling_eri}
\begin{ruledtabular}
\begin{tabular}{rrd{2.1}d{2.1}r}
$l_{\max}$ & $M$ & \multicolumn{1}{c}{Uncoupled (\%)} &
  \multicolumn{1}{c}{H-set (\%)} & Reduction \\
\hline
1 & 4  & 12.5 & 9.2  & 26\% \\
2 & 9  & 5.1  & 4.2  & 18\% \\
\end{tabular}
\end{ruledtabular}
\end{table}

The block-diagonal structure is a consequence of angular momentum
conservation and recoupling algebra: it cannot be destroyed by any
continuous deformation of the hyperangular coordinates, and it applies to
\emph{any} encoding that uses angular momentum pair coupling---not
specifically to the GeoVac basis.


\subsection{Energy accuracy}
\label{sec:accuracy}

To confirm that the Pauli count comparison is between methods of
comparable chemical accuracy, Table~\ref{tab:accuracy} reports ground
state energies.

\begin{table}[t]
\caption{Ground state energies (Hartree) for He.  Exact:
$E_{\mathrm{He}} = -2.9034$~Ha~\cite{nist_he}.  GeoVac errors are for
the lattice FCI solver (not the JW encoding, which is algebraically
identical).}
\label{tab:accuracy}
\begin{ruledtabular}
\begin{tabular}{lrd{4.4}d{2.2}}
Method   & $Q$ & \multicolumn{1}{c}{$E_0$ (Ha)} & \multicolumn{1}{c}{Error (\%)} \\
\hline
GeoVac $n_{\max}{=}2$ &  10 & -2.8876 & 0.55 \\
GeoVac $n_{\max}{=}3$ &  28 & -2.8920 & 0.39 \\
GeoVac $n_{\max}{=}4$ &  60 & -2.8961 & 0.25 \\
GeoVac $n_{\max}{=}5$ & 110 & -2.8982 & 0.18 \\
\end{tabular}
\end{ruledtabular}
\end{table}

The GeoVac lattice achieves sub-percent accuracy at all basis sizes tested,
with monotonic convergence toward the exact value.  At $n_{\max} = 5$
($Q = 110$), the error is 0.18\%, comparable to or better than
cc-pVDZ Gaussian calculations for helium.  The accuracy is therefore
competitive, but the term count is not:\ at matched qubits the lattice
carries \emph{more} Pauli terms than cc-pVDZ at $Q=10$, so no favorable
accuracy-per-Pauli-term ratio is claimed under the exact rule.


\subsection{Measurement cost: QWC grouping}
\label{sec:measurement}

On near-term quantum hardware, the energy expectation value
$\langle H \rangle = \sum_i c_i \langle P_i \rangle$ is estimated by
measuring each Pauli string~$P_i$ independently.  Pauli strings that
qubitwise-commute (QWC)---i.e., share the same Pauli operator on every
qubit where both act non-trivially---can be measured simultaneously in
a single circuit~\cite{verteletskyi2020}.  The number of QWC groups
therefore determines the minimum number of distinct measurement circuits,
and reducing it lowers the total shot budget.

We partition each qubit Hamiltonian into QWC groups using greedy
sorted insertion (largest $|c_i|$ first), following the algorithm of
Verteletskyi, Yen, and Izmaylov~\cite{verteletskyi2020}.
Table~\ref{tab:qwc} reports the results.

\begin{table}[t]
\caption{Qubitwise-commuting (QWC) measurement groups for GeoVac He
($Z = 2$, 2~electrons) and Gaussian baselines.  The groups/terms ratio
measures how efficiently terms pack into simultaneous-measurement sets.
QWC grouping for $n_{\max} = 5$ was omitted ($\gvq{he_n5_pauli}{2{,}434{,}441}$ terms
under the exact rule; the $O(N^2)$ grouping algorithm is prohibitively slow
at this scale, and the correction made it more so---the retired
pair-diagonal count was 227{,}338).
The $n_{\max} = 2$ and $3$ rows were re-measured 2026-08-30 under the exact
rule (retired pair-diagonal values:\ 120/25/0.208 and 2{,}659/791/0.297).
The $n_{\max}=4$ row was re-measured as well (retired:\
31{,}039/10{,}199/0.329); its grouping pass is expensive but completed.
The groups/terms ratio rises at all three points, so the exact rule makes
measurement grouping \emph{less} efficient, not more.}
\label{tab:qwc}
\begin{ruledtabular}
\begin{tabular}{lrrrr}
System & $Q$ & Pauli terms & QWC groups & Groups/terms \\
\hline
GeoVac $n_{\max}{=}2$ &  10 &         287 &      67 & 0.233 \\
GeoVac $n_{\max}{=}3$ &  28 &      14{,}078 &  5{,}569 & 0.396 \\
GeoVac $n_{\max}{=}4$ &  60 &     250{,}402 & 86{,}224 & 0.344 \\
\hline
Gauss He STO-3G   &   2 &           4 &       1 & 0.250 \\
Gauss H$_2$ STO-3G &  4 &          15 &       5 & 0.333 \\
\end{tabular}
\end{ruledtabular}
\end{table}

A power-law fit over the GeoVac He data ($n_{\max} = 2$--$4$) gives
\begin{equation}
  N_{\mathrm{QWC}} \sim Q^{4.01}, \quad R^2 = 0.998,
  \label{eq:qwc_scaling}
\end{equation}
slightly \emph{above} the Pauli term exponent of $3.77$, and inside
rather than below the Gaussian $O(Q^{4\text{--}5})$ band.  Under the
retired pair-diagonal rule this exponent read $3.36$---below the term
exponent and well below the Gaussian regime---so the measurement-grouping
advantage reported there does not survive the correction.  The groups-per-term ratio
is higher at every basis than under the retired rule:\
$0.233 / 0.396 / 0.344$ at $n_{\max}=2,3,4$ against the retired
$0.208 / 0.297 / 0.329$.  The sequence is not monotone in the basis, so no
trend is claimed; what the correction establishes is that grouping is
uniformly \emph{less} efficient than previously reported, and the earlier
reading---that selection-rule sparsity translates directly into measurement
savings---does not survive it.

Note that the greedy QWC grouping is a heuristic upper bound; the
optimal minimum clique cover is NP-hard.  More sophisticated grouping
strategies (e.g., general commutativity or unitary
partitioning~\cite{yen2020}) would further reduce the group count, and
these improvements compose with the structural sparsity reported here.


\subsection{Simulation cost: Pauli 1-norm and Trotter bounds}
\label{sec:onenorm}

The Pauli 1-norm,
\begin{equation}
  \lambda = \sum_i |c_i|,
  \label{eq:onenorm}
\end{equation}
is the operationally most important cost metric for quantum
simulation~\cite{childs2021,campbell2019}.  For Trotterization, the
first-order error is bounded by
$\epsilon \le (t/r)^2 \lambda^2 / 2$, giving a minimum step count
\begin{equation}
  r \ge \frac{t\lambda}{\sqrt{2\epsilon}}.
  \label{eq:trotter_steps}
\end{equation}
For qubitization (linear combination of unitaries), the query complexity
scales as $O(\lambda t / \epsilon)$~\cite{childs2021}.  In both
paradigms, smaller~$\lambda$ means lower cost.

Table~\ref{tab:onenorm} reports $\lambda$, the normalized
$\lambda/Q$, and Trotter step counts at two precision targets:
$\epsilon = 10^{-3}$ (NISQ-era) and $\epsilon = 10^{-6}$
(fault-tolerant) for unit simulation time $t = 1$.

\begin{table}[t]
\caption{Pauli 1-norm and first-order Trotter step counts for GeoVac He
($Z = 2$) and Gaussian baselines.  $\lambda$: Pauli 1-norm;
$r(\epsilon)$: minimum Trotter steps for error~$\le \epsilon$ at
$t = 1$~a.u.  The 1-norm bound is a standard upper bound
(Childs et al.~\cite{childs2021}, Eq.~48); the tighter commutator-based
bound is reported in Table~\ref{tab:comm_bound}.}
\label{tab:onenorm}
\begin{ruledtabular}
\begin{tabular}{lrd{3.2}d{2.3}rr}
System & $Q$ & \multicolumn{1}{c}{$\lambda$} & \multicolumn{1}{c}{$\lambda/Q$}
 & $r(10^{-3})$ & $r(10^{-6})$ \\
\hline
GeoVac $n_{\max}{=}2$ &  10 &   11.18 & 1.118 &       253 &     7{,}987 \\
GeoVac $n_{\max}{=}3$ &  28 &   74.21 & 2.650 &     1{,}753 &    55{,}410 \\
GeoVac $n_{\max}{=}4$ &  60 &  275.72 & 4.595 &     5{,}849 &   184{,}959 \\
GeoVac $n_{\max}{=}5$ & 110 &  790.01 & 7.182 &    14{,}693 &   464{,}616 \\
\hline
Gauss He STO-3G   &   2 &    3.38 & 1.688 &        76 &     2{,}387 \\
Gauss H$_2$ STO-3G &  4 &    1.98 & 0.496 &        45 &     1{,}403 \\
Gauss He cc-pVDZ  &  10 &   42.95 & 4.295 &       961 &    30{,}380 \\
Gauss He cc-pVTZ  &  28 &  530.47 & 18.945 &    11{,}862 &   375{,}109 \\
\end{tabular}
\end{ruledtabular}
\end{table}

Power-law fits over the GeoVac He series yield the central quantitative
result of this work:
\begin{align}
  \lambda       &\sim Q^{1.77}, \quad R^2 = 0.9997, \label{eq:onenorm_exp} \\
  r(10^{-3})    &\sim Q^{1.77}, \quad R^2 = 0.9997, \label{eq:trotter3_exp} \\
  r(10^{-6})    &\sim Q^{1.77}, \quad R^2 = 0.9997. \label{eq:trotter6_exp}
\end{align}
The identical exponents for $\lambda$ and~$r$ follow from
Eq.~\eqref{eq:trotter_steps}: for fixed~$t$ and~$\epsilon$,
$r \propto \lambda$.

\emph{Re-measured 2026-08-30.}  The exponents in
Eqs.~\eqref{eq:onenorm_exp}--\eqref{eq:trotter6_exp} and the QWC exponent of
Eq.~\eqref{eq:qwc_scaling} were originally fitted under the retired
pair-diagonal ERI rule.  Re-fitting the same He series under the exact rule gives
\begin{equation}
  N_{\mathrm{Pauli}} \sim Q^{3.77},\qquad
  \lambda \sim Q^{1.77},\qquad
  N_{\mathrm{QWC}} \sim Q^{4.01},
  \label{eq:exact_exponents}
\end{equation}
with $R^2 = 1.0000$, $0.9997$ and $0.9976$ (maximum $|\log|$ residuals
$0.0064$, $0.042$, $0.202$), against the retired $3.15$, $1.69$ and $3.36$.
The first two are four-point fits over $n_{\max} = 2$--$5$
($Q = 10$--$110$); the QWC exponent has three points, because greedy
grouping is $O(N^2)$ and the $n_{\max}=5$ system carries $2.4$~million
Pauli terms.
Two points are worth stating plainly.  The Pauli exponent $3.78$ was
obtained here independently of the value quoted elsewhere in this paper and
agrees with it.  The QWC exponent, by contrast, moves \emph{above} the term
exponent and into the Gaussian $O(Q^{4\text{--}5})$ band, so the
measurement-group scaling advantage claimed under the retired rule does not
survive.

The fourth point confirms rather than moves the fit:\ on three points
($Q \le 60$) the same series gives $3.779$ and $1.792$, so adding
$Q = 110$ shifts the exponents by $0.006$ and $0.018$.  The
$\lambda$-derived Trotter-step columns of Table~\ref{tab:onenorm} remain
retired-rule values, since they were computed from the retired fit rather
than re-derived.

The genuine Gaussian baselines now enable equal-qubit 1-norm
comparisons (Table~\ref{tab:onenorm}).  At $Q = 10$, GeoVac
($\lambda = \gvq{he_n2_lambda}{11.18}$) has $3.8\times$ lower 1-norm than Gaussian
cc-pVDZ ($\lambda = 42.95$), requiring $3.8\times$ fewer Trotter
steps.  At $Q = 28$, GeoVac ($\lambda = \gvq{he_n3_lambda}{74.21}$) has $7.1\times$
lower 1-norm than Gaussian cc-pVTZ ($\lambda = 530.47$).  This
gap grows because the Gaussian 1-norm per qubit ($\lambda/Q$)
increases from 4.3 to 18.9, while GeoVac's grows more slowly (1.12 to
2.65).  (GeoVac 1-norms re-measured 2026-08-30 under the exact rule;
the retired pair-diagonal values were 11.29 and 78.36, giving
$6.8\times$ at $Q=28$.  The 1-norm advantage is one of the few
comparisons the correction leaves intact, and at $Q=28$ it improves.)

For comparison, a dense Hamiltonian with $O(Q^4)$ nonzero ERIs of
comparable magnitude would have $\lambda \sim Q^{2+}$, giving
$Q^{0.3+}$ more Trotter steps at matched system size.  At
$Q = 110$, the GeoVac 1-norm of $\lambda = \gvq{he_n5_lambda}{790}$ is already modest
enough for practical Trotter simulation with $r \approx 17{,}700$
steps at $\epsilon = 10^{-3}$.

As a control, the hydrogen atom (1~electron, no ERI) exhibits a
lower 1-norm exponent of~1.57, confirming that the two-electron
integrals contribute to---but do not dominate---the 1-norm growth.

\paragraph{Commutator-based bound.}
The $\lambda^2$ bound used above counts \emph{all} pairs of Pauli
terms, but two Pauli strings $P_j$ and $P_k$ that commute contribute
zero to the first-order Trotter error.  A tighter bound retains only
anticommuting pairs~\cite{childs2021}:
\begin{equation}
  \epsilon_{\mathrm{comm}} \le \frac{t^2}{r}
    \sum_{\substack{j < k \\ [P_j, P_k] \neq 0}} |c_j|\,|c_k|.
  \label{eq:comm_bound}
\end{equation}
Whether two $Q$-qubit Pauli strings commute or anticommute is
determined by their symplectic inner product: each Pauli operator is
encoded as a pair of $Q$-bit binary vectors $(x, z)$ via
$I \mapsto (0,0)$, $X \mapsto (1,0)$, $Y \mapsto (1,1)$,
$Z \mapsto (0,1)$, and the strings anticommute iff
\begin{equation}
  \sum_{q=1}^{Q} \bigl[ x_j^{(q)} z_k^{(q)} + z_j^{(q)} x_k^{(q)} \bigr]
  \equiv 1 \pmod{2}.
  \label{eq:symplectic}
\end{equation}
For $N$ Pauli terms, the full anticommutativity matrix is computed via
modular matrix arithmetic on two $N \times Q$ binary arrays at cost
$O(N^2 Q)$---no Hilbert-space matrices are constructed.  This is
feasible up to $N \sim 30{,}000$ ($n_{\max} = 4$).

Table~\ref{tab:comm_bound} reports the commutator-based bound
alongside the standard $\lambda^2$ bound for GeoVac He at
$n_{\max} = 2$--$4$.

\begin{table}[t]
\caption{Commutator-based Trotter bound for GeoVac He ($Z = 2$).
$f_{\mathrm{anti}}$: fraction of Pauli term pairs that anticommute;
$r_{\mathrm{comm}}$ and $r_{\lambda^2}$: minimum Trotter steps at
$\epsilon = 10^{-3}$, $t = 1$;
$\rho$: tightening ratio
$\epsilon_{\mathrm{comm}} / \epsilon_{\lambda^2}$.
Three data points ($n_{\max} = 2$--$4$).}
\label{tab:comm_bound}
\begin{ruledtabular}
\begin{tabular}{lrd{2.1}rrr}
System & $Q$ & \multicolumn{1}{c}{$f_{\mathrm{anti}}$ (\%)}
 & $r_{\mathrm{comm}}$ & $r_{\lambda^2}$ & \multicolumn{1}{c}{$\rho$} \\
\hline
$n_{\max}{=}2$ &  10 & 28.0 &        60 &      253 & 0.056 \\
$n_{\max}{=}3$ &  28 & 21.1 &       303 &    1{,}753 & 0.030 \\
$n_{\max}{=}4$ &  60 & 12.2 &       821 &    5{,}849 & 0.020 \\
\end{tabular}
\end{ruledtabular}
\end{table}

A power-law fit over the three data points yields
\begin{equation}
  r_{\mathrm{comm}}(10^{-3}) \sim Q^{1.47},
  \label{eq:comm_exp}
\end{equation}
compared to $r_{\lambda^2} \sim Q^{1.77}$
(Eq.~\ref{eq:trotter3_exp}).  The commutator bound reduces the
effective scaling exponent by~0.30 ($1.77 \to 1.47$; the retired $\lambda$
exponent $1.69$ gave $0.22$).

The anticommuting fraction $f_{\mathrm{anti}}$ decreases as
$Q^{-0.45}$ over the range tested: from 28\% at $Q = 10$ to 12\% at
$Q = 60$.  This reflects the physical mechanism behind the tighter
bound: angular momentum selection rules (Gaunt integral triangle
inequality) cause an increasing fraction of Pauli term pairs to act
on disjoint qubit subsets as the system grows, so their commutators
vanish.

At $n_{\max} = 4$ ($Q = 60$), the commutator bound requires
$r = 821$ Trotter steps versus $r = 5{,}849$ from the $\lambda^2$
bound---a $7\times$ reduction.  Because $\alpha_{\mathrm{comm}} <
\alpha_{\lambda^2}$, this advantage grows with system size.

Table~\ref{tab:scaling_summary} collects all scaling exponents.

\begin{table}[t]
\caption{Power-law scaling exponents for GeoVac He ($Z = 2$, 2~electrons)
over $n_{\max} = 2$--$5$ ($Q = 10$--$110$).  Each metric is fit as
$\mathrm{metric} = a \cdot Q^\alpha$ in log--log space.
\emph{Re-measured 2026-08-30} under the exact ERI rule; the retired
pair-diagonal values were $3.147$, $3.355$ and $1.694$ (Pauli, QWC,
$\lambda$).  The Pauli and $\lambda$ fits use four points; the QWC fit uses
three ($\ddagger$), since greedy grouping is $O(N^2)$ and the $n_{\max}=5$
system carries $2.4$~million Pauli terms.  The commutator bound ($\dagger$)
is unchanged and remains a retired-rule value.}
\label{tab:scaling_summary}
\begin{ruledtabular}
\begin{tabular}{ld{2.3}d{2.4}}
Metric & \multicolumn{1}{c}{$\alpha$} & \multicolumn{1}{c}{$R^2$} \\
\hline
Pauli terms      & 3.773 & 1.0000 \\
QWC groups$^{\ddagger}$ & \gvq{exp_qwc_3pt}{4.013} & 0.9976 \\
1-norm $\lambda$  & \gvq{exp_lambda_4pt}{1.774} & 0.9997 \\
Trotter $r(10^{-3})$ & 1.774 & 0.9997 \\
Trotter $r(10^{-6})$ & 1.774 & 0.9997 \\
Commutator $r_{\mathrm{comm}}(10^{-3})$$^\dagger$ & 1.47 & \\
\end{tabular}
\end{ruledtabular}
\end{table}


\subsection{Gaussian baselines from published integrals}
\label{sec:genuine_gaussian}

The Gaussian baselines in Tables~\ref{tab:qwc} and~\ref{tab:onenorm}
use genuine published integral values, replacing the synthetic
integrals used in earlier versions of Table~\ref{tab:pauli}.
Specifically:
\begin{itemize}
\item \textbf{H$_2$ STO-3G} ($M = 2$, $Q = 4$): MO integrals from
  Szabo and Ostlund~\cite{szabo1996}, Table~3.15.  FCI energy
  $-1.1373$~Ha, validated by diagonalizing the Jordan--Wigner
  Hamiltonian to $< 10^{-4}$~Ha agreement.  Produces exactly 15~Pauli
  terms, 5~QWC groups, $\lambda = 1.98$.
\item \textbf{He STO-3G} ($M = 1$, $Q = 2$): analytical STO integrals
  with $\zeta = 1.6875$~\cite{hehre1969}.  With a single spatial
  orbital, this is a trivial (single-determinant) system; it serves
  as a lower bound on Gaussian costs.  Produces 4~Pauli terms,
  1~QWC group, $\lambda = 3.38$.
\item \textbf{He cc-pVDZ} ($M = 5$, $Q = 10$): contracted $[2s1p]$
  basis with computed integrals from uncontracted Gaussians.  FCI
  energy $-2.8877$~Ha (within 0.001~Ha of published
  values~\cite{nist_he}).  Produces 156~Pauli terms,
  $\lambda = 42.95$.  This provides the first equal-qubit comparison
  with GeoVac $n_{\max} = 2$ (also $Q = 10$): GeoVac has
  $1.8\times$ \emph{more} Pauli terms (\gvq{he_n2_pauli}{288} vs.\ 156) and $3.8\times$
  lower 1-norm.  (Corrected 2026-08-30: the retired pair-diagonal count,
  120, had claimed a $1.3\times$ \emph{reduction}; the exact rule reverses
  the direction of this comparison.  See Sec.~\ref{sec:eri_rule}.)
\item \textbf{He cc-pVTZ} ($M = 14$, $Q = 28$): contracted
  $[3s2p1d]$ basis with cached computed integrals.  FCI energy
  $-2.9003$~Ha (within 0.0002~Ha of published values).  Produces
  21{,}607~Pauli terms at 100\% ERI density, $\lambda = 530.47$.
  At the same qubit count, GeoVac $n_{\max} = 3$ produces
  \gvq{he_n3_pauli}{14{,}078}~terms ($1.5\times$ fewer) with $\lambda = \gvq{he_n3_lambda}{74.21}$
  ($7.1\times$ lower).  (Corrected 2026-08-30: the retired
  pair-diagonal rule gave 2{,}659~terms, $8.1\times$ fewer, and
  $\lambda = 78.36$ at $6.8\times$.  The term-count advantage shrinks
  sharply; the 1-norm advantage improves.)
\end{itemize}

For molecular systems beyond minimal basis, we take the published
Gaussian Pauli counts (Jordan--Wigner with 2-qubit reduction,
exploiting particle number and spin symmetry) for LiH
and H$_2$O across three Gaussian basis sets from
Trenev~\textit{et al.}~\cite{trenev2025} (Table~5, Appendix~B).
These data are reported in Table~\ref{tab:gaussian_published}.

\begin{table}[t]
\caption{Published Gaussian Pauli term counts from
Trenev~\textit{et al.}~\cite{trenev2025} (Table~5, Appendix~B),
Jordan--Wigner encoding with 2-qubit reduction.  Fitted exponents are
our own log-log fit using all three basis sets per molecule.}
\label{tab:gaussian_published}
\begin{ruledtabular}
\begin{tabular}{llrrr}
Molecule & Basis & $Q$ & Pauli terms & $R^2$ (fit) \\
\hline
LiH & STO-3G   & 10 &       276 & \\
LiH & 6-31G    & 20 &     5\,851 & \\
LiH & cc-pVDZ  & 36 &    \gvq{lih_ccpvdz_pauli}{63\,519} & \\
\multicolumn{2}{l}{\textit{Fitted:} $\alpha = 4.25$} & & & 0.999 \\
\hline
H$_2$O & STO-3G   & 12 &       551 & \\
H$_2$O & 6-31G    & 24 &     8\,921 & \\
H$_2$O & cc-pVDZ  & 46 &   \gvq{h2o_ccpvdz_pauli}{107\,382} & \\
\multicolumn{2}{l}{\textit{Fitted:} $\alpha = 3.92$} & & & $>0.999$ \\
\end{tabular}
\end{ruledtabular}
\end{table}

\noindent
\textbf{Note on qubit reduction:} The published Gaussian counts of
Trenev~\textit{et al.}~\cite{trenev2025} (Table~5, Appendix~B) use a
2-qubit reduction (exploiting particle number and spin symmetry),
giving $Q = 10$ for LiH STO-3G rather than the raw Jordan--Wigner count
of $Q = 12$.  The GeoVac values throughout this paper use raw
Jordan--Wigner encoding without this reduction, making the comparison
conservative (GeoVac uses more qubits per orbital than the Gaussian
baseline).

These recomputed molecular exponents ($\alpha = 4.25$ for LiH,
$\alpha = 3.92$ for H$_2$O; both within the $Q^{3.9}$--$Q^{4.3}$ span
of our recomputation) supersede the earlier estimate of
$\alpha \approx 4.6$ obtained from synthetic integrals.  Both remain
well above the GeoVac single-geometry exponent of~3.8 and the
composed-system within-molecule exponent of~3.17 (small-basis fit; the
local slope steepens toward $\sim\!3.8$ by $n_{\max}=4$).  Corrected
2026-08-29; the retired pair-diagonal rule gave~2.5
(Sec.~\ref{sec:composed}).


%% ========================================================================
\section{Composed Systems}
\label{sec:composed}
%% ========================================================================

The structural sparsity demonstrated for single-geometry systems
(Secs.~\ref{sec:results}--\ref{sec:genuine_gaussian}) extends naturally
to multi-center molecules via the composed natural geometry
framework~\cite{loutey_paper17}.  We demonstrate this for four
systems of increasing complexity: H$_2$ (1~orbital block, the
single-block limit), LiH (3~blocks), BeH$_2$ (5~blocks), and H$_2$O
(7~blocks), using the pseudopotential approximation to decompose each
many-electron problem into coupled single-center blocks.


\subsection{Composition principle}
\label{sec:composition_principle}

The composed Hamiltonian partitions the orbital space into independent
single-center blocks, each retaining the angular quantum numbers and
Gaunt integral selection rules of the single-geometry case.  The
pseudopotential approximation (Ref.~\cite{loutey_paper17}, Sec.~II)
replaces explicit core-valence two-body Coulomb interaction with
one-body effective-$Z$ screening plus a Phillips--Kleinman repulsive
barrier, making the ERI tensor block-diagonal: cross-block ERIs are
exactly zero.

The block decomposition for each molecule is:
\begin{itemize}
\item \textbf{H$_2$} (1~block): Single bond-pair block at $Z_{\mathrm{eff}} = 1$
  (2~electrons, no core, no Phillips--Kleinman barrier).  This is the
  single-block limit of the composed framework, structurally identical
  to the He encoding but at $Z = 1$ with R-dependent nuclear repulsion
  $V_{NN} = 1/R$.
\item \textbf{LiH} (3~blocks): Li core ($Z = 3$, 2~electrons) +
  Li-valence ($Z_{\mathrm{eff}} = 1$) + H-valence ($Z = 1$).
\item \textbf{BeH$_2$} (5~blocks): Be core ($Z = 4$, 2~electrons) +
  2~Be--H bond blocks, each containing Be-side
  ($Z_{\mathrm{eff}} = 2$) and H-side ($Z = 1$) orbitals.
\item \textbf{H$_2$O} (7~blocks): O core ($Z = 8$, 2~electrons) +
  2~O--H bond blocks (O-side $Z_{\mathrm{eff}} = 6$ + H-side
  $Z = 1$) + 2~lone pair blocks ($Z_{\mathrm{eff}} = 6$).
\end{itemize}
Each block is constructed from hydrogenic orbitals up to
$n_{\max}$, giving $M = n_{\max}^2$ spatial orbitals per block and
$Q = 2 \sum_b M_b$ total qubits.


\subsection{Scaling results}
\label{sec:composed_scaling}

Table~\ref{tab:composed_pauli} reports the Pauli term counts for all
three molecules at $n_{\max} = 1, 2, 3$.

\begin{table}[t]
\caption{Pauli term counts for composed multi-center systems after
Jordan--Wigner transformation.  $N_b$: number of orbital blocks;
$M$: total spatial orbitals; $Q = 2M$: qubits.  Counts include the
identity operator (add~$-1$ for non-identity count used in
Eq.~\ref{eq:linear_pauli}).  Power-law fits use
all three data points per molecule.  All counts under the exact
global-$M_L$ rule (corrected 2026-08-29); the fitted exponent is 3.17
for \emph{every} system --- the count factorizes exactly as
$N_{\mathrm{Pauli}} = c(n_{\max})\,Q$, so all sweeps share one log--log
shape --- and the local slope rises with basis (LiH $n_{\max}=4$:
$Q=180$, $755{,}867$ terms; slope $3.77$ on the last interval), so 3.17
is a small-basis fit, not an asymptote.}
\label{tab:composed_pauli}
\begin{ruledtabular}
\begin{tabular}{llrrrrd{2.3}}
Molecule & $N_b$ & $n_{\max}$ & $M$ & $Q$ & Pauli terms
  & \multicolumn{1}{c}{$\alpha$ (fit)} \\
\hline
H$_2$\textsuperscript{a} & 1 & 1 &  1 &   2 &         4 & \\
      &   & 2 &  5 &  10 &       280 & \\
      &   & 3 & 14 &  28 &    14\,179 & 3.17 \\
\hline
LiH   & 3 & 1 &  3 &   6 &        10 & \\
      &   & 2 & 15 &  30 &       838 & \\
      &   & 3 & 42 &  84 &    42\,535 & 3.17 \\
\hline
BeH$_2$ & 5 & 1 &  5 &  10 &        16 & \\
      &   & 2 & 25 &  50 &     1\,396 & \\
      &   & 3 & 70 & 140 &    70\,891 & 3.17 \\
\hline
H$_2$O & 7 & 1 &  7 &  14 &        22 & \\
      &   & 2 & 35 &  70 &     1\,954 & \\
      &   & 3 & 98 & 196 &    99\,247 & 3.17 \\
\end{tabular}
\end{ruledtabular}
\end{table}

{\footnotesize \textsuperscript{a}H$_2$ is the single-block limit of
the composed framework: no core, no PK, one bond-pair block at
$Z_{\mathrm{eff}} = 1$.  Under the retired pair-diagonal rule it showed
an anomalous $\alpha = 3.13$ against the molecules' $\approx 2.5$; the
exact rule dissolves the anomaly (all systems 3.17; see below).
The Pauli count is R-independent (verified at $R = 0.5$--$3.0$~bohr).}

The central result is that \emph{every} composed system---from the
single-block H$_2$ to the 7-block H$_2$O---exhibits the \emph{same}
Pauli scaling exponent to all reported digits:
\begin{align}
  \alpha_{\mathrm{H}_2} = \alpha_{\mathrm{LiH}} = \alpha_{\mathrm{BeH}_2}
    = \alpha_{\mathrm{H}_2\mathrm{O}} &= 3.17
  \label{eq:alpha_lih}
\end{align}
(fit over $n_{\max} = 1$--$3$; the same value is obtained for NaH,
MgH$_2$, HCl and SiH$_4$; the former per-molecule equation labels are
retired along with the per-molecule spread).  The exact
identity of the exponents is a consequence of the factorization
$N_{\mathrm{Pauli}} = c(n_{\max})\,Q$ with molecule-independent
$c$ (verified exactly for $n_{\max} \le 3$; at $n_{\max} = 4$ the
coefficients differ across molecules by $\sim 0.03\%$): block-diagonal
ERI structure makes the total count the \emph{sum} of per-block counts,
each obeying the same Gaunt selection rules.  The 3-point exponent sits
below the recomputed Gaussian molecular exponents
($\alpha = 3.92$--$4.25$), but the local slope rises with basis
(3.77 on the LiH $n_{\max} = 3 \to 4$ interval, comparable to the
atomic path), so the basis-axis gap narrows at large $n_{\max}$; the
robust advantage is the linear system-size axis.


\subsection{Composed 1-norm and Trotter cost}
\label{sec:composed_onenorm}

Table~\ref{tab:composed_onenorm} reports the Pauli 1-norm and
first-order Trotter step counts for the composed systems.

\begin{table}[t]
\caption{Pauli 1-norm and Trotter step counts for composed systems at
$t = 1$~a.u.  H$_2$O 1-norms are inflated by the $Z_{\mathrm{eff}} = 6$
Phillips--Kleinman barrier and should be interpreted cautiously.
\emph{Re-measured 2026-08-30}, closing the standing vintage note:\ every
$\lambda$ is the live composed builder's identity-included 1-norm with PK in
the Hamiltonian, and the Trotter columns are regenerated from it via
$r = \lambda t/\sqrt{2\epsilon}$ (the relation the previous table already
satisfied at every cell).  The retired pre-exact-rule column read
$8.17 / 42.08 / 133.37$, $16.01 / 37.33 / 202.49$,
$198.20 / 354.89 / 735.60$, $18{,}913 / 28{,}055$; the $354.89$ cell was
independently identified in Paper~20 as a deprecated legacy BeH$_2$
builder value.  Pauli counts are order-independent and were never affected.}
\label{tab:composed_onenorm}
\begin{ruledtabular}
\begin{tabular}{llrd{3.2}rr}
Molecule & $n_{\max}$ & $Q$ & \multicolumn{1}{c}{$\lambda$}
 & $r(10^{-3})$ & $r(10^{-4})$ \\
\hline
H$_2$ & 2 &  10 &     8.31 &        186 &        587 \\
      & 3 &  28 &    46.50 &      1{,}040 &      3{,}288 \\
      & 4 &  60 &   163.59 &      3{,}658 &     11{,}567 \\
\hline
LiH   & 1 &   6 &    15.35 &        343 &      1{,}085 \\
      & 2 &  30 &    38.55 &        862 &      2{,}726 \\
      & 3 &  84 &   225.24 &      5{,}036 &     15{,}927 \\
\hline
BeH$_2$ & 1 &  10 &   211.20 &      4{,}723 &     14{,}934 \\
      & 2 &  50 &   374.87 &      8{,}382 &     26{,}507 \\
      & 3 & 140 &   801.69 &     17{,}926 &     56{,}688 \\
\hline
H$_2$O$^\dagger$ & 1 &  14 & 18{,}915 &  422{,}950 & 1{,}337{,}485 \\
      & 2 &  70 & 28{,}062 &  627{,}483 & 1{,}984{,}274 \\
\end{tabular}
\end{ruledtabular}
\end{table}

For LiH, the 1-norm grows from $\lambda = 15$ at $Q = 6$ to
$\lambda = 225$ at $Q = 84$, giving Trotter step counts of
$r \approx 5{,}000$ at $\epsilon = 10^{-3}$ for the largest basis
tested.  At $Q = 30$ ($n_{\max} = 2$), LiH requires only 862~Trotter
steps at chemical accuracy ($\epsilon = 10^{-3}$)---compared to the
published Trenev Gaussian LiH STO-3G value of 276~Pauli terms at $Q = 10$,
where the 1-norm (not published) would require a comparable or larger
step count on fewer qubits.

For BeH$_2$, the 1-norm is dominated by the $Z^2$-scaled PK barrier
on the Be core, giving larger absolute values but moderate growth:
$\lambda = 375$ at $Q = 50$ to $\lambda = 802$ at $Q = 140$.

$^\dagger$\textbf{H$_2$O 1-norm.}  The raw H$_2$O 1-norms are
dominated by the $Z^2$-scaled Phillips--Kleinman barrier at
$Z_{\mathrm{eff}} = 6$.  PK contributes 98.7\% of the total 1-norm:
the single PK diagonal element on the O-side 1s orbital is
$\sim 2{,}387$~Ha, replicated across four O-side valence blocks
(2~bond + 2~lone pair)---measured 2026-08-30 as $2386.769$~Ha exactly
four times.  Excluding PK, the electronic 1-norm is
$\lambda_{\mathrm{elec}} = 250$~Ha ($n_{\max} = 1$) and $372$~Ha
($n_{\max} = 2$), comparable to BeH$_2$'s total 1-norm.  For LiH
($Z = 3$), PK contributes only 11.7\% of the 1-norm ($\lambda_{\mathrm{elec}} = 34$
vs.\ $\lambda_{\mathrm{total}} = 39$ at $Q = 30$).  The $Z^2$ scaling
of PK parameters inflates the raw 1-norm for heavier atoms.

\textbf{PK classical partitioning.}
Because PK is a one-body operator, its energy contribution
$E_{\mathrm{PK}} = \mathrm{Tr}(h^{(1)}_{\mathrm{PK}} \cdot \gamma)$
can be computed classically from the 1-RDM $\gamma$ measured during VQE,
with zero additional quantum circuits.  The operator decomposition
$H_{\mathrm{full}} = H_{\mathrm{elec}} + H_{\mathrm{PK}}$ is verified
to machine precision ($< 10^{-12}$) for all three composed systems,
and the algebraic identity
$E_{\mathrm{full}} = \langle\psi|H_{\mathrm{elec}}|\psi\rangle
+ \langle\psi|H_{\mathrm{PK}}|\psi\rangle$
holds with residual $< 10^{-13}$~Ha.
This quantum--classical partitioning reduces the H$_2$O quantum
Hamiltonian 1-norm from $28{,}061.9$~Ha to $372.5$~Ha ($75.3\times$),
with no approximation.  The partitioned electronic 1-norms across the
composed systems all sit far below their with-PK values: LiH $34.0$~Ha,
BeH$_2$ $68.8$~Ha, H$_2$O $372.5$~Ha (re-measured 2026-08-30; the
retired figures were $359$ / $78\times$ / $33$ / $66$).

The commutator-based bound (Sec.~\ref{sec:onenorm}) provides an
additional $4$--$5\times$ step reduction for composed systems at
$n_{\max} \ge 2$, consistent with the single-geometry tightening
ratio.


\subsection{Comparison with Gaussian bases}
\label{sec:composed_gaussian}

We compare the composed GeoVac Pauli counts against the published
Gaussian data of Trenev~\textit{et al.}~\cite{trenev2025}
(Table~5, Appendix~B) for LiH and H$_2$O
(Table~\ref{tab:gaussian_published}).

For LiH, the three published Gaussian data points (STO-3G at $Q = 10$,
6-31G at $Q = 20$, cc-pVDZ at $Q = 36$) yield a fitted exponent
$\alpha_G = 4.25$ ($R^2 = 0.999$).  For H$_2$O, the corresponding
points ($Q = 12, 24, 46$) give $\alpha_G = 3.92$ ($R^2 > 0.999$).

Table~\ref{tab:equal_qubit} shows the equal-qubit comparison for
H$_2$O, using our fitted Gaussian power law for interpolation.

\begin{table}[t]
\caption{Equal-qubit Pauli term comparison for H$_2$O.  Gaussian
values are interpolated from the $Q^{3.92}$ power law (our log-log fit
of the published counts of Trenev~\textit{et al.}~\cite{trenev2025}).
GeoVac uses raw Jordan--Wigner; Gaussian uses 2-qubit reduction
(conservative for GeoVac).}
\label{tab:equal_qubit}
\begin{ruledtabular}
\begin{tabular}{rrrd{4.0}d{3.1}}
$n_{\max}$ & $Q$ & \multicolumn{1}{c}{GeoVac} & \multicolumn{1}{c}{Gaussian (interp.)}
  & \multicolumn{1}{c}{Ratio} \\
\hline
1 &  14 &         22 &       1\,126 & 51 \\
2 &  70 &      1\,954 &     580\,688 & 297 \\
3 & 196 &     99\,247 &  31\,457\,102 & 317 \\
\end{tabular}
\end{ruledtabular}
\end{table}

At $Q = 196$, the composed GeoVac Hamiltonian for H$_2$O requires
$\sim 320\times$ fewer Pauli terms than the Gaussian power-law
interpolation at the same qubit count, and $\sim 50\times$ at the
smallest basis ($Q = 14$).  (Corrected 2026-08-29: the retired
pair-diagonal counts gave $51\times$--$1{,}712\times$ here.)  The
multiplier grows from $Q = 14$ to $Q = 196$ because the exact-rule
exponent ($3.17$ over this range) still sits below the Gaussian
$3.92$; the paper's decisive comparison, however, is the
\emph{system-size} axis, where the composed count is exactly linear
in~$Q$ at fixed basis while the Gaussian count is not.

For direct (non-interpolated) comparison: at $Q = 70$ ($n_{\max} = 2$),
GeoVac produces \gvq{h2o_composed_pauli}{1\,954}~Pauli terms for H$_2$O, while the recomputed
Gaussian STO-3G value is 551~terms at $Q = 12$.  At $Q = 196$
($n_{\max} = 3$), GeoVac's 99\,247~terms sits just below the
published Trenev Gaussian cc-pVDZ count of \gvq{h2o_ccpvdz_pauli}{107\,382} at $Q = 46$
($0.92\times$; under the retired rule this read ``far below'' at
$5.8\times$).  These
comparisons at \emph{different} qubit counts illustrate that GeoVac
achieves comparable term counts while encoding substantially more
orbital degrees of freedom.


\subsection{Cross-block ERI bound}
\label{sec:cross_block}

The pseudopotential approximation sets cross-block valence ERIs to
zero.  To bound the impact of this approximation, we count the number
of cross-block index quadruplets satisfying the Gaunt selection rules.

For two valence blocks with identical angular quantum numbers, the
cross-block ERI count equals the same-block count exactly, because
both blocks use the same hydrogenic orbital set.  Including
cross-block ERIs would at most double the ERI Pauli contribution.
Since $N_{\mathrm{ERI}}$ dominates $N_{\mathrm{h1}}$ at all basis
sizes, this would at most double the total Pauli count---shifting
the exponent by at most $\log 2 / \log Q_{\max} \approx 0.13$,
keeping it well below~3.0 even in the worst case.

Empirical verification (Track~CB, v2.0.37, measured under the retired
pair-diagonal rule --- the ratio, not the absolute counts, is the
claim): for LiH at $n_{\max}=2$
($Q = 30$), the cross-block Coulomb ERI count was~130 versus~195
within-block (ratio~$2/3$), and the Pauli count with cross-block
ERIs included was~854 ($2.56\times$ the pair-diagonal composed~334;
the exact-rule composed count is~838).  The $\le
2\times$ bound is slightly exceeded because the theoretical bound
counts only same-angular-set pairs, while the core block
($Z = 3$, max $l = 1$) has different radial extent than valence
($Z_{\mathrm{eff}} = 1$), producing additional nonzero quadruplets.
The 1-norm increases from~37.3 to~85.7~Ha ($2.30\times$), and QWC
groups from~21 to~88 ($4.19\times$).

Importantly, simply adding cross-block ERIs while removing the PK
pseudopotential does \emph{not} improve accuracy: a 4-electron FCI
test gives 29\% error (coupled) versus 15\% (composed with PK)
versus 10.9\% (no PK, no cross-block).  The root cause is that the
composed orbital basis uses different effective charges per block,
so cross-block ERIs add electron--electron repulsion without the
balancing cross-center nuclear attraction (one-body terms where
each orbital feels all nuclei).  This confirms that the
pseudopotential is not redundant in second quantization for the
composed basis: PK encodes a one-body effective potential that
approximates both core--valence repulsion and orthogonality
enforcement, and replacing it with explicit two-body terms is
incomplete without also adding two-center one-body terms.
The block-diagonal ERI structure thus remains the correct
operating point for the composed architecture.

\textit{Note added (v2.0.42).}
Paper~19~\cite{loutey_paper19} resolves the incompleteness identified above.
The ``balanced coupled'' Hamiltonian adds both cross-block ERIs
\emph{and} cross-center nuclear attraction integrals
$\langle \psi_{nlm}^A | V_B | \psi_{n'l'm'}^A \rangle$, computed
analytically via incomplete gamma functions (machine precision).
The multipole expansion of $V_{ne}$ terminates exactly at
$L_{\max} = 2 l_{\max}$ by Gaunt selection rules.  The resulting
Hamiltonian produces the only bound LiH configuration in
4-electron FCI (1.8\% energy error vs PK's 15.0\%) and achieves
0.20\% energy error at $n_{\max} = 3$ with analytical integrals.
The resource tradeoff is regime-dependent: for VQE/NISQ (Pauli-count
dominated), composed with PK remains cheaper; for QPE (1-norm
dominated), balanced coupled eliminates PK entirely, which is
decisive for heavy atoms where PK's $Z^2$ barrier dominates
the 1-norm.  See Sec.~\ref{sec:balanced_coupled} for the
resource census.


\subsection{Limitations of the composed-system analysis}
\label{sec:composed_limitations}

The comparison in this section quantifies quantum \emph{simulation cost}
(Pauli terms, measurement groups, 1-norm), not basis-set quality.
GeoVac uses hydrogenic spin-orbitals in natural geometry blocks while
Gaussian methods use contracted GTOs; the physical content per qubit
differs.  The metrics that determine quantum simulation cost---Pauli
term count, 1-norm, and qubit count---are independent of the classical
PES accuracy.  Several approximations enter the composed-system
results, and we document them below, distinguishing those that affect
simulation cost from those that affect only energy accuracy:
\begin{enumerate}
\item \textbf{Cross-block ERIs set to zero:} The pseudopotential
  approximation eliminates all cross-block two-electron integrals.
  Selection-rule counting (Sec.~\ref{sec:cross_block}) bounds the
  impact at $\le 2\times$ the reported Pauli counts.
  \textit{Resolved:} Paper~19~\cite{loutey_paper19} includes all
  cross-block ERIs with Gaunt selection rules preserved across
  block boundaries.
\item \textbf{Cross-center nuclear attraction omitted:} One-body
  cross-center terms (e.g., H~orbital feeling the Li~nucleus) are
  not included.  These affect $N_{\mathrm{h1}}$ but not ERI scaling;
  at typical bond lengths the cross-center overlap decays
  exponentially.
  \textit{Resolved:} Paper~19 computes these integrals analytically
  via incomplete gamma functions.  33~terms at $n_{\max}=2$,
  168 at $n_{\max}=3$.
\item \textbf{Phillips--Kleinman parameters $Z^2$-scaled:} The
  repulsive barrier coefficients are $Z^2$-scaled from the
  Li$^{2+}$ core, not computed ab initio for each system.  This
  affects energy accuracy but not the Pauli count scaling.
\item \textbf{Valence $Z_{\mathrm{eff}}$ is constant:} The effective
  nuclear charge seen by valence electrons is position-independent,
  not $r$-dependent.  A self-consistent $Z_{\mathrm{eff}}(r)$ would
  change energies but not the orbital block structure.
\item \textbf{Gaussian comparison uses 2-qubit reduction:} The
  published Trenev Gaussian Pauli counts use a 2-qubit reduction
  (particle number and spin symmetry); GeoVac does not.  This
  makes the comparison conservative for GeoVac.
\item \textbf{$R_{\mathrm{eq}}$ error measures PES shape, not
  single-point quality:} The reported $R_{\mathrm{eq}}$ errors
  (5.3\% for LiH, 11.7\% for BeH$_2$, 26\% for H$_2$O) reflect
  the equilibrium distance of the composed PES, not the energy
  accuracy at any particular geometry.  For quantum simulation at
  a fixed nuclear configuration---the standard use case in
  fault-tolerant resource estimation~\cite{lee2021,reiher2017}---the
  relevant comparison is simulation cost at the same geometry and
  qubit count.
\item \textbf{Valence-valence coupling not captured:} The
  block-diagonal approximation eliminates inter-bond-pair electron
  interactions.  For BeH$_2$, this produces a qualitatively correct PES
  (bound state at $R_{\mathrm{eq}} = 3.31$~bohr) but quantitatively
  poor accuracy (32\% error vs.\ experiment).  Two correction approaches
  were tested: (i)~$Z_{\mathrm{eff}}$ partitioning (dividing screened
  charge among bonds) worsened $R_{\mathrm{eq}}$ to 103\% error;
  (ii)~classical point-charge inter-bond repulsion moved
  $R_{\mathrm{eq}}$ in the wrong direction.  The root cause is that bond
  pairs sharing a nucleus require orbital-level coupling (exchange and
  orthogonalization), not scalar energy corrections.  This limitation
  applies to the \emph{accuracy} of composed polyatomic PES, not to the
  Pauli term scaling, which remains valid as a lower bound.
  \textit{Partially resolved:} Paper~19 includes cross-block ERIs
  and cross-center $V_{ne}$ between all block pairs.  BeH$_2$
  4-electron FCI improves from 27.4\% (PK) to 10.7\% (balanced).
\end{enumerate}


\subsection{Transcorrelated variant}
\label{sec:tc_composed}

The electron-electron cusp ($1/r_{12}$ singularity) limits basis
convergence at each level of the natural geometry hierarchy
(Paper~18, Sec.~II).  The transcorrelated (TC) method removes
this singularity via a similarity transformation
$H_{\mathrm{TC}} = e^{-J} H\, e^{J}$ with Jastrow factor
$J = \tfrac{1}{2} r_{12}$.  For the composed qubit pipeline,
the TC transformation replaces the standard Coulomb ERI with a
gradient operator $G = -\tfrac{1}{2}\,\hat{r}_{12} \cdot
(\nabla_1 - \nabla_2)$, while leaving one-body integrals unchanged.
A constant shift of $-\tfrac{1}{4}$ per electron pair arises from
the Baker-Campbell-Hausdorff second-order term
$\tfrac{1}{2}|\nabla J|^2$.  The BCH expansion terminates exactly
at second order for two-electron blocks, and three-body terms vanish
between composed blocks (no shared electrons).

The TC Hamiltonian is non-Hermitian: the gradient operator $G$ breaks
the chemist symmetry $\langle pq|rs\rangle \ne \langle rs|pq\rangle$,
so the Jordan-Wigner transformation produces additional Pauli terms
from the antisymmetric part of the ERI.  However, the Gaunt angular
selection rules---which determine \emph{which} index quadruplets are
nonzero---are preserved, because $J = \tfrac{1}{2}r_{12}$ is
rotationally invariant.  The TC transformation densifies the radial
integrals within each allowed angular channel but does not create
new angular channels.  The quantum-algorithm cost of this
non-Hermiticity---and whether it is worth avoiding---is analysed in
Sec.~\ref{sec:tc_nonhermitian_cost}.

Table~\ref{tab:tc_composed} compares standard and TC composed
Hamiltonians at the production operating point ($n_{\max} = 2$).

\begin{table}[t]
\caption{Standard vs.\ transcorrelated (TC) composed Hamiltonians
at $n_{\max} = 2$.  ``Pauli ratio'' is TC/Standard Pauli term count.
``1-norm ratio (elec.)''\ compares electronic-only 1-norms with PK
classically partitioned.  The TC constant factor of $1.616\times$ is
uniform across all molecules.  \emph{Re-measured 2026-08-30} under the
exact global-$M_L$ rule, discharging the vintage note that stood here:\
the retired pair-diagonal row read $334/562$, $556/934$, $778/1{,}307$
with 1-norm ratios $1.09/1.16/1.00$, and the constant was $1.68\times$.
Pauli counts exclude the identity term.}
\label{tab:tc_composed}
\begin{ruledtabular}
\begin{tabular}{lrrrr}
Molecule & $Q$ & {Std Pauli} & {TC Pauli}
  & {1-norm ratio} \\
\hline
LiH    &  30 &    837 & 1\,353 & 1.28 \\
BeH$_2$ &  50 & 1\,395 & 2\,255 & 1.29 \\
H$_2$O  &  70 & 1\,953 & 3\,157 & 1.25 \\
\end{tabular}
\end{ruledtabular}
\end{table}

The Pauli ratio is $1.616\times$ for all three molecules, identical to
three decimals --- so the cross-molecule uniformity survived the exact-rule
correction, with the constant moving from the retired $1.68\times$.
\emph{The scope caveat from the 2026-08-29 vintage note is retained, not
retired with the numbers:}\ this is a fixed-basis ($n_{\max} = 2$)
cross-molecule measurement, so it establishes uniformity \emph{across
molecules} and not that the TC overhead leaves the basis-axis exponent
unchanged --- the same single-point-fallacy risk the correction exposed
elsewhere.  Electronic-only 1-norm overhead is a tight $25$--$29\%$
($1.283$ / $1.293$ / $1.253\times$ for LiH / BeH$_2$ / H$_2$O).  The
retired figures were $9$--$16\%$ with H$_2$O at exactly $1.00\times$;
that coincidence, and the PK-dominance explanation built on it, do not
survive re-measurement.

For He at the atomic level, the TC variant eliminates a basis
divergence: standard hydrogenic FCI error grows from
5.3\% ($n_{\max} = 1$) to 8.2\% ($n_{\max} = 3$), while TC
error converges to a 3.3--3.6\% plateau.  The plateau persists
even with the full gradient operator (Track BX-4); the
radial-only restriction accounts for only 0.01~pp of the
$\sim$3.3\% floor, which is intrinsic to the hydrogenic basis
at this truncation level.

\textit{Angular gradient benchmark (Track BX-4).}
The full angular gradient term of $G$ was implemented and benchmarked
for $l > 0$ orbitals.  For He, including the angular gradient reduces
the FCI error by 0.010--0.014~pp: radial-only gives 3.621\% at
$n_{\max} = 2$ and 3.639\% at $n_{\max} = 3$; full TC gives 3.611\%
and 3.625\%, respectively.  The accuracy gain is negligible.  The
Pauli cost, however, is substantial: 188 $\to$ 500 terms at
$n_{\max} = 2$, and 1{,}492 $\to$ 3{,}449 ERIs at $n_{\max} = 3$.
\emph{The composed-pipeline angular-gradient multipliers are withdrawn}
(2026-08-30):\ they read $562 \to 1{,}498$ (LiH), $936 \to 2{,}496$
(BeH$_2$), $1{,}310 \to 3{,}494$ (H$_2$O), a uniform $2.67\times$ and a
$4.49\times$ total, but the operator producing them is defective.  Every
entry the angular gradient adds to the composed ERI tensor violates
$L_z$:\ measured $78/78$ added entries for LiH, against a control in
which the plain composed and radial-only TC tensors have \emph{zero}
violations under the same test.  A two-body Coulomb-type operator cannot
add $L_z$-violating support, so those entries are spurious and the
multipliers measure the defect as much as the operator.  They are
withdrawn rather than re-measured, because re-measuring a broken operator
produces a number that has to be disclaimed in the same sentence.
The angular gradient remains a \emph{negative result for quantum
efficiency} --- the verdict is now over-determined, since the accuracy
gain was already negligible (0.010--0.014~pp) and the added support is
unphysical.  The radial-only TC Hamiltonian
(\texttt{include\_angular=False}) remains the default for the composed
pipeline.

A PES scan for LiH confirms that both the standard and TC Pauli
counts are $R$-independent (837 and 1{,}353 at all sampled bond lengths
under the exact rule; the scan was run at the retired 334 and 562, and
$R$-independence is a support property that the value correction does not
disturb),
consistent with the selection-rule origin of the sparsity.

\subsubsection{Quantum-algorithm cost of the non-Hermitian form}
\label{sec:tc_nonhermitian_cost}

The similarity transformation makes $H_{\mathrm{TC}}$ non-Hermitian,
placing it outside the Hermitian resource model used above
(Rayleigh-quotient VQE, real-time Trotter, and textbook QPE all assume a
Hermitian generator).  We characterise this cost on an atomic
Coulomb--Sturmian realisation and reach two conclusions: the
non-Hermitian form is cheap to run, and the naive Hermitisation that
would avoid it sacrifices accuracy.

\emph{The non-Hermiticity is localised and benign.}  For a
three-electron reference (Li~$1s^2 2s$) the electron--electron cusp
activates a genuine three-body operator, which we contract to an
effective two-body operator using the reference one-particle density
matrix (the xTC scheme~\cite{christlmaier2023xtc}).  This contracted
operator is Hermitian to machine precision
($\lVert A - A^\dagger\rVert \le 3\times10^{-18}$ for an $s$ reference,
$3\times10^{-17}$ for an open-$p$ reference); the only non-Hermitian
term is the generic convective operator $G$, present in any
transcorrelated Hamiltonian.  For the full operator the spectrum is
entirely real, the ground state real, and the departure from normality
small ($0.015$ for Li, $0.069$ for He).  The cost of eigenvalue
estimation for a non-normal block-encoded operator is governed by the
eigenvector condition number $\kappa_V$ (the Bauer--Fike factor), which
enters \emph{linearly} in the optimal quantum eigenvalue-processing
route~\cite{lowsu2024qeve}.  We measure $\kappa_V = 1.21$ for Li
(grid-stable) and $\kappa_V = O(1)$, $\lesssim 2$, for He, so the penalty
is a small constant factor rather than a change of scaling; the Pauli
1-norm is barely affected
($\lambda_{\mathrm{TC}}/\lambda_{\mathrm{std}} \le 1.13$).  Independent
work on transcorrelated electronic Hamiltonians in Gaussian bases
reaches the same conclusion ($\kappa_V \approx 1$--$3$, reduced further
by xTC)~\cite{uvarov2025qeve_tc}.  The one genuine overhead is the Pauli
\emph{count}: the non-symmetric $G$ tensor roughly doubles it
($117 \to 249$ on the atomic $s$-only engine).

\emph{Hermitisation by symmetrisation is a false economy.}  The
symmetrised operator
$\tfrac12(H_{\mathrm{TC}} + H_{\mathrm{TC}}^\dagger)$ is attractive on
cost---it restores the plain Pauli count, lowers the 1-norm slightly
below plain ($0.98\times$), and is Hermitian ($\kappa_V = 1$)---but it is
a different, less accurate operator.  Discarding the anti-Hermitian part
of $G$ shifts the energy downward past the exact value, and the
overshoot \emph{worsens} as the basis improves: at the largest basis
tested the genuine transcorrelated energy lies $+5.4$~mHa (He) and
$+7.7$~mHa (Li) above the exact non-relativistic energy and approaches it
from the variationally correct side, whereas the symmetrised energy lies
$-6.6$ and $-21.4$~mHa \emph{below} exact.  No stationary Jastrow range
anchors the symmetrised operator ($E_{\mathrm{TC}}$ is monotone in the
geminal exponent).  We therefore run the operator in its native
non-Hermitian form---affordable by the $\kappa_V = O(1)$ result
above---rather than symmetrising.  Deleting $G$ altogether is not an
alternative (the remaining two-body term then over-binds by
$\sim0.6$~Ha), and folding to
$H_{\mathrm{TC}}^\dagger H_{\mathrm{TC}}$ is excluded because it returns
the smallest singular value, not the ground eigenvalue, and squares the
conditioning.  These measurements use the atomic $s$-only
Coulomb--Sturmian engine (He, Li) with a single spin-independent
geminal; the robust content is the \emph{ordering} (native non-Hermitian
more accurate than symmetrised, both correcting the cusp relative to
plain), not an absolute accuracy figure for this minimal basis.

\subsubsection{Angular sparsity survives the cusp correction (atomic xTC)}
\label{sec:tc_atomic_sparsity}

The cost analysis above establishes that the non-Hermitian TC operator is
affordable.  The complementary question is whether it is still \emph{sparse}:
when the electron--electron cusp activates a genuine three-body operator and we
contract it to an effective two-body operator against the reference density (the
xTC scheme~\cite{christlmaier2023xtc}), does the contracted operator inherit the
angular Gaunt sparsity of the plain Coulomb operator (Paper~22), or does the
three-body $\to$ two-body contraction fill in the zero angular blocks?  The
question is sharp because the transcorrelated vertex is \emph{non-abelian}: a leg
harmonic and a correlator harmonic Clebsch--Gordan couple to a \emph{multiplet}
of resultants, not a single one, which is exactly why the three-body operator does
not collapse to two-body the way the abelian plane-wave TC does.

The answer is settled at the level of angular \emph{support}---which orbital
quadruplets are nonzero---which is radial-independent and therefore exact (integer
block counts, no quadrature).  Table~\ref{tab:xtc_fillin} gives the fill-in (blocks
xTC makes nonzero where Coulomb was zero) for spherical, open-$p$, and open-$d$
references across increasingly rich bases, swept over the whole monoatomic library
$Z=1$--$56$.

\begin{table}[t]
\caption{Exact angular fill-in of the xTC-contracted two-body operator relative to
the plain Coulomb operator, swept across the monoatomic library.  ``Fill-in'' counts
blocks nonzero under xTC but zero under Coulomb; fill-out is zero in every case, so
zero fill-in means the support (hence the Jordan--Wigner Pauli term set) is preserved
exactly.  Spherical reference: any closed, half-filled high-spin, or $s$-open shell
(Unsold; e.g.\ Be $1s^2 2s^2$, N $2p^3$, Mn $3d^5$); open-$p$: C/O ($2p^2$/$2p^4$);
open-$d$: Sc/Ti-type ($3d^1$/$3d^2$).  Counts are exact integers (angular support is
radial-independent).}
\label{tab:xtc_fillin}
\begin{ruledtabular}
\begin{tabular}{lrrrr}
Basis & {spherical} & {open-$p$} & {open-$d$} & {Coulomb density} \\
\hline
$s+p$       & 0 & 0   & 0   & 14.8\% \\
$s+p+d$     & 0 & 4   & 4   & 8.5\%  \\
$s+p+d+f$   & 0 & 96  & 96  & 6.1\%  \\
$s+p+d+f+g$ & 0 & 264 & 268 & 4.8\%  \\
\end{tabular}
\end{ruledtabular}
\end{table}

The fill-in is governed by a single exact rule: it is \emph{zero at every basis},
bit-exactly, precisely when the reference density is spherical.  By Uns\"old's theorem
the spherical shells are the closed subshells, the half-filled high-spin subshells,
and the $s$-open shells---so the zero-fill-in guarantee covers not only the framework's
headline molecules H$_2$, LiH, BeH$_2$ (all $s$-referenced, Be is $1s^2 2s^2$) but
also, e.g., nitrogen ($2p^3$) and chromium/manganese ($3d^5$).  Across the
$Z=1$--$56$ library this is 29 of the 56 atoms; the remaining 27 split into 16
open-$p$ and 11 open-$d$ references.  The mechanism is the same monopole collapse in
every spherical case: a spherical reference density carries only the $L'=0$ multipole
(Uns\"old), so contracting a correlator leg against it forces the leg to the monopole
and collapses the non-abelian vertex to a single Coulomb-like multipole with identical
angular support.  Because the support is unchanged (and no block is removed either),
the Jordan--Wigner Pauli term count is preserved exactly; on the $s+p$ radial engine
(Li reference) the plain and xTC operators both carry \gvq{block_sp_n2_pauli}{279} Pauli terms, and the
two-body 1-norm is \emph{lowered} to $0.76$--$0.91\times$ plain across geminal widths
$\gamma \in [0.6, 1.5]$.  So xTC handles the cusp without inflating the qubit cost.

A non-spherical open-$p$ reference (partially filled $2p$ that is not half-filled:
C, O, F, and the oxygen of H$_2$O---but \emph{not} N, whose $2p^3$ is half-filled and
therefore spherical) carries the additional $L'=2$ density multipole, which lets the
vertex bridge an even-multipole gap between external pairs.  The fill-in is still zero
through $s+p$ (at $l \le 1$ every same-parity pair already contains the monopole, so
Coulomb misses nothing), and only a small, structured fill-in appears once $d$
orbitals enter---4 blocks at $s+p+d$, 96 at $s+p+d+f$, 264 at $s+p+d+f+g$.  An
open-$d$ reference (the transition metals whose $d$-shell is neither empty, half, nor
full) additionally carries the $L'=4$ multipole; this makes \emph{no} difference
through $s+p+d+f$ (identical $0/4/96$) and first bites at $s+p+d+f+g$, where the
open-$d$ fill-in is 268 against the open-$p$ 264---a genuine but tiny $L'=4$ effect
($+4$ blocks out of $390{,}625$).  In every case the xTC density stays within $2.5\%$
relative of the Coulomb density (both fall from $14.8\%$ at $s+p$ to $4.8\%$ at
$s+p+d+f+g$), so the operator does not densify across the physical range; the honest
difference from the spherical case is that the preservation is no longer bit-exact and
the fill-in fraction rises slowly with $l_{\max}$.  The whole pattern is a periodicity
statement (Paper~16): the fill-in class is fixed by the open shell's angular momentum,
so it repeats down each group of the periodic table.

\emph{This is an atomic property.}  The angular sparsity that xTC preserves is the
single-center $l$-selection of Paper~22, and that sparsity does not survive to two
centers: a two-center Coulomb operator already sits at $100\%$ of its
$m$-conservation limit (LiH) against the atom's $55\%$, so a molecule has lost the
rich $l$-sparsity \emph{before} xTC is applied.  For a molecule xTC therefore
preserves only the $m$-conservation sparsity plus a modest 1-norm reduction; the
combined ``cusp corrected and angular sparsity preserved'' result is strongest for
atoms.  As throughout this section, these are single-common-$k$ Coulomb--Sturmian
measurements with a single spin-independent geminal on a classical non-Hermitian
proof of concept---structural (support and resource) results, not absolute accuracy
claims for the minimal basis.

\subsection{Balanced coupled variant}
\label{sec:balanced_coupled}

Paper~19~\cite{loutey_paper19} presents a PK-free alternative: the
balanced coupled Hamiltonian includes all four interaction
types---within-block $h_1$ and ERI, cross-block ERI, and cross-center
nuclear attraction $\langle \psi^A | V_B | \psi^A \rangle$---with no
pseudopotential.  Cross-center $V_{ne}$ integrals are computed
analytically via incomplete gamma functions; the multipole expansion
terminates exactly at $L_{\max} = 2 l_{\max}$ by Gaunt selection
rules.

Table~\ref{tab:balanced} compares composed and balanced encodings
at $n_{\max} = 2$.

\begin{table}[t]
\caption{Composed vs.\ balanced coupled resource census at
$n_{\max} = 2$.  \emph{Pauli counts and $\lambda$ both exclude the
identity term}, on both sides --- re-synced 2026-08-30, when the balanced
columns were found to be identity-\emph{included} and the composed
$\lambda$ column identity-included, under a caption declaring otherwise.
Composed~$\lambda$ values are electronic-only 1-norms
(PK classically partitioned via 1-RDM; the partitioning machinery is
implemented in \texttt{geovac/pk\_partitioning.py} per CLAUDE.md~\S~7).
Balanced~$\lambda$ values include cross-center $V_{ne}$ (no PK).
Paper~20's resource table reports the same systems with the identity
term \emph{in}, which is why its cells are larger.}
\label{tab:balanced}
\begin{ruledtabular}
\begin{tabular}{lrrrrr}
System & $Q$ & \multicolumn{2}{c}{Pauli terms}
  & \multicolumn{2}{c}{$\lambda$ (Ha)} \\
  & & Comp. & Bal. & Comp. & Bal. \\
\hline
LiH     &  30 &    837 & 2{,}725 &  27.7 &   75.2 \\
BeH$_2$ &  50 &  1{,}395 & 8{,}867 &  54.3 &  289.6 \\
H$_2$O  &  70 &  1{,}953 & 19{,}741 & 186.8 & 1{,}439.8 \\
\end{tabular}
\end{ruledtabular}
\end{table}

The Pauli ratio grows with block count ($3.26\times$ to $10.1\times$),
but the 1-norm ratio favors balanced for BeH$_2$ ($0.88\times$
composed with PK in-Hamiltonian, the QPE-relevant convention; live
builders, re-measured 2026-08-30 --- $328.1$ vs.\ $374.9$~Ha, both
identity-included, since the composed-with-PK total is only available
that way) where PK's $Z^2$-scaled barrier is costly.  For H$_2$O,
the balanced 1-norm ($1{,}612.0$~Ha) is $17.4\times$ lower than the
full composed 1-norm with PK ($28{,}061.9$~Ha).  Composed and balanced
are thus complementary encoding strategies: composed is cheaper for
VQE (Pauli-count dominated), balanced is cheaper for QPE on heavy
atoms (1-norm dominated, PK eliminated).  The balanced coupled
encoding is available via \texttt{build\_balanced\_hamiltonian()}
in the production codebase.

\subsection{Discussion of composed results}
\label{sec:composed_discussion}

The composed-geometry approach preserves and improves the structural
sparsity of single-geometry GeoVac Hamiltonians.  Three factors drive
this:
\begin{enumerate}
\item \textbf{Block-diagonal ERIs:} The pseudopotential approximation
  eliminates all cross-block two-electron integrals, reducing the
  ERI count from $O(M^4)$ to $O(\sum_b M_b^4)$.
\item \textbf{Angular selection rules preserved:} Each block retains
  the Gaunt integral sparsity from the single-center case, with ERI
  density decreasing as $\sim 1/M_b$ within each block.
\item \textbf{Composition is additive:} Adding a new atomic center
  adds a new ERI block but does not create cross-block entries.
\end{enumerate}
The universality of the within-molecule exponent ($3.17$, identical
across all eight systems tested, spanning 1--7 orbital blocks and three
periodic-table rows) is a consequence of an exact factorization: the
composed Pauli count is $N_{\mathrm{Pauli}} = c(n_{\max}) \times Q$
with molecule-independent rational coefficients $c(1) = 3/2$,
$c(2) = 279/10 = 27.90$, $c(3) = 7089/14 \approx 506.36$, so every
molecule's $(Q, N_{\mathrm{Pauli}})$ sweep is a horizontal translate of
every other's on the log--log plane.  Across the composed molecules at
$n_{\max} = 2$ the count is exactly linear in~$Q$:
$N_{\mathrm{Pauli}} = \gvq{composed_coeff}{27.90} \times Q$, confirming that the linear
system-size scaling extends to arbitrary polyatomic systems via the
fiber bundle decomposition~\cite{loutey_paper17}.  (Corrected
2026-08-29 from the retired pair-diagonal $11.10 \times Q$ and
$Q^{2.5}$; under the retired rule H$_2$ showed an anomalous exponent
$3.13$ vs.\ the molecules' $2.50$--$2.52$ --- the exact rule dissolves
the anomaly.)


\subsection{Second-row molecules via frozen-core composition}
\label{sec:second_row}

The composed architecture extends naturally to second-row atoms
($Z = 11$--18) by replacing the Level~3-solved He-like core
with an analytical frozen core.  Each [Ne]-like core
(1$s^2$\,2$s^2$\,2$p^6$, 10~electrons) is described by hydrogenic
wavefunctions with Clementi--Raimondi orbital exponents, producing
a smooth $Z_{\text{eff}}(r)$ screening function ($Z$ at the origin,
$Z-10$ at infinity).  The frozen-core energy enters the qubit
Hamiltonian only through the identity operator and does not affect
quantum measurement cost.

Table~\ref{tab:second_row} lists balanced coupled resource metrics
for six second-row molecules at $n_{\max} = 2$.

\begin{table}[t]
\caption{Balanced coupled resource metrics for second-row molecules
  at $n_{\max} = 2$, with composed metrics for comparison.
  All molecules use [Ne] frozen cores with Clementi--Raimondi
  orbital exponents.  Re-measured 2026-08-29 under the exact
  global-$M_L$ rule; the balanced QWC column and two $\lambda_{\text{ni}}$
  cells (NaH $19.6 \to 20.6$, HCl $866.1 \to 869.7$) were measured
  2026-08-30, closing a standing ``pending'' note.  Pauli counts include
  the identity term; $\lambda_{\text{ni}}$ excludes it.}
\label{tab:second_row}
\begin{ruledtabular}
\begin{tabular}{lrrrrr}
Molecule & $Q$ & \multicolumn{2}{c}{Pauli terms}
  & $\lambda_{\text{ni}}$ (Ha) & QWC \\
  & & Comp. & Bal. & Bal. & Bal. \\
\hline
NaH      &  20 &    559 &    575 &  20.6 &    69 \\
MgH$_2$  &  40 & 1{,}117 & 4{,}861 & 110.5 &   903 \\
HCl      &  50 & 1{,}396 & 9{,}824 & 869.7 & 1{,}173 \\
H$_2$S   &  60 & 1{,}675 & 13{,}863 & 879.6 & 1{,}282 \\
PH$_3$   &  70 & 1{,}954 & 18{,}854 & 895.3 & 1{,}477 \\
SiH$_4$  &  80 & 2{,}233 & 24{,}745 & 914.1 & 1{,}551 \\
\end{tabular}
\end{ruledtabular}
\end{table}

The within-molecule Pauli scaling of all four second-row molecules
(NaH, MgH$_2$, HCl, SiH$_4$) yields an exponent of $3.17$,
identical to the first-row result to all reported digits (the
$N_{\mathrm{Pauli}} = c(n_{\max})\,Q$ factorization).
This confirms that the scaling is universal across
periodic table rows---driven by Gaunt selection rules that are
independent of the core treatment.

Because second-row molecules use frozen cores (no encoded core block),
they require fewer qubits than their first-row isostructural
analogues: NaH uses $Q = 20$ versus LiH's $Q = 30$.  Within a
given row, however, the Pauli count depends only on the orbital
block topology and $n_{\max}$.  The key structural observation is
that the Pauli \emph{scaling exponent} ($3.17$) is universal
across both rows, driven by Gaunt selection rules that are
independent of nuclear charge.

For 1-norm comparisons, we report
$\lambda_{\text{ni}}$---the non-identity component of the Pauli
1-norm---which excludes the frozen-core energy offset and nuclear
repulsion (classical constants that do not contribute to quantum
measurement cost or Trotter error).  For NaH, the identity term
accounts for the bulk of the total 1-norm ($193.4$~Ha total vs.\
$20.6$~Ha non-identity); the simulation-relevant
$\lambda_{\text{ni}} = 20.6$~Ha is comparable
to first-row LiH ($\lambda_{\text{ni}} = 27.7$~Ha).

The framework extends further to third-row main-group atoms
($Z = 19$--20 and $Z = 31$--36) via [Ar] (18-electron) and
[Ar]3$d^{10}$ (28-electron) frozen cores with Clementi--Raimondi
orbital exponents.  Six additional molecules---KH, CaH$_2$, GeH$_4$,
AsH$_3$, H$_2$Se, HBr---bringing the single- and two-center library
to 18 molecules across three periodic table rows (extended to
28 with multi-center systems in Sec.~\ref{sec:multi_center}).  The universal
$3.17$ scaling exponent is maintained.

A strong structural invariance is confirmed: frozen-core molecules
with the same block topology produce \emph{identical} Pauli counts.
For example, NaH and KH both yield 559 composed / 575 balanced
Pauli terms at $Q = 20$, despite different core types ([Ne] vs.\ [Ar])
(Table~\ref{tab:second_row}; re-measured 2026-08-30 under the exact
global-$M_L$ rule, which raised both counts from the retired
pair-diagonal $223/239$ while leaving the invariance itself intact).
This invariance holds exactly across all six isostructural pairs
spanning rows~2 and~3.

Basis convergence at $n_{\max} = 3$ (requiring $\ell = 2$ rotation
matrices, now implemented via the Wigner $D$-matrix) produces NaH
Hamiltonians at $Q = 56$ with 5{,}349 balanced Pauli terms.
Two-electron FCI gives a 0.40~Ha energy improvement over $n_{\max} = 2$
at the reference geometry, confirming variational convergence.
However, the potential energy surface retains the overattraction
observed at $n_{\max} = 2$---no equilibrium geometry is found---indicating
that PES shape accuracy is limited by the frozen-core
approximation, not the angular basis.


\subsection{Multi-center molecules}
\label{sec:multi_center}

The composed architecture handles molecules with more than two nuclei
through the same block-diagonal construction: each electron pair
occupies an independent orbital block, and cross-block ERIs are zero
by the pseudopotential approximation.  Table~\ref{tab:multi_center}
reports composed Pauli counts for the five multi-center systems
in the shipping library at $n_{\max} = 2$.

\begin{table}[t]
\caption{Composed (block-diagonal) Pauli term counts for multi-center
  molecules at $n_{\max} = 2$.  $N_e$: total electrons;
  $N_b$: orbital blocks; $Q = 2M$: qubits.  Frozen cores are used
  where applicable (Na: [Ne], Cl: [Ne]).  Counts are non-identity terms,
  measured from the shipping library; each is $\gvq{composed_coeff}{27.90} \times Q$ exactly,
  the block coefficient of Sec.~\ref{sec:block_pauli_count}.  Corrected
  2026-08-29: the retired pair-diagonal rule gave $11.10 \times Q + 1$.}
\label{tab:multi_center}
\begin{ruledtabular}
\begin{tabular}{lrrrr}
Molecule & $N_e$ & $N_b$ & $Q$ & Pauli terms \\
\hline
LiF      & 12 & 6 &  70 &   1{,}953 \\
CO       & 14 & 7 & 100 &   2{,}790 \\
N$_2$    & 14 & 7 & 100 &   2{,}790 \\
F$_2$    & 18 & 9 & 100 &   2{,}790 \\
NaCl     &  8 & 4 &  50 &   1{,}395 \\
\end{tabular}
\end{ruledtabular}
\end{table}

The most striking feature of Table~\ref{tab:multi_center} is
\emph{isostructural invariance}: CO and N$_2$ produce identical
Pauli counts (2{,}790) despite different nuclear charges, confirming
that the invariance observed for frozen-core pairs
(Sec.~\ref{sec:second_row}) extends to multi-center topologies.
Similarly, CO, N$_2$, and F$_2$ all yield 2{,}790 terms at
$Q = 100$.  The Pauli count is determined entirely by the qubit
count~$Q$, not by the number of blocks, the block structure, or the
atomic species.

Across the 35 composed molecules in the library---spanning
$Z = 1$--56 and including the transition metal hydrides
(Sec.~\ref{sec:d_orbitals})---the composed Pauli count obeys
\begin{equation}
  N_{\mathrm{Pauli}} = \gvq{composed_coeff}{27.90} \times Q,
  \label{eq:linear_pauli}
\end{equation}
where $N_{\mathrm{Pauli}}$ excludes the identity term.
The coefficient \gvq{composed_coeff}{27.90} ($= 279/10$) is exact (not a fit) across all main-group
molecules (corrected 2026-08-29; the retired pair-diagonal rule gave 11.10).
This exact linearity in~$Q$
is a stronger result than the within-molecule basis scaling
(Eq.~\ref{eq:alpha_lih}), which is measured by
varying $n_{\max}$ at fixed molecular topology.
Equation~\eqref{eq:linear_pauli} describes the \emph{across-molecule}
scaling at fixed $n_{\max} = 2$: adding more orbital blocks (more
nuclei, more electron pairs) increases~$Q$ and~$N_{\mathrm{Pauli}}$
in exact proportion, because each block contributes an independent
set of Pauli terms with no cross-block coupling.

With these five multi-center systems, the GeoVac library spans
37 systems (35 composed molecules plus the He and H$_2$ reference atoms)
across $Z = 1$--56 (H through Ba), including first-row
(H$_2$ through CH$_4$), second-row (NaH through SiH$_4$),
third-row (KH through HBr), multi-center (LiF, CO, N$_2$, F$_2$, NaCl),
and the ten transition metal hydrides (ScH--ZnH;
Sec.~\ref{sec:d_orbitals}).  All are accessible via the
\texttt{hamiltonian()} API with ecosystem export to Qiskit,
OpenFermion, and PennyLane.


\subsection{d-Orbital blocks and transition metal hydrides}
\label{sec:d_orbitals}

Higher angular momentum admits \emph{more} multipole orders, not
fewer:\ for $d$-orbital blocks ($l = 2$) the allowed orders are
$k = 0, 2, 4$ (three values) against $k = 0, 1$ for $s$+$p$.
\emph{Corrected 2026-08-30:} this section previously opened ``the Gaunt
selection rules become \emph{more} restrictive at higher angular
momentum'' and quoted 20\% against 38.5\% nonzero angular coefficients.
Both were artifacts of the retired pair-diagonal rule, which deleted
proportionally more of the $d$-block's genuine $m$-changing couplings.
Table~\ref{tab:d_block_eri} quantifies the
effect on ERI density and Pauli term counts for isolated blocks at
fixed qubit counts.

\begin{table}[t]
\caption{ERI density and Pauli term counts for single orbital blocks
  of different angular momentum composition at $n_{\max} = 2$.
  $M$: spatial orbitals; $Q = 2M$: qubits.}
\label{tab:d_block_eri}
\begin{ruledtabular}
\begin{tabular}{lrrrrr}
Block type & $M$ & $Q$ & ERI density & Pauli & Pauli/$Q$ \\
\hline
$s$-only ($l_{\max} = 0$) & 3 & 6 & 100\% & 118 & 19.67 \\
$s$+$p$ ($l_{\max} = 1$) & 9 & 18 & 15.7\% & 2{,}967 & 164.83 \\
$d$-only ($l = 2$) & 5 & 10 & 13.6\% & 343 & 34.30 \\
\end{tabular}
\end{ruledtabular}
\end{table}

The $d$-only block is \emph{not} the sparsest.  Its ERI density
(13.6\%) is close to the $s$+$p$ block's (15.7\%), and per qubit it is
the denser of the two:\ Pauli/$Q = 34.30$ against $\gvq{composed_coeff}{27.90}$ for the
$s$+$p$ block at $n_{\max}=2$ that dominates main-group molecules.
Under the retired pair-diagonal rule these read 4.0\% / 5.60 against
8.9\% / 54.22, making the $d$-block look dramatically leaner.  That was
an artifact:\ higher angular momentum has \emph{more} $m$-changing
integrals to drop, so the pair-diagonal restriction deleted
proportionally more of the $d$-block's genuine couplings.  The
correction is consistent with the composed $d$-block coefficient
($30.03$ against main-group $27.90$, Sec.~\ref{sec:eri_rule}):\ the
$d$-block is denser in both measurements.

Transition metal hydrides are constructed via frozen-core
composition.  ScH ($Z = 21$) and TiH ($Z = 22$) carry
an [Ar] frozen core (18 electrons) and valence blocks that include an
isolated $d$-orbital block ($l_{\min} = 2$, $l_{\max} = 2$) alongside
an $s$-orbital bond pair.  At $n_{\max} = 2$, both molecules require
$Q = 30$ qubits and produce 901 non-identity Pauli terms---more than
LiH ($\gvq{lih_composed_pauli}{837}$ at the same~$Q$)---giving Pauli/$Q = \gvq{composed_coeff_d}{30.03}$, slightly
\emph{above} the main-group coefficient of~27.90.  Under the exact
Coulomb rule $d$-blocks are marginally \emph{denser} than $s/p$
blocks, not sparser; the retired pair-diagonal figures (277 terms,
Pauli/$Q = 9.23$, ``below the universal 11.10'') inverted this
ordering.

The atomic classifier now automates the full first transition series
($Z = 21$--30, type-F dispatch), and all ten hydrides (ScH--ZnH) ship
in the 37-system library; the Hamiltonians remain resource-level
results---they have not been benchmarked against FCI or Gaussian
baselines.
Nevertheless, the structural finding---that $d$-orbital blocks have a
lower \emph{pair-diagonal} Pauli/$Q$ ratio than $s$/$p$ blocks
(Sec.~\ref{sec:eri_rule})---is an artifact of the pair-diagonal rule
dropping more $m$-swap integrals at higher angular momentum (under the
exact global-$M_L$ rule the $d$-block is denser, not sparser), not a
sparser exact selection rule.


\subsection{Compact nested encoding}
\label{sec:nested_encoding}

For single-center $N$-electron systems, a compact ``nested'' encoding
places all electrons in a single orbital block on $S^{3N-1}$, using
H-set coupled hyperspherical harmonics (Sec.~\ref{sec:recoupling_sparsity})
as the angular basis.  This eliminates the core/valence decomposition and
the PK pseudopotential entirely, reducing the qubit count by
${\sim}3{\times}$ for 4-electron atoms.

Table~\ref{tab:nested_be} compares the nested and composed encodings for
beryllium (4~electrons).

\begin{table}[t]
\caption{Nested vs.\ composed encoding for Be ($Z = 4$, 4~electrons,
  $n_{\max} = 2$).  $\lambda_{\rm ni}$: non-identity 1-norm.}
\label{tab:nested_be}
\begin{ruledtabular}
\begin{tabular}{lrrrr}
Encoding & $Q$ & Pauli & $\lambda_{\rm ni}$ (Ha) & FCI err \\
\hline
Nested (H-set) & 10 & 112 & 18.95 & 4.9\% \\
Composed + PK  & 12 & 115 & 121.35 & 1.0\% \\
\end{tabular}
\end{ruledtabular}
\end{table}

The nested encoding achieves comparable Pauli count (112 vs.~115) with
$6.4{\times}$ lower non-identity 1-norm (18.95 vs.~121.35~Ha) due to PK
elimination.  The accuracy tradeoff (4.9\% vs.~1.0\% FCI error) reflects
the single-center basis truncation at $n_{\max} = 2$.

For molecules, the nested encoding requires per-pair effective nuclear
charges (heterogeneous basis), which introduces L\"owdin
orthogonalization costs that destroy Gaunt-level sparsity: LiH with
heterogeneous nested encoding produces 1{,}711 Pauli terms (measured
under the retired pair-diagonal rule, vs.\ 120 uncoupled or 334 composed
under that same rule; the exact-rule composed count is 838 --- the
L\"owdin-inflation comparison is rule-consistent as quoted).  The composed architecture's block-diagonal
factorization avoids this by maintaining independent orthonormal bases at
each length scale.

The nested encoding is most advantageous for atomic systems and
single-point molecular calculations where qubit count is the binding
hardware constraint.


%% ========================================================================
%% §V Spinor composed encoding — Dirac-on-S^3 Tier 2 sprint drop-in (T6).
%% ========================================================================

\section{Spinor composed encoding}
\label{sec:spinor_composed}

The composed qubit Hamiltonians of \S\ref{sec:composed} use a spinless
orbital basis with a trivial factor-of-two spin doubling at
Jordan--Wigner encoding. This section reports the
\emph{spin-ful} extension: one qubit per Dirac label $(n, \kappa, m_j)$
with the angular structure provided by Szmytkowski's spinor spherical
harmonics and the spin-orbit coupling provided by the Breit--Pauli
leading-order reduction of the Dirac equation.
The construction is an engineering upgrade of
\S\ref{sec:composed}'s composed pipeline, not a new Dirac Fock
projection: the Schr\"odinger Fock map on $S^3$ does not lift to the
Dirac sector (a $\mathbb{Z}_2$-graded superalgebra
obstruction), so the spinor basis rides on the existing
scalar $S^3$ graph with relativistic coefficients added on top. All
matrix elements are closed-form in $(\kappa, m_j)$-native labels: no
numerical quadrature enters the angular layer, and the radial layer
reuses the exact-rational $R^k$ evaluator from
\texttt{hypergeometric\_slater.py}.

\subsection{Resource metrics for LiH, BeH, CaH}
\label{subsec:spinor_resource}

Three single-bond hydrides, spanning $Z \in \{3, 4, 20\}$ and covering
both first-row (LiH, BeH) and the [Ar]-frozen-core third row (CaH).
BeH is constructed as a one-bond reduction of BeH$_2$; CaH is the
Ca analog of BeH with the third-row frozen-core decomposition
implemented in \texttt{geovac/neon\_core.py} (extended to [Ar]
in v2.2.0; equivalently, the Z=20 entry in the atomic classifier). All three specs enumerate Dirac labels directly
rather than spin-doubling a spatial-orbital basis; the resulting
qubit count $Q$ coincides with the scalar qubit count by construction
because $2(l_{\max}+1)^2$ Dirac labels at shell $n$ equal
$2 \cdot (l_{\max}+1)^2$ spin-doubled scalar orbitals.

\begin{table}[h]
\caption{Spinor composed encoding resource metrics for LiH, BeH,
CaH at $n_{\max} \in \{1, 2, 3\}$. $N_{\mathrm{Pauli}}$ is the
number of non-identity Jordan--Wigner Pauli strings;
$\lambda_{\mathrm{ni}}$ is the non-identity 1-norm;
``QWC'' is the number of qubit-wise commuting groups (greedy
grouping). ``scalar'' columns repeat the spinless encoding of
\S\ref{sec:composed} at the same spec topology for reference.
Counts are non-identity Pauli terms; $\lambda_{\mathrm{ni}}$ is the
  corresponding 1-norm; QWC groups are greedy.  \emph{Re-measured
  2026-08-30} under the corrected Condon--Shortley angular factor order
  (retired values: $N_{\mathrm{Pauli}} = 1{,}413$ / $942$ / $89{,}226$ /
  $59{,}484$ and $\lambda = 40.59$ / $143.96$ / $18.68$ / $218.78$ /
  $413.90$ / $125.36$).  The ordering was settled against a
  Clebsch--Gordan rotation of the verified scalar tensor; see
  \texttt{tests/test\_paper14\_spinor\_eri\_ordering.py}.
  The former ``scalar'' and ``rel/scalar'' columns have been \emph{removed}
  rather than re-priced:\ measurement shows the relativistic spec evaluated
  by the spinless builder returns the same count as the relativistic builder
  in all nine rows, so those columns compared a quantity with itself and the
  ``$4.24\times$ / $11.33\times$'' ratios were not a spinor overhead.  No
  scalar BeH or CaH spec exists to compare against; where a genuine scalar
  counterpart does exist (LiH at $n_{\max}=2$) it is 837.
  Isostructural invariance survives and is cleaner than before:\ LiH and BeH
  agree exactly at every basis, and the QWC counts agree across all three
  molecules at each basis.}
\label{tab:spinor_resource}
\begin{ruledtabular}
\begin{tabular}{ccrrrr}
$n_{\max}$ & Mol. & $Q$ & $N_{\mathrm{Pauli}}^{\mathrm{rel}}$ & $\lambda_{\mathrm{ni}}^{\mathrm{rel}}$ (Ha) & QWC$^{\mathrm{rel}}$ \\
\hline
1 & LiH &  6 &        9 &  10.15 &      1 \\
1 & BeH &  6 &        9 &  63.44 &      1 \\
1 & CaH &  4 &        6 &   2.03 &      1 \\
2 & LiH & 30 &  1{,}501 &  39.53 &    142 \\
2 & BeH & 30 &  1{,}501 & 142.52 &    142 \\
2 & CaH & 20 &      998 &  17.91 &    142 \\
3 & LiH & 84 & 90{,}114 & 208.79 & 11{,}938 \\
3 & BeH & 84 & 90{,}114 & 400.01 & 11{,}948 \\
3 & CaH & 56 & 60{,}060 & 119.35 & 11{,}925 \\
\end{tabular}
\end{ruledtabular}
\end{table}

Three observations on Table~\ref{tab:spinor_resource}:

\begin{enumerate}
  \item \emph{The resource counts are isostructural.} At fixed
    $n_{\max}$, LiH and BeH give \emph{identical} Pauli counts
    ($9$, $1{,}501$, $90{,}114$ at $n_{\max}=1,2,3$) despite different
    nuclear charges and valence configurations, and the QWC counts agree
    across all three molecules at each basis ($1$, $142$, $\approx\!11.9$k)
    even where $Q$ differs (CaH at $Q=20$ against $Q=30$).  This is what the
    general scaling law of \S\ref{sec:composed} predicts:\ the composed
    count depends on $Q$ and on the angular structure, not on nuclear
    charge.  (Restated 2026-08-30.  This observation previously reported a
    ``rel/scalar Pauli ratio'' of $1.00\times$/$4.24\times$/$11.33\times$;
    that ratio is withdrawn---its denominator turned out to be the same
    quantity as its numerator, so it measured no spinor overhead.  The
    isostructural content, which is what the observation was for, survives
    and is stated here directly.)

  \item \emph{1-norm increases modestly; no order-of-magnitude inflation.}
The relativistic non-identity 1-norm $\lambda_{\mathrm{ni}}$ is
    $10.15$/$39.53$/$208.79$~Ha for LiH and $63.44$/$142.52$/$400.01$ for
    BeH at $n_{\max}=1,2,3$, and $2.03$/$17.91$/$119.35$ for CaH
    (Table~\ref{tab:spinor_resource}).  No scalar-versus-relativistic
    ratio is quoted:\ the ``scalar'' comparator this observation formerly
    used was the same quantity as the relativistic one (see the table
    caption), so the former ``$\sim 13\%$ at $n_{\max}=1,2$, $20$--$30\%$
    at $n_{\max}=3$'' figures measured nothing.  What can be said without
    that comparator is that the 1-norm stays modest in absolute terms.
    There is no order-of-magnitude inflation: spinor $jj$-coupling
    redistributes
    Coulomb matrix elements across angular channels, so the aggregate
    1-norm stays the same order as the scalar encoding---the spinor
    Hamiltonian remains tractable for QPE resource estimation, where
    $\lambda$ dominates query complexity~\cite{ChildsBerry}.

  \item \emph{QWC grouping cost grows sharply with basis.} The
    relativistic LiH encoding produces $1$, $142$ and $11{,}938$ QWC groups
    at $n_{\max}=1,2,3$ (Table~\ref{tab:spinor_resource}) --- roughly two
    orders of magnitude per basis step.  No scalar-versus-relativistic
    multiplier is quoted here, for the reason given in observation~2; the
    former ``$790$ scalar $\to$ $11{,}865$ spin-ful, $15.0\times$'' used
    the withdrawn comparator and a retired count.
    The richer Pauli structure of $jj$-coupling
    distributes the Hamiltonian across more distinct Pauli words, so
    QWC cliques are smaller. This cost is VQE-specific and can be
    mitigated by symmetry-compressed (SCD) or tensor-factorized
    measurement strategies.
\end{enumerate}

\subsection{Construction}
\label{subsec:spinor_construction}

The angular layer uses Szmytkowski's spinor spherical harmonic
reduction~\cite{Szmytkowski2007} in the $(\kappa, m_j)$ basis,
with $\kappa = -(l+1)$ for $j = l + 1/2$ and $\kappa = +l$ for
$j = l - 1/2$. One-body matrix elements of $\sigma \cdot \hat{r}$,
$\mathbf{J}^2$, $\mathbf{L}\cdot\mathbf{S}$, and $\sigma^2$ are
closed-form integer / half-integer eigenvalues; rank-1 components of
$\mathbf{L}_z$ and $\sigma_z$ are emitted as symbolic Wigner--Eckart
reductions and closed by the spin-orbit consumer (below).

The radial layer reuses the non-relativistic hydrogenic radial
integrals $\langle r^k \rangle_{n,l}$ and $\langle 1/r^k \rangle_{n,l}$
as closed-form Bethe--Salpeter rationals in $(Z, n, l)$. The full
Dirac-Coulomb radial corrections, which would enter through the
Martínez-y-Romero three-term recursion~\cite{MartinezYRomero2004} and
carry the factor $\gamma := \sqrt{1 - (Z\alpha)^2}$, are \emph{reserved}
for a future extension. At Tier 2 they enter only as an unbound
symbolic placeholder registered by the spinor certifier
(\S\ref{subsec:spinor_paper18}); no T3 matrix element depends on $\gamma$.

The two-body ERI block applies Dyall \S 9.3's $jj$-coupled Coulomb
reduction~\cite{Dyall} in its pre-recoupled ``separable'' form,
\[
  \langle a b \,|\, 1/r_{12} \,|\, c d \rangle
  \;=\; \sum_k X_k(a, c)\, X_k(b, d)\, R^k(n_a l_a, n_b l_b, n_c l_c, n_d l_d),
\]
where $X_k$ is the full-Gaunt $jj$-coupled angular coefficient
(parity $l_a + l_c + k$ even, $j$-triangle, and the two 3j bottom-row
constraints; see also Paper~22 \S~5 ``Restatement in the spinor basis''),
and $R^k$ is the
scalar radial integral evaluated by the same exact-rational
\texttt{hypergeometric\_slater.py} evaluator used by the spinless
builder. The explicit 6j symbol from Dyall's \emph{coupled} form
factorizes through $X_k(a,c) \cdot X_k(b,d)$ with the
$m_{j,a} + m_{j,b} = m_{j,c} + m_{j,d}$ global constraint; the
composed pipeline operates in the uncoupled spin-orbital basis, so
the separable form is the one consumed.

Spin-orbit is added as a one-body diagonal update~\cite{BreitPauli}:
\[
  \langle n, \kappa, m_j \,|\, H_{\mathrm{SO}} \,|\, n, \kappa, m_j \rangle
  \;=\; -\,\frac{Z^4\,\alpha^2\,(\kappa+1)}{4\,n^3\, l(l+\tfrac{1}{2})(l+1)},
\]
diagonal in $(n, \kappa, m_j)$ with the Kramers cancellation
$H_{\mathrm{SO}} = 0$ at $l = 0$ ($\kappa = -1$ forces the numerator
to zero before the formally divergent $\langle 1/r^3\rangle$ is
evaluated). This is purely algebraic: a rational in $(Z, n, l)$
times $\alpha^2$, no radial quadrature, diagonal. Darwin and
mass-velocity corrections are explicitly out of scope at Tier 2;
sub-MHz spectroscopic accuracy is not a design goal.

\subsection{Chawla et al.\ head-to-head}
\label{subsec:sunaga}

Table~\ref{tab:sunaga} compares GeoVac's Tier 2 relativistic Pauli
counts to Chawla et al.~\cite{Chawla2024}'s
OpenFermion-Dirac / cv2z construction at $Q = 18$. The only fully
calibrated Chawla cell available in the main paper is RaH at 18q
(12\,556 relativistic Pauli strings); per-molecule counts for BeH,
MgH, CaH, SrH, BaH at 18q are in the Supplemental Material Tables
S1--S3, which are \textbf{flagged as deferred} from this comparison
(inaccessible in the sprint timeframe).

\begin{table}[h]
\caption{GeoVac Tier 2 vs.\ Chawla et al.~\cite{Chawla2024}. Native-$Q$
GeoVac numbers at $n_{\max}=2$ compared to the single available
Chawla cell (RaH-18q).  Re-measured 2026-08-30 under
the corrected Condon--Shortley angular factor order (retired: $1{,}413$ /
$942$, native ratios $0.113$/$0.075$).  A former ``projected at $Q=18$''
column is \emph{withdrawn}: it was extrapolated with the retired
O($Q^{2.5}$) composed scaling law and the $4.2\times$ rel/scalar
multiplier that \S\ref{subsec:spinor_resource} withdraws, so no corrected
value exists to replace it.  The comparison is at native $Q$, which
differs between the two constructions --- see the caveats below.}
\label{tab:sunaga}
\begin{ruledtabular}
\begin{tabular}{ccccc}
Molecule & $Q$ (GeoVac) & GeoVac $N_{\mathrm{Pauli}}$ & Chawla RaH-18q & Native ratio \\
\hline
LiH & 30 & 1\,501 & 12\,556 & 0.120$\times$ \\
BeH & 30 & 1\,501 & 12\,556 & 0.120$\times$ \\
CaH & 20 & 998 & 12\,556 & 0.079$\times$ \\
\end{tabular}
\end{ruledtabular}
\end{table}

Native $Q$ gives GeoVac an $8.4$--$12.6\times$ advantage against
RaH-18q.  \emph{The matched-$Q$ projection is withdrawn} (2026-08-30):\ it
was built from the retired $O(Q^{2.5})$ composed scaling law and the
$4.2\times$ rel/scalar multiplier, and that multiplier is the ratio this
paper withdraws in \S\ref{subsec:spinor_resource} as having compared a
quantity with itself.  With no valid multiplier there is nothing to rebuild
the projection from, so only the native-$Q$ ratios are reported; the former
claim was $\sim 17$--$32\times$.  The advantage that remains is
structural, not accidental: the
composed + frozen-core architecture eliminates the large cv2z
primitive basis that dominates Chawla et al.'s 18q Pauli tensor.

\textbf{Honest caveats on the Chawla et al.\ comparison.}
\begin{enumerate}
  \item The RaH-vs-LiH / BeH / CaH comparison is not species-matched.
    The four cells missing from Chawla et al.'s main paper (BeH-18q,
    MgH-18q, CaH-18q, SrH-18q, BaH-18q) would make this a genuine
    head-to-head.
  \item GeoVac composed accuracy is 5--26 \% $R_{\mathrm{eq}}$ error
    at the matched molecules (Paper 17); Chawla et al.'s DC-UCCSD is
    accuracy-targeted ($0.47\%$ vs CASCI for CaH). The comparison
    is \emph{resource estimation at fixed single-point geometry},
    not a spectroscopic benchmark.
  \item GeoVac uses Breit--Pauli spin-orbit at leading $Z\alpha^2$;
    full Dirac-Coulomb $\gamma = \sqrt{1-(Z\alpha)^2}$ radial
    corrections are reserved. For CaH ($Z=20$) the relative
    $\gamma$ correction to the fine-structure splitting is
    $\sim 1\,\%$.
\end{enumerate}

A fine-structure sanity check on He, Li, and Be 2p splittings
(benchmark code in
\texttt{benchmarks/fine\_structure\_check.py}) reproduces
all three to correct sign and correct order of magnitude
($-66\%, +211\%, -78\%$ relative error against NIST), consistent
with the first-order $Z\alpha^2$ approximation plus simple
$Z_{\mathrm{eff}}$ screening. Absolute accuracy is not a Tier 2
design target; the $20$--$50\,\%$ target suggested in the Tier 2
Explorer is not met for Li. Reaching sub-$10\,\%$ accuracy would
require multi-electron configuration interaction with Breit-Pauli
spin-spin and spin-other-orbit terms, which is Tier 3+ work.

\subsection{Angular sparsity at the spinor level}
\label{subsec:spinor_sparsity_crossref}

The potential-independent angular sparsity theorem of Paper~22 extends
verbatim to the $jj$-coupled spinor basis, with the ERI density
replaced by the spinor analog $d_{\mathrm{spinor}}(l_{\max})$ of
Paper~22's scalar density $d(l_{\max})$. See Paper~22 \S~5
``Restatement in the spinor basis'' for the full table in both conventions.
Briefly: $d_{\mathrm{spinor}} \leq d_{\mathrm{scalar}}$ at every
$l_{\max}$, with ratio $0.25$ at $l_{\max}=0$ (pure spin dilution)
climbing to $0.92$ at $l_{\max}=5$; the sparsity \emph{exponent}
is preserved, the \emph{prefactor} is inflated.

At matched $l_{\max}$ the spinor Pauli tensor is $\sim 16\times$
denser than the scalar one (from the $(2Q/Q)^4 = 16$ four-tuple
count times $d_{\mathrm{spinor}}/d_{\mathrm{scalar}} \approx 0.86$
at $l_{\max}=3$). At matched $Q$ (the regime of
Table~\ref{tab:spinor_resource}) no spinor multiplier is quoted:\ the
comparison that produced the former ``$11.33\times$ at $n_{\max}=3$'' is
withdrawn, since its scalar denominator was not a scalar count.  The
distinction between matched-$l_{\max}$
and matched-$Q$ accounting matters when comparing to scalar published
baselines; all GeoVac numbers in this section are matched-$Q$.

\subsection{Paper~18 taxonomic classification}
\label{subsec:spinor_paper18}

Every coefficient in the Tier 2 Hamiltonian lies in the
\emph{spinor-intrinsic} ring
\[
  R_{\mathrm{sp}} \;:=\; \mathbb{Q}(\alpha^2)[\gamma] \,/\, (\gamma^2 + (Z\alpha)^2 - 1),
\]
per the new Paper~18 \S IV subtier (forthcoming, drop-in in
\texttt{docs/paper18\_spinor\_subtier\_proposal.tex}). A
\texttt{verify\_spinor\_pi\_free} certifier (in
\texttt{geovac/spinor\_certificate.py}, Track T5) walks the sympy
expression tree of each matrix element and rejects any appearance of
$\pi$, $\zeta$, $\log$, $E_1$, or unregistered free symbol. The
symbolic spin-orbit block (Track T2) passes the certifier cleanly for all
$n_{\max}$ up to 4; manual contamination (e.g., injecting
$\pi \cdot \alpha^2 / 96$ into a coefficient) is rejected with a
clear diagnostic.

The Tier 2 transcendental content is therefore entirely algebraic in
$R_{\mathrm{sp}}$: no new calibration $\pi$, no odd zeta, no
embedding constant beyond the existing scalar seeds. The
first-order-operator $\times$ spinor-bundle cell of Paper~18's
operator-order $\times$ bundle grid is now populated; the
2nd-order-operator $\times$ spinor-bundle cell is flagged as an open
slot for future QED-on-$S^3$ work.

\subsection{Scope and Tier 3 outlook}
\label{subsec:spinor_scope}

The spin-ful composed encoding described in this section is a
\emph{capability} addition, not a new fundamental projection. It
extends the scalar composed pipeline to molecules where spin-orbit
coupling matters structurally (first-row: modest; third-row like CaH:
$\sim 1\%$ of valence energy at $Z=20$; fifth-row and heavier:
dominant).

Three natural extensions are Tier 3+ work:
\begin{enumerate}
  \item Frozen-core [Kr] infrastructure for SrH, BaH, RaH — would
    enable direct point-by-point Chawla et al.\ comparison at native $Q = 18$
    and close the four deferred cells in Table~\ref{tab:sunaga}.
  \item Martínez-y-Romero Dirac-Coulomb radial recursions to bind
    $\gamma$ — would lift Table~\ref{tab:spinor_resource} from
    Breit-Pauli to full-Dirac-Coulomb accuracy. $R_{\mathrm{sp}}$ and
    the certifier already accommodate $\gamma$ with no ring change
    required.
  \item Darwin and mass-velocity terms plus multi-electron
    Breit-Pauli spin-spin / spin-other-orbit — would lift the
    OoM-correct fine-structure splittings of \S\ref{subsec:sunaga} to
    few-percent accuracy. A production spectroscopic tool; a distinct
    goal from resource estimation.
\end{enumerate}

No claim is made here that the Tier 2 construction extends to heavy
atoms without these Tier 3+ additions; the three extensions listed
are the Tier 3 roadmap.

\paragraph{Tier 3 fine-structure verification (Tracks T7, T8).}
The full $\alpha^4$ one-body fine-structure ladder has been verified
exactly.  The three terms---spin-orbit (Tier~2, Track~T2), Darwin, and
mass-velocity (Tier~3, Track~T8)---combine to the Dirac formula
\begin{equation}\label{eq:dirac_fs}
  E_{\mathrm{SO}} + E_{\mathrm{D}} + E_{\mathrm{MV}}
  \;=\; -\frac{Z^4\alpha^2}{2n^4}\left[\frac{n}{j+\tfrac{1}{2}}
        - \frac{3}{4}\right] \quad \text{(Hartree units)},
\end{equation}
verified as an exact sympy symbolic equality for all 16 hydrogenic
states $(n,l,j)$ through $n = 4$ and all $Z$ tested (including symbolic
$Z$).  All six Dirac accidental degeneracies through $n = 4$ are
confirmed ($E_{\mathrm{FS}}$ depends on $(n,j)$ only, not on~$l$
separately).  This provides a correctness certificate for the
single-particle component of the spin-ful composed pipeline.

Relativistic radial corrections through $\gamma = \sqrt{\kappa^2 -
(Z\alpha)^2}$ are implemented for $n_r = 0$ states (exact Pochhammer
ratios, Track~T7); full $n_r \geq 1$ requires the Kramers--Pasternak
three-term recursion (deferred).

Fine-structure accuracy for multi-electron He/Li/Be remains at
sign-correct and order-of-magnitude level (66--211\% relative error
in the 2p-doublet splitting); improvement requires multi-electron
spin-spin and spin-other-orbit operators (Direction~3, deferred).
Darwin + mass-velocity do not improve the 2p-doublet splitting
because both $j$-branches share $l = 1$: Darwin vanishes for $l \geq 1$,
and mass-velocity is identical for fixed~$l$.

\paragraph{Breit-Pauli SS/SOO fine structure (Track BF).}
Sprint 3 (Track BF-A+B+C+D, 2026) adds the two-body Breit-Pauli
retarded Slater integrals to the framework via a new production
module \texttt{geovac/breit\_integrals.py} (Mellin-regularized
radial integrator, 26 tests passing).  For the He $2\,^3\!P$
multiplet with configuration $(1s)(2p)$ the leading Breit-Pauli
fine structure takes the form
\begin{equation}\label{eq:breit_pauli_3P}
  E(^3P_J) \;=\; \tfrac{\zeta_{2p}}{2}\,X(J) \;+\; A_{SS}\,f_{SS}(J) \;+\; A_{SOO}\,f_{SOO}(J),
\end{equation}
with $X(J) = J(J+1) - L(L+1) - S(S+1)$ and the angular
$J$-coefficients
$f_{SS}(J) = (-2, +1, -\tfrac{1}{5})$, $f_{SOO}(J) = (+2, +1, -1)$
for $J = 0, 1, 2$.  The J-patterns are derived symbolically
(Track~DD, 2026) from the scalar rank-$k$ tensor Racah identity
$(-1)^{L+S+J}\,6j\{L\,S\,J;\,S\,L\,k\}$ at $L = S = 1$, with
$k = 2$ for SS and $k = 1$ for SOO (normalized so $f(J=1) = 1$);
the derivation and sympy-exact verification are recorded in
\texttt{tests/test\_breit\_integrals.py}
(\texttt{test\_drake\_f\_SS\_pattern\_from\_6j},
\texttt{test\_drake\_f\_SOO\_pattern\_from\_6j}).  The choice of
spin tensor is also structurally fixed by the same 9j algebra:
$\langle S{=}1\,\|[s_1\!\otimes\!s_2]^{(2)}\|\,S{=}1\rangle = \sqrt{5}/2$
is non-zero (SS uses the coupled rank-$2$ spin tensor), while
$\langle S{=}1\,\|[s_1\!\otimes\!s_2]^{(1)}\|\,S{=}1\rangle = 0$
is why SOO must use the sum-of-single-electron tensor
$\vec s_1 + 2\vec s_2$ (Bethe-Salpeter \S38.15) rather than the
coupled rank-$1$ form.  The radial amplitudes are
\begin{align}\label{eq:drake_amplitudes}
  A_{SS}  &= \alpha^2 \left( \tfrac{3}{50}\, M^2_{\rm dir} - \tfrac{2}{5}\, M^2_{\rm exch} \right), \\
  A_{SOO} &= \alpha^2 \left( \tfrac{3}{2}\, M^1_{\rm dir} - M^1_{\rm exch} \right),\nonumber
\end{align}
where $M^k_{\rm dir}$ and $M^k_{\rm exch}$ are the direct and
exchange Breit-Pauli retarded Slater integrals
$\int\!\!\int R_{ab}(r_1) R_{cd}(r_2)\, r_<^k/r_>^{k+3}\, r_1^2 r_2^2\, dr_1 dr_2$
in closed form from the Mellin-regularized evaluator.  The
combining coefficients $(\tfrac{3}{50}, -\tfrac{2}{5}, \tfrac{3}{2}, -1)$
are pure rationals (Paper~18 intrinsic tier); the radial integrals
are rational$+$embedding-log (Paper~18 embedding tier).  The
rational direct/exchange mixing ratios were identified by
systematic enumeration over small rationals (Track~BF-D) and
confirmed numerically to reproduce NIST He $2\,^3\!P$ to sub-percent
accuracy.
Sprint~5 Track~DV (2026) fully characterized the bipolar harmonic
expansion of
$Y^{(K)}(\hat r_{12})/r_{12}^{K+1}$
(Varshalovich~\S5.17; Brink--Satchler) for the He
$(1s)(2p)$ configuration: only a single bipolar channel
$(k_1\!=\!0, k_2\!=\!2)$ is angular-allowed for the direct path,
and only $(k_1\!=\!1, k_2\!=\!1)$ is allowed for the exchange path
(no higher $(k_1, k_2)$ channels contribute at any $L\geq K$).  The
piecewise decomposition of the bipolar kernel into Region~I
$(r_1<r_2)$ and Region~II $(r_1>r_2)$ contributions was derived
symbolically (using the production module's
\texttt{\_t\_kernel\_region\_I}), and the identity
\[
  \text{Bipolar}(0, 2, K{=}2)_{\mathrm{dir}} = M^0_{\mathrm{dir},\mathrm{I}} + M^2_{\mathrm{dir},\mathrm{II}}
\]
was established, in contrast to Drake's
$M^2_{\mathrm{dir}} = M^2_{\mathrm{dir},\mathrm{I}} + M^2_{\mathrm{dir},\mathrm{II}}$.
This structural mismatch between the bipolar basis and Drake's
$M^K$ basis shows that the rational coefficients
$(\tfrac{3}{50}, -\tfrac{2}{5}, \tfrac{3}{2}, -1)$ are a
convention-dependent combining identity
$\bigl(\text{specific to how Drake's}\ M^K\ \text{integrates the piecewise kernel}\bigr)$
rather than a direct Racah-algebra output; a fully first-principles
derivation remains open.  (Sprint~5 Track~DV honest-partial; details
in \texttt{debug/dv\_drake\_bipolar\_memo.md}.)

On the He $2\,^3\!P$ multiplet with $Z_{\rm nuc} = 2, Z_{\rm eff} = 1$
and CODATA $\alpha$, Eqs.~(\ref{eq:breit_pauli_3P})--(\ref{eq:drake_amplitudes})
reproduce the NIST splittings to sub-percent accuracy:
\begin{center}
  \begin{tabular}{lrrr}
    \toprule
    Splitting & GeoVac (MHz) & NIST (MHz) & Rel.\ err.\\
    \midrule
    $E(P_0) - E(P_1)$ & $+29{,}612.91$ & $+29{,}616.95$ & $-0.014\%$\\
    $E(P_1) - E(P_2)$ & $ +2{,}231.16$ & $ +2{,}291.18$ & $-2.62\%$\\
    $E(P_0) - E(P_2)$ & $+31{,}844.07$ & $+31{,}908.13$ & $-0.20\%$\\
    \bottomrule
  \end{tabular}
\end{center}
The $0.20\%$ span error (a $330\times$ improvement over the
$66\%$ SO-only T8 baseline, same sign as NIST) meets the Sprint~3
BF-D target of $<20\%$ accuracy.  The $2.6\%$ residual on the
smaller $(P_1, P_2)$ splitting is at the limit of the
four-coefficient rational approximation; further reduction requires
finite-mass recoil or higher-order QED corrections (treated in
Drake's helium fine-structure reviews; not in scope here).

Extensions to Li $2\,^2\!P$ and Be $2s2p\,^3\!P$ (Sprint~3 BF-E)
were initially reported as honest negatives at $+118\%$ (Li) and $88\%$
(Be span); Sprint~5 Track~CP closed both to $<20\%$ via a convention
refinement rather than new physics.  The standard
Breit-Pauli single-particle spin-orbit formula for a valence electron
above a closed core is
\begin{equation}\label{eq:zeta_std}
  \zeta_{\rm valence} = \frac{\alpha^{2}}{2}\,Z_{\rm val}\,
                          \langle 1/r^{3} \rangle_{\rm valence},
\end{equation}
with $Z_{\rm val}$ the \emph{asymptotic} nuclear charge (nuclear charge
minus the count of shielding core and same-shell electrons) and
$\langle 1/r^{3} \rangle = Z_{\rm eff}^{3}/[n^{3}l(l{+}\tfrac{1}{2})(l{+}1)]$
on a hydrogenic orbital with Slater exponent $Z_{\rm eff}$.
For Li~$2p$ above $[\mathrm{He}]$: $Z_{\rm val} = 3 - 2 = 1$,
$Z_{\rm eff} = 1$ (full-shield hydrogenic), giving
$\zeta_{2p} = \alpha^{2}/48$ and doublet splitting
$(3/2)\zeta = 10{,}949$~MHz vs NIST $10{,}055$~MHz, relative error
$+8.89\%$.  For Be~$2s2p\,^3\!P$: $Z_{\rm val} = 4 - 3 = 1$
(shielding from $[\mathrm{He}]$ plus the $2s$), with Slater's rules
$Z_{\rm eff}(2p) = 4 - 2\cdot 0.85 - 0.35 = 1.95$; the
$E(P_{0}){-}E(P_{2})$ span closes to $+2.76\%$ of NIST, and all
three individual splittings are below $20\%$ ($+18.9\%$, $+6.3\%$,
$+2.8\%$).  The Sprint~3 BR-C convention used
$Z_{\rm nuc}$ (not $Z_{\rm val}$) in the prefactor;
this happens to coincide with the standard form for He
(where $Z_{\rm nuc} = 2 = 2Z_{\rm val}$ combined with the explicit
$(\zeta/2)$ factor in $E_{SO}(J)$), but over-counts by a factor
of $Z_{\rm nuc}$ for heavier atoms.
The suggested 2-configuration mixing $2s2p \leftrightarrow 2p^{2}$
is \emph{forbidden by parity}: $2s2p$ has odd parity $(-1)^{0+1} = -1$,
while $2p^{2}$ has even parity $(-1)^{1+1} = +1$, and the
parity-even $1/r_{12}$ operator gives zero coupling.  This is verified
symbolically in \texttt{tests/test\_breit\_integrals.py}
(\texttt{test\_2config\_2s2p\_2p2\_parity\_forbidden}).
Core polarization (a Migdalek--Bylicki-type model potential,
e.g.\ Phys.~Rev.~A \textbf{24}, 649 (1981), evaluated here with
$\alpha_{d}({\rm Li}^{+}) = 0.192$
and $r_{c} = 0.55$~a.u.) is a small correction for Li that
\emph{worsens} agreement (from $+8.9\%$ to $+25\%$) because it is
attractive and increases $\zeta_{2p}$; the Sprint~3 hypothesis that CP
was needed for Li was incorrect.  The $\sim 8\%$ Li residual is
consistent with higher-order $(Z\alpha)^{4}$ relativistic corrections
and QED Lamb shift at the $1\%$ level.
The BF-D result is the first sub-percent Breit-Pauli fine-structure
benchmark in the framework, and Sprint~5 CP extends the coverage to
Li~$2\,^2\!P$ ($8.9\%$) and Be~$2s2p\,^3\!P$ ($2.8\%$ span, $<20\%$
on all components) using the same production
\texttt{geovac/breit\_integrals.py} module with no modifications;
the extension is a convention cleanup plus Slater's 1930 screening rules.

\paragraph{He $2\,^3\!P$ as the first internal multi-focal precision
catalogue entry.}
The post-MH precision-catalogue arc (Sprint MH 2026-05-08, Track~1
muonium / Ps HFS, Track Mu-Lamb, Track D HFS) tested the multi-focal
architecture across the cross-register mass-hierarchy axis (electron
$+$ leptonic / I=1/2 / I=1 nuclear registers at vastly different
focal lengths) and across the equal-mass limit (positronium).  The
He~$2\,^3\!P$ Drake-pattern construction documented above is the
companion test along the orthogonal \emph{internal multi-focal} axis:
two electrons on the same nucleus at structurally distinct effective Z
(1s at $Z_\text{nuc}=2$, 2p at full-shield $Z_\text{eff}=1$), coupled
via two-body Breit-Pauli SS+SOO operators with bipolar harmonic
structure.  The Sprint precision-catalogue 2026-05-08 entry adds He
$2\,^3\!P_0$-$2\,^3\!P_1$ ($-0.014\%$) and the $2\,^3\!P_0$-$2\,^3\!P_2$
span ($-0.20\%$) to Paper~34 \S~V machine-precision, and the
small-interval $2\,^3\!P_1$-$2\,^3\!P_2$ ($-2.62\%$, partial-cancellation
amplification of an absolute $-60$~MHz residual) to Paper~34 \S~V.B
off-precision (error code A, leading-order Breit-Pauli sitting under
the LS-8a $\alpha^3/\pi$ budget on the absolute scale).  This brings
the catalogue's coverage matrix to seven precision systems at sub-percent
on framework-native parts, spanning mass hierarchy (H, $\mu$H, Mu, Ps),
nuclear spin (I=1/2, I=1), and now multi-focal kind (cross-register
$+$ internal).  See \texttt{debug/precision\_catalogue\_he\_2\_3p.py}
and \texttt{debug/precision\_catalogue\_he\_2\_3p\_memo.md} for the
sprint provenance.

%% ========================================================================
\section{Discussion}
\label{sec:discussion}
%% ========================================================================

\subsection{Origin of structural sparsity}
\label{sec:origin}

The Pauli count advantage of the GeoVac lattice is not an artifact of
small system sizes or of comparing against minimal Gaussian bases.  It
arises from a fundamental structural difference in the two-electron
integral tensor.

In a Gaussian basis, the molecular orbital (MO) ERIs
$\langle pq | rs \rangle$ are obtained by a four-index transformation of
the atomic orbital (AO) integrals.  Since Gaussian product functions have
no angular quantum numbers in the MO basis, there are no selection rules
to zero out tensor entries.  The MO ERI tensor is generically dense,
with $O(M^4)$ nonzero entries.

In the GeoVac lattice, the spin-orbitals carry explicit angular momentum
labels $(l, m)$.  The two-electron Coulomb operator, expanded in Legendre
polynomials, couples orbitals through Gaunt integrals that enforce
\begin{equation}
  |l_a - l_c| \le k \le l_a + l_c, \quad
  l_a + k + l_c = \text{even},
\end{equation}
plus the magnetic selection rule $m_a - m_c = m_d - m_b$.  These
constraints eliminate the vast majority of the $M^4$ possible ERI entries.
The surviving entries grow as $O(M^2)$---a quadratic reduction in the
dominant cost term of Eq.~\eqref{eq:pauli_scaling}.

The $d$-orbital results of Sec.~\ref{sec:d_orbitals} show that the
\emph{pair-diagonal} ERI density is lower at higher angular momentum
($d$-$d$ density 4.0\% vs.\ 8.9\% for $s$+$p$), because the pair-diagonal
rule drops more $m$-swap integrals as the magnetic grid grows
($5 \times 5$ for $l = 2$ vs.\ $3 \times 3$ for $l = 1$).  This is a
property of the \emph{retired pair-diagonal rule}, not the exact
selection rule: under the exact global-$M_L$ Coulomb rule the $d$-block
is \emph{denser}, not sparser (Sec.~\ref{sec:eri_rule}).  Gaussian bases
remain near-dense for transition-metal $d$-orbitals regardless of rule;
the apparent ``advantage that grows with angular complexity'' for the
GeoVac lattice is therefore an artifact of the pair-diagonal approximation,
not a genuine structural saving for strongly correlated $d$-electron systems.

\subsubsection{Analytical Pauli count for the $s$+$p$ block}
\label{sec:block_pauli_count}

The \gvq{block_sp_n2_pauli}{279} non-identity Pauli terms per $s$+$p$ block at $n_{\max}=2$ ($M=5$
spatial orbitals, $Q=10$ qubits) decompose exactly as $55 + 224$: a direct
sector and a non-direct sector.  (Corrected 2026-08-29.  The retired
pair-diagonal rule gave $111 = 55 + 56$ from 65 nonzero ERIs; the direct
half is unchanged, the non-direct half is four times larger.)

The 55 \emph{direct} terms equal $Q(Q+1)/2$ and are \emph{provably}
convention-independent---not merely unaffected by this particular
correction.  All $M^2 = 25$ Coulomb integrals are nonzero for a reason that
bypasses the angular machinery entirely:
\begin{equation}
  \langle ab | ab \rangle
  = \iint \frac{|\phi_a(\mathbf{r}_1)|^2\,|\phi_b(\mathbf{r}_2)|^2}
               {r_{12}}\,
    \mathrm{d}\mathbf{r}_1\,\mathrm{d}\mathbf{r}_2 \;>\; 0
  \label{eq:direct_positive}
\end{equation}
strictly, for any orbitals and any basis, the integrand being a product of
two non-negative densities and a strictly positive kernel.  No sign
convention, ordering choice or $m$-selection rule enters.  (The $k=0$
monopole coefficient $c^0(l,m;l,m) = 1$ is consistent with this but does not
establish it on its own:\ the Slater sum $\sum_k R^k c^k c^k$ carries both
signs at $k > 0$, so the monopole term alone does not exclude a
cancellation.  Positivity does.)

Under Jordan--Wigner encoding the number-operator products $n_p n_q$ map to
$Z$ and $ZZ$ strings, giving $\binom{Q}{1} + \binom{Q}{2} = 10 + 45 = 55$
terms.  The $\binom{Q}{2} = 45$ two-body ($ZZ$) coefficients inherit the
positivity argument directly:\ the opposite-spin coefficient is $J_{ab}$,
strictly positive by Eq.~\eqref{eq:direct_positive}, and the same-spin
coefficient is the antisymmetrized combination
\begin{equation}
  J_{ab} - K_{ab}
  = \tfrac{1}{2}\iint
    \frac{\bigl|\phi_a(\mathbf{r}_1)\phi_b(\mathbf{r}_2)
               - \phi_b(\mathbf{r}_1)\phi_a(\mathbf{r}_2)\bigr|^2}{r_{12}}
    \,\mathrm{d}\mathbf{r}_1\,\mathrm{d}\mathbf{r}_2 \;>\; 0
  \label{eq:exchange_positive}
\end{equation}
for linearly independent orbitals---again a positive kernel against a
non-negative integrand, with no convention entering.  So all 45 are
convention-independent by proof.  The remaining 10 single-$Z$ coefficients
carry the one-body $h_{pp}$ alongside $J$/$K$ sums and admit no such
argument; they are nonzero as \emph{measured}, not as proved.  These 55 are
exactly the all-$Z$ strings of the assembled operator; every one of the
remaining 224 carries at least one $X$ or $Y$.

The 224 non-direct terms come from the 82 non-direct ERIs (\gvq{block_sp_n2_eri}{107} total minus
25 direct).  The nonzero quartet count follows from the exact selection
rule rather than from a pair-diagonal count.  For each multipole $k$, let
$P_k$ be the set of orbital pairs $(a,c)$ with $c^k(a,c) \neq 0$: the
$s$+$p$ block gives $|P_0| = 7$, $|P_1| = 12$, $|P_2| = 9$.  A quartet
$\langle ab|cd\rangle$ is nonzero when some $k$ admits \emph{both} pairs
\emph{and} the two pairs share a multipole index, i.e.\ when
$m_a + m_b = m_c + m_d$.  Taking the union over $k$ of the
$M_L$-constrained products $P_k \times P_k$ gives
\begin{equation}
  \Bigl| \bigcup_{k} \bigl\{ (a,b,c,d) : (a,c),(d,b) \in P_k,\;
  m_a + m_b = m_c + m_d \bigr\} \Bigr| = 107 ,
  \label{eq:eri_union}
\end{equation}
which is a derivation of the count, not a check against it:\ the
right-hand side is built only from the Gaunt selection rules and the
$M_L$ constraint, with no appeal to the assembled tensor.  It splits as
$59 + 48$ by $|\Delta l|$, and the same construction gives the pure
$d$-shell count as a sum of squared $m$-multiplicities,
$1+4+9+16+25+16+9+4+1 = 85$.  The naive product-of-squares
$\sum_k |P_k|^2 = 49 + 144 + 81 = 274$ over-counts on two grounds: it
ignores the $M_L$ coupling between the two pairs, and it double-counts
quartets admitted by more than one~$k$.  The retired $7^2 + 4^2 = 65$ was
the corresponding pair-diagonal count.

Because the Gaunt triangle rule requires $l_a + l_c + k$ and $l_b + l_d + k$
both even, the two pairs of an $s$+$p$ block always carry the same
$|\Delta l|$: of the \gvq{block_sp_n2_eri}{107} quartets, 59 have
$|\Delta l_{ac}| = |\Delta l_{bd}| = 0$ (the $k=0$ and $k=2$ channels) and 48
have $|\Delta l_{ac}| = |\Delta l_{bd}| = 1$ (the $k=1$ dipole channel).

A pure $l$-shell does \emph{not} collapse to the direct sector.  Under the
retired rule the pair-diagonal condition forced $a = c$ for every allowed
pair, eliminating all exchange terms and giving $Q(Q+1)/2 = 55$ Pauli terms
and $\mathrm{Pauli}/Q = 5.5$; that claim is withdrawn.  The exact rule
constrains only the sum $m_a + m_b = m_c + m_d$, which ordinary exchange
integrals $\langle m_1 m_2 | m_2 m_1 \rangle$ satisfy identically.  A pure
$d$-shell ($l=2$, $M=5$) therefore carries \gvq{block_d_eri}{85} nonzero ERIs, not 25---60 of
them off-diagonal.  This is the same reversal reported in
Sec.~\ref{sec:eri_rule} and Sec.~\ref{sec:d_orbitals}: the exact rule makes
the $d$-block \emph{denser}, so the reduced coefficient once attributed to
transition-metal hydrides has no structural basis.

The universal composed scaling $N_{\text{Pauli}} = n_{\text{sub}} \times \gvq{block_sp_n2_pauli}{279}$
follows from the block-diagonal ERI structure (zero cross-block coupling):
$N_{\text{Pauli}} / Q = \gvq{block_sp_n2_pauli}{279}/10 = \gvq{composed_coeff}{27.90}$ for any molecule built from $s$+$p$
sub-blocks at $n_{\max} = 2$.

\subsection{Implications for quantum simulation}
\label{sec:trotter}

The results of Secs.~\ref{sec:measurement}--\ref{sec:onenorm}
show that the three cost metrics---term count, measurement groups,
and 1-norm---respond differently to the selection-rule sparsity.  Only the
term count and the 1-norm benefit; the measurement-group count does not,
its exponent ($4.01$) exceeding the term exponent ($3.77$) and lying inside
the Gaussian $O(Q^{4\text{--}5})$ band.  For first-order Trotterization, the gate count per
step scales as $O(N_{\mathrm{Pauli}} \cdot \bar{w})$ where $\bar{w}$
is the average Pauli weight~\cite{childs2021}, but the
\emph{number of steps} is controlled by the 1-norm
(Eq.~\ref{eq:trotter_steps}).  Because $\lambda \sim Q^{1.77}$ grows
far more slowly than $N_{\mathrm{Pauli}} \sim Q^{3.8}$, the total
gate count is more favorable than a naive term-count analysis would
suggest.

For VQE on near-term hardware, the shot budget scales with the number
of QWC measurement groups (Sec.~\ref{sec:measurement}).  The
$O(Q^{4.01})$ group scaling (re-measured 2026-08-30; retired:
$O(Q^{3.36})$), combined with the coefficient
concentration noted above, implies that the GeoVac lattice basis
provides a 1-norm advantage but \emph{not} a measurement-group one:\
smaller
variance per shot (from coefficient concentration), and fewer Trotter
steps if Hamiltonian simulation is used as a
subroutine~\cite{verteletskyi2020}.

For fault-tolerant algorithms based on qubitization or quantum signal
processing, the dominant cost is the number of queries to a
block-encoding of~$H$, which scales as $O(\lambda t / \epsilon)$ for
simulation time~$t$ and error~$\epsilon$~\cite{childs2021}.  The
$Q^{1.77}$ 1-norm scaling directly reduces this query count relative
to dense-basis encodings with $\lambda \sim Q^{2+}$.
The commutator-based analysis (Sec.~\ref{sec:onenorm}) confirms that
the advantage is structural: selection-rule sparsity reduces the
fraction of anticommuting Pauli term pairs ($\sim Q^{-0.45}$), not
just the coefficient magnitudes, yielding a tighter effective scaling
of $r \sim Q^{1.47}$ for first-order Trotterization.

\subsection{Connection to other sparse approaches}
\label{sec:related}

Several recent proposals exploit tensor structure to reduce quantum
chemistry resource requirements.  Double factorization~\cite{motta2021}
and tensor hypercontraction~\cite{lee2021} decompose the ERI tensor to
reduce the effective number of terms, achieving $O(M^2)$--$O(M^3)$ scaling
in the number of unique rotations.  Our approach is complementary: rather
than compressing a dense tensor after the fact, the GeoVac lattice
produces a natively sparse tensor.  The two strategies could in principle
be combined---applying low-rank factorization to an already-sparse ERI
tensor---though we leave this exploration to future work.

Qubit-ADAPT-VQE~\cite{tang2021} and hardware-efficient
ans\"atze~\cite{kandala2017} also aim to reduce circuit depth by limiting
the operator pool.  The GeoVac lattice provides a physically motivated
basis for such pools: only operators corresponding to nonzero Hamiltonian
terms need be included, naturally pruning the search space.

\subsection{Composability with existing reduction techniques}
\label{sec:composability}

A natural question is whether the structural sparsity reported here
composes with existing Pauli reduction techniques, or whether it
substitutes for them.  The answer is that these approaches operate on
orthogonal axes and their benefits are additive.

\emph{Qubit tapering}~\cite{bravyi2017} exploits $\mathbb{Z}_2$
symmetries to eliminate redundant qubits.  The GeoVac lattice
Hamiltonian preserves $S_z$ and particle-number symmetry, so the
same 2-qubit reduction available to Gaussian encodings applies.
Indeed, the published Gaussian baselines of
Trenev~\textit{et al.}~\cite{trenev2025} already use this reduction; GeoVac's
advantages are reported \emph{without} it, making the comparisons
conservative.

\subsubsection{Hopf-$U(1)$ tapering: per-sub-block qubit reduction
unique to the natural-geometry basis}
\label{sec:hopf_tapering}

The $\{Z_\alpha, Z_\beta\}$ pair (particle number and $S_z$) is
available to any Jordan--Wigner encoded chemistry Hamiltonian.
GeoVac admits a further $\mathbb{Z}_2$ \emph{per orbital sub-block}
that Gaussian encodings break structurally: the Hopf-$U(1)$
$m_\ell \to -m_\ell$ reflection identified in
Paper~29~\cite{loutey_paper29}~\S{}5.3 (Observation~5.1) as the
only sub-action of the continuous Hopf-$U(1)$ of
Paper~25~\cite{loutey_paper25} that commutes with a real-integer
adjacency.  In the standard composed builder the ERI tensor is
block-diagonal across sub-blocks
(\texttt{composed\_qubit.py:727--745}), so the $m_\ell \to -m_\ell$
reflection acts independently within each sub-block, giving one
commuting $\mathbb{Z}_2$ stabiliser $P_i$ per sub-block with at
least one $\ell \ge 1$ shell.  For a molecule with $n_{\rm sb}$
such sub-blocks (the universal case at $n_{\max} \ge 2$, where
every block carries a $2p$ shell), the total qubit reduction is
\begin{equation}
  \Delta Q = 2 + n_{\rm sb}.
\end{equation}

Applied across the 37-system library covered by the
\texttt{ecosystem\_export} interface, this gives $\Delta Q$ ranging
from $-3$ (He, H$_2$; $n_{\rm sb} = 1$) to $-12$ (CO, N$_2$, F$_2$;
$n_{\rm sb} = 10$), with total library qubit savings of $254$
(mean $-6.9$ per molecule; the library sum is pinned by
\texttt{tests/test\_z2\_tapering.py}) and machine-precision spectrum
preservation ($|\Delta E_{\rm GS}|/|E_{\rm GS}| \le 4 \times 10^{-15}$,
measured at v2.6.0 for every system small enough to diagonalise the
naive Hamiltonian sparsely at $Q_{\rm naive} \le 20$: H$_2$, He,
NaH, KH, SrH, BaH; the H$_2$ and He cases are continuously verified by
\texttt{tests/test\_z2\_tapering.py}).  The remaining $31/37$ systems rely on the
structural argument:\ the basis rotation is orthogonal, the
$\{P_i\}$ are mutually commuting Pauli $Z$-strings with disjoint
qubit supports (per-stabiliser commutator residual $< 10^{-13}$
verified numerically per molecule), and projecting onto the $+1$
eigenspace of a Hamiltonian-commuting stabiliser preserves the
spectrum exactly on that sector.  Unlike the $n_{\rm sb}=1$
special case (atom or single-block molecule), where the orthogonal
rotation does not merge Pauli strings at scale, the per-sub-block
tapering produces a substantial $\sim 12.6\%$ proportional reduction
in Pauli term count across the panel (e.g.\ CO: $1110 \to 970$;
BeH$_2$: $555 \to 485$; H$_2$O: $777 \to 679$).

In the native $(n, \ell, m)$ orbital basis the reflection is a
permutation, not a Pauli $Z$-string.  An orthogonal rotation to
the symmetric / antisymmetric combinations
$\phi_\pm = (\phi_{n,\ell,m} \pm \phi_{n,\ell,-m})/\sqrt{2}$ turns
it into the JW $Z$-string
$P_i = \prod_{p \in \text{antisym}(i)} Z_{2p}\,Z_{2p+1}$ per
sub-block~$i$, which commutes with the rotated Hamiltonian by
construction (cross-block ERIs vanish in the standard composed
builder; within each sub-block the Gaunt selection rule
$m_1 + m_2 = m_3 + m_4$ is enforced at the ERI tensor level).
Together with $\{Z_\alpha, Z_\beta\}$ this yields $2 + n_{\rm sb}$
commuting stabilisers.

The reduction is library-universal because the symmetry is
basis-intrinsic.  Gaussian encodings break the $\mathbb{Z}_2$
structurally via the multi-center AO orthogonalisation that mixes
$m$-channels; only an angular-momentum eigenbasis that does not
cross-center-mix---such as GeoVac's natural composed build---preserves
the symmetry bit-exact, and only such a basis admits the further
decomposition into $n_{\rm sb}$ independent per-sub-block
reflections.  The $\mathbb{Z}_2$'s that survive JW encoding are a
direct structural consequence of the framework's geometric origin
(the Hopf bundle structure of $S^3$), not accidental basis-choice
symmetries.

The scaling pattern is notable.  Most tapering schemes give
diminishing returns at scale; for the GeoVac composed encoding
specifically, $\Delta Q$ scales linearly with sub-block count,
which scales linearly with chemical complexity (atoms, bonds,
lone pairs).  The largest beneficiaries are precisely the
production-relevant 100-qubit diatomics (CO, N$_2$, F$_2$ at
$\Delta Q = -12$), where qubit count is the binding constraint
for NISQ and early-FT simulation.  The per-sub-block construction
degrades gracefully under future extensions:\ if cross-block ERIs
become non-zero (e.g.\ \texttt{cross\_block\_h1=True}, or
balanced-coupled multipole), the per-stabiliser commutator audit
catches the offending $P_i$ and drops it from the stabiliser list,
recovering the global $n_{\rm sb}=1$ result rather than producing
an incorrect spectrum.  Per-molecule tapering panel documented in
\texttt{debug/sprint\_z2\_per\_subblock\_memo.md} (data file
\texttt{debug/data/sprint\_z2\_per\_subblock.json}; same-day
predecessor at \texttt{debug/sprint\_z2\_tapering\_memo.md}
covers the $n_{\rm sb}=1$ global construction).

\paragraph*{Triply-additive blocking with spin adaptation.}
The Hopf-$U(1)$ $\mathbb{Z}_2$ stabilisers $\{P_i\}_{i=1}^{n_{\rm sb}}$
commute pairwise with the total spin operator $S^2$ and with the
particle-number-and-$S_z$ stabilisers $\{Z_\alpha, Z_\beta\}$,
bit-exactly on LiH, BeH$_2$, and H$_2$O (every commutator coefficient
is exact zero across 21, 36, and 55 pairs respectively; see
\texttt{debug/sprint\_qpt\_hopf\_stacking\_memo.md}).  The Hopf tapering
therefore stacks additively with $S^2$-spin-adaptation techniques such
as the recent Quantum Paldus
Transform~\cite{burkat_fitzpatrick2025}, which block-diagonalises a JW
Hamiltonian by $u(d)\!\times\!SU(2)$ irrep at $O(d^3)$ Toffoli depth.
On LiH (4 electrons, 3 active Hopf sub-blocks), the joint $(N, S_z,
P_{\rm Hopf})$ sector is $5.7\times$ smaller than the
particle-number-only sector of standard Bravyi~2017
tapering~\cite{bravyi2017} ($4{,}791$ Slater determinants vs $27{,}405$),
with further reduction available from full QPT adaptation to $S^2$
eigenstates.  At spectral-triple level the structural reason is that
$\{P_i\}$ acts on spatial labels only (the $m_\ell \to -m_\ell$
reflection within each sub-block), while $S^2$ acts on spin labels
only; the two operators factor cleanly through the spatial--spin
tensor decomposition of the GeoVac one-body bundle.  The compatibility
extends to the spinor-coupled Tier 2 builder
(Sec.~\ref{sec:spinor_composed}) with $J^2$ replacing $S^2$:\ on LiH
relativistic at CODATA $\alpha$, $\|H_{\rm rel}(\alpha) - H_{\rm
rel}(0)\|_1 / \alpha^2 = 3.5417$ bit-identical at $\alpha/2, \alpha,
2\alpha$ (ratio max/min $= 1.0000$).  The natural spinor-side analog
of the non-relativistic $m_\ell \to -m_\ell$ Hopf reflection ---
the $m_j \to -m_j$ Z$_2$ stabiliser on the relativistic Dirac
basis --- does NOT commute with $H_{\rm rel}$ (residuals $1.7
\times 10^{-2}$ to $6.8 \times 10^{-2}$ on LiH/BeH/CaH$_{\rm rel}$
in both the original Dirac basis and after the (sym, antisym)
rotation; densification check bit-exact at $0$, ruling out
implementation bug). The structural reason is the load-bearing
finding of this section's tapering story:\ the non-relativistic
$m_\ell$ Z$_2$ is \emph{protected} by the integer-$\ell$ Gaunt
parity selection rule, whereas the half-integer-$j$ jj-coupled
chemistry has no analog (see Sec.~\ref{sec:gaunt_parity_protection}).
The relativistic-tapering frontier is therefore closed for Z$_2$
candidates at sprint scale; non-Z$_2$ machinery ($j^2$ sector
restriction analog of $S^2$; Kramers projection) remains a separate
research direction.  Across the $18$-molecule
hydride library at $n_{\max}=2$, the per-molecule QPT cost
($\sim M^3$ Toffoli, range $1{,}000$--$91{,}125$, total $561{,}000$)
is $4$--$6$ orders of magnitude below typical fault-tolerant QPE
budgets~\cite{lee2021,reiher2017} ($10^9$--$10^{11}$),
making spin-adaptation effectively free at GeoVac scale.

\subsubsection{$\ell$-parity Z$_2$:\ direct Gaunt-parity consequence
(strict improvement over Hopf-only)}
\label{sec:ell_parity_tapering}

The Gaunt selection rule for the chemistry ERI requires
$(\ell_a + \ell_c + k)$ even AND $(\ell_b + \ell_d + k)$ even on
every nonzero $(ab|cd)$, which forces
$\ell_a + \ell_b + \ell_c + \ell_d$ to be even on every nonzero
ERI element.  Equivalently, the operator
$P_\ell^{(i)} = (-1)^{N_{\ell\text{-odd}}^{(i)}}$, counting electrons
in $\ell$-odd orbitals of sub-block $i$, commutes with the
two-body Coulomb sector exactly.  Within-block $h_1$ in the
composed builder is $\ell$-diagonal so $[P_\ell^{(i)}, h_1] = 0$
trivially.

This is an additional $\mathbb{Z}_2$ stabiliser per sub-block,
GF($2$)-independent of the Hopf $m_\ell$ stabiliser of
Sec.~\ref{sec:hopf_tapering} (an $\ell=1$, $m_\ell=0$ orbital
carries Hopf parity $+1$ but $\ell$-parity $-1$).  No basis
rotation required:\ orbitals already carry definite $\ell$, so
$P_\ell^{(i)}$ is a JW $Z$-string directly.  Production module
\texttt{geovac/extended\_tapering.py} ships this as
\texttt{use\_ell\_parity=True} (default).

\paragraph*{Verification panel} (composed builder, $n_{\max} = 2$):
\begin{center}
\begin{tabular}{lrrrrr}
System & $Q_{\rm naive}$ & $Q_{\rm Hopf}$ & $Q_{\rm Hopf}+\ell$ &
$\Delta Q_\ell$ & Pauli ratio \\
\hline
LiH    & 30 & 25 & \textbf{22} & $+3$ & $0.97$ \\
HF     & 60 & 52 & \textbf{46} & $+6$ & $0.97$ \\
BeH$_2$  & 50 & 43 & \textbf{38} & $+5$ & $0.97$ \\
H$_2$O   & 70 & 61 & \textbf{54} & $+7$ & $0.97$ \\
NH$_3$   & 80 & 70 & \textbf{62} & $+8$ & $0.97$ \\
CH$_4$   & 90 & 79 & \textbf{70} & $+9$ & $0.97$ \\
\end{tabular}
\end{center}
$\Delta Q_\ell = n_{\rm sb}$ (one stabiliser per sub-block with
$\ell$-odd content) is matched bit-exactly across the panel.  The
Pauli-count ratio is uniformly $\sim 0.97$ (a $\sim 3\%$ reduction
beyond Hopf-only), making this a \emph{strict improvement} over
Hopf-only on both axes.  Ground-state spectrum preservation is
verified on H$_2$: the minimum over tapered sector choices matches
the naive JW ground energy to $<10^{-10}$
(test \texttt{tests/test\_extended\_tapering.py},
\texttt{test\_extended\_hopf\_ell\_spectrum\_preserved\_h2}).

\subsubsection{Hidden per-sub-block particle-conservation Z$_2$:\
meta-investigation via symmetry-adapted basis}
\label{sec:hidden_particle_z2}

The combined $\mathbb{Z}_2^3$ irrep decomposition under
$(p_{\rm Hopf}, p_\ell, p_{\rm swap})$ exposes additional hidden
$\mathbb{Z}_2$ symmetries that the piecewise stabiliser search of
Secs.~\ref{sec:hopf_tapering}--\ref{sec:ell_parity_tapering} would
not surface.  Concretely, in the composed builder the ERI is
block-diagonal across sub-blocks (no cross-block ERI by default,
\texttt{composed\_qubit.py:727--745}), so the per-sub-block electron
count $N_b$ is conserved \emph{individually}, and
\begin{equation}
  P_N^{(b)} = (-1)^{N_b}
\label{eq:per_block_particle_z2}
\end{equation}
is a $\mathbb{Z}_2$ stabiliser commuting with $H$, independent in
GF($2$) of the Hopf $m_\ell$ stabiliser (which only $Z$-strings the
$\ell$-odd antisymmetric orbitals) and of the $\ell$-parity
stabiliser (which only $Z$-strings the $\ell$-odd orbitals).
$P_N^{(b)}$ $Z$-strings \emph{all} orbitals in block $b$, a
strictly larger support.  Production module
\texttt{geovac/symmetry\_adapted\_basis.py} ships this with the
\texttt{extended\_plus\_hidden\_tapered\_from\_spec} entry point;
wired into the production API at
\texttt{ecosystem\_export.hamiltonian(tapered='full')}.

Across the same panel (extending Sec.~\ref{sec:ell_parity_tapering}):
\begin{center}
\begin{tabular}{lrrr}
System & $Q$ Hopf$+\ell$ & $Q$ full & hidden $\Delta Q$ \\
\hline
LiH    & 22 & \textbf{20} & $+2$ \\
HF     & 46 & \textbf{41} & $+5$ \\
BeH$_2$  & 38 & \textbf{37} & $+1$ \\
H$_2$O   & 54 & \textbf{53} & $+1$ \\
NH$_3$   & 62 & \textbf{55} & $+7$ \\
CH$_4$   & 70 & \textbf{62} & $+8$ \\
\hline
Total  & 292 & \textbf{268} & $+24$ \\
\end{tabular}
\end{center}
The audit gate verifies GF($2$)-independence of each candidate
stabiliser; $24$ of $132$ panel candidates pass and contribute.
Ground-sector agreement between the extended and extended$+$hidden
tapered forms is verified to $<10^{-10}$~Ha on LiH max\_n=$1$ via full
diagonalisation, alongside block-diagonality and stabiliser-audit
checks (test \texttt{tests/test\_symmetry\_adapted\_basis.py}); the
naive-vs-tapered spectrum comparison is carried by the H$_2$
$\ell$-parity test above.  Combined with
Hopf $+ \ell$-parity, the production
\texttt{tapered='full'} pipeline saves $30 \to 20$ on LiH (33\%
reduction) and $90 \to 62$ on CH$_4$ (31\%) over naive JW --- a
multiplicative win across the chemistry-encoding-relevant
$\sim Q=30$--$100$ qubit regime.

\subsubsection{Structural reason:\ Gaunt parity protects integer-$\ell$
parity Z$_2$'s but not their half-integer-$j$ analogs}
\label{sec:gaunt_parity_protection}

The Hopf $m_\ell$ Z$_2$ (Sec.~\ref{sec:hopf_tapering}) and the
$\ell$-parity Z$_2$ (Sec.~\ref{sec:ell_parity_tapering}) both
commute with $H$ for the \emph{same} underlying structural reason:\
under all-orbital sign-flip of the parity quantum number in a
4-orbital ERI element, the angular factor accumulates a sign
$(-1)^{\sum_i p_i}$ from the four $3j$-symbol parity factors.  For
integer $\ell$ with sign flip $m_\ell \to -m_\ell$ (Hopf) the
multipole Gaunt rule $(\ell_a + \ell_c + k)$ even AND
$(\ell_b + \ell_d + k)$ even forces $\ell_a + \ell_b + \ell_c + \ell_d$
even.  For integer $\ell$ with the $\ell$-parity sign flip directly,
the same rule forces the analog sum even.  Both Z$_2$'s are
\emph{protected} by the Gaunt parity selection rule on the
multipole-expanded ERI:\ the four-orbital sign factor is
\emph{forced} to be $+1$, ensuring the ERI tensor is invariant
under the reflection.

The half-integer-$j$ jj-coupled analog has no such protection.
Under all-orbital $m_j \to -m_j$ flip the four 3j-symbol sign
factors give $(-1)^{j_a + j_b + j_c + j_d}$, and the jj-coupled
selection rules do not force this sum to be even:\ a quadruple
$(j_a, j_b, j_c, j_d) = (3/2, 1/2, 1/2, 1/2)$ has sum $3$ (odd),
factor $-1$, and such ERI elements are present at non-zero
amplitude in $H_{\rm rel}$.  This is the load-bearing structural
reason the relativistic $m_j$ Z$_2$ candidates fail (both
direct-basis and rotated-basis; see the $1.7$--$6.8 \times 10^{-2}$
residuals reported above).

The general principle is:\ for any parity-style $\mathbb{Z}_2$
tapering candidate in a chemistry construction, the candidate
commutes with the two-body Coulomb sector if and only if the four
angular-momentum-parity factors are \emph{forced} to be $+1$ by the
ERI's angular selection rules.  Integer-$\ell$ Gaunt parity
provides this protection cleanly (yielding the Hopf + $\ell$-parity
Z$_2$'s); half-integer-$j$ jj-coupling does not.  This is the
underlying reason GeoVac's chemistry-side computational story
(O($Q^{2.5}$) Pauli scaling, $\sim 30\%$ qubit reduction at
production via four-axis tapering, Pauli-count reduction) is the
non-relativistic composed builder rather than its relativistic
Tier 2 sibling.

\subsubsection{MPO bond rank: structural inheritance from the
multipole expansion}
\label{sec:mpo_bond_rank}

The operator Schmidt rank $\chi_k$ at qubit cut $k$ of a chemistry
Hamiltonian under Jordan--Wigner ordering equals the MPO bond
dimension after maximally compact decomposition.  For chemistry
Hamiltonians this can be computed directly from the Pauli expansion
without running DMRG, by SVD of the coefficient matrix
$M[P_{\rm L}, P_{\rm R}]$ that decomposes each Pauli string into left
and right parts at the cut.

\paragraph*{Bond rank decomposition of the GeoVac composed
Hamiltonian (Thm.~3.2.A.unified).}
\label{thm:mpo_bond_rank}
Let $H = h_1 + V_{ee}$ be a GeoVac composed Hamiltonian on $N$
spin-orbitals.  At every Jordan--Wigner cut $k \in \{1, \ldots, N-1\}$,
the operator Schmidt rank $\chi_k$ decomposes into five components.

\emph{(A) One-body exact closed form.}
\begin{equation}
  \chi_k^{h_1}  =  2\,\mathrm{rank}\big(h_{p\le k,\,q > k}\big)
  + \mathbf{1}_{LL} + \mathbf{1}_{RR},
\label{eq:chi_h1_closed}
\end{equation}
where $h_{p\le k,\,q>k}$ is the $k \times (N-k)$ cross-cut sub-matrix
of $h_1$ in spin-orbital indexing and $\mathbf{1}_{LL}, \mathbf{1}_{RR}$
indicate the existence of nonzero pure-left and pure-right one-body
content.  The factor of $2$ is the JW $XX + YY$ structure of the
cross-cut hopping.

\emph{(B) Two-piece subadditivity.}
\begin{equation}
  \chi_k^H  \le  \chi_k^{h_1} + \chi_k^{V_{ee}},
\label{eq:chi_subadd}
\end{equation}
with slack absorbed by additive baseline mode sharing between $h_1$
and $V_{ee}$.

\emph{(C) Per-multipole subadditivity.}
\begin{equation}
  \chi_k^{V_{ee}}  \le  \sum_L \chi_k^{V_{ee}^L},
\label{eq:chi_perL_subadd}
\end{equation}
where $V_{ee}^L$ is the $L$-multipole component of $V_{ee}$ under the
Slater--Gaunt decomposition (Sec.~\ref{sec:methods}).  Slack absorbs
cross-$L$ operator overlap when a fermionic pair $a^\dagger_p a_q$
carries weight in multiple Gaunt-allowed $L$.

\emph{(D) F3 inheritance.}  The per-$L$ rank $\chi_k^{V_{ee}^L}$
inherits the graded singular-value spectrum of the radial-integral
matrix $R^L(\rho_\alpha; \rho_\beta)$ that controls the
double-factorization rank reported in Sec.~\ref{sec:intro} (the same
hydrogenic-basis Laguerre-recurrence ill-conditioning).

\emph{(E) Boundary saturation.}  At every sub-block boundary the
composed builder has zero cross-block ERI content (F4 of the
DF--multipole equivalence), all cross-cut content vanishes, and
$\chi_k^H = 2$ (additive baseline of $H_{\rm left\ block} \otimes I
+ I \otimes H_{\rm right\ block}$).

\paragraph*{Universal interior profile.}
For the standard $\{1s, 2s, 2p_{-1,0,+1}\}$ valence at $n_{\max}=2$,
the interior bond rank within a sub-block follows the universal core
sequence $\chi_k^H = \{16, 16, 9, 9, 9, 6, 3\}$ at local cuts $k =
2, \ldots, 8$, depending only on the orbital basis; the edge cuts are
position-dependent (first cut: $4$ for the leading sub-block, $5$ for
sub-blocks with a completed block to their left --- one extra additive
baseline mode; the final sub-block's last cut saturates to the
chain-end value $2$).  The leading-sub-block sequence at cuts
$1,\ldots,10$ is $\{4, 16, 16, 9, 9, 9, 6, 3, 3, 2\}$.  Verified
bit-exactly on LiH at $29/29$ cuts (recomputed from the live builder
and pinned by \texttt{tests/test\_paper14\_mpo\_rank.py}); the
$\chi_k = 2$ boundary saturation is verified across LiH
(3 sub-blocks), BeH$_2$ (5 sub-blocks), and H$_2$O (7 sub-blocks) at
$12/12$ sub-block boundaries, giving $8\times$
interior/boundary ratio.  The dominant rank carrier transitions with
the cut position:\ $L=0$ governs cuts $1$--$3$ (1s--2s cross-cut
active), $L=1$ governs cuts $4$--$6$
(1s,\,2s\,$\leftrightarrow$\,2p active), $L=2$ governs cuts $7$--$9$
(sequential 2p closure).  The profile is determined by the orbital
structure and the multipole hierarchy ($L_{\max} = 2 \ell_{\max}$),
with no fit parameters.

\paragraph*{Bit-exact verification panel.}
On LiH, statement~(A) holds at $29/29$ cuts
(eq.~\eqref{eq:chi_h1_closed} exact); (B) and (C) hold at $29/29$
with average slack $2.3$ and $5.9$ respectively, both from shared
additive baseline modes between $h_1, V_{ee}^L$.  Driver and data
files (permanent):\ \texttt{tests/mpo\_rank\_support/sprint\_s2\_v2\_unified\_panel.py} +
\texttt{tests/mpo\_rank\_support/data/sprint\_s2\_v2\_unified\_panel.json};
the full LiH panel is recomputed from the live builder and the shipped
panel pinned by \texttt{tests/test\_paper14\_mpo\_rank.py}
(sprint chronicle:\ CHANGELOG v3.54.0).  Negative control extended to
the balanced-coupled library:\ across LiH, BeH$_2$, H$_2$O at $12/12$
sub-block boundaries the chi$=2$ floor is broken into a clean
two-value quantization $\chi_k^{H,\,{\rm balanced}} \in \{9, 16\}$ that
decomposes as $\chi_k^{H,\,{\rm balanced}} - 2 = N_{\rm cross}(k)
\times 7$, where $N_{\rm cross}(k) \in \{1, 2\}$ counts cross-block
multipole couplings straddling the boundary and $7$ equals the
double-factorisation rank of the $\{1s, 2s, 2p\}$ valence basis
(Sec.~\ref{sec:intro}).  The boundary lift therefore directly
inherits the F3 graded structure: the same Laguerre--Gaunt ring that
fixes DF rank fixes the per-boundary lift, bit-identically across
three molecules and twelve boundaries (Sprint S1 / S2-v2 negative
controls; panel data
\texttt{tests/mpo\_rank\_support/data/sprint\_s2\_v2\_balanced\_library\_panel.json},
pinned by \texttt{tests/test\_paper14\_mpo\_rank.py}; chronicle
CHANGELOG v3.54.0).

This bond-rank inheritance is a chemistry-side manifestation of the
spectral-truncation GH-convergence structure underlying the GeoVac
spectral triple:\ the $L$-multipole density operators are the natural
carriers of the operator system at the cut, and the MPO bond rank
inherits the graded rank of the radial-integral matrix $R^L$ from the
same Laguerre--Gaunt ring that supplies double-factorization
compression (Sec.~\ref{sec:intro}).  Properties~(A)--(E) above provide
a quantitative DMRG-side realisation of the qualitative ``$V_{ee} \in
\mathcal{A}\otimes \mathcal{A}$'' statement of Sec.~\ref{sec:intro}.

\emph{Pauli grouping} (QWC, general commutativity, unitary
partitioning~\cite{verteletskyi2020,yen2020}) reduces measurement
overhead by identifying simultaneously measurable Pauli strings.
These algorithms take a set of Pauli strings as input and group them;
a smaller input set produces fewer groups.  GeoVac provides the
smaller input set---$Q^{3.8}$ instead of $Q^{4+}$ terms---so
grouping starts from a more favorable position.

\emph{Low-rank factorization}~\cite{motta2021,lee2021} decomposes the
ERI tensor to reduce the number of distinct rotation circuits.  The
rank of the GeoVac ERI tensor is bounded by the number of allowed
Gaunt channels, which grows as $O(l_{\max}^2)$---already much lower
than the $O(M^2)$ rank of a generic dense tensor.  Applying low-rank
factorization to a natively sparse tensor would start from a lower
rank, potentially achieving further compression.

In short, GeoVac's structural sparsity provides a better starting
point for every downstream optimization in the standard quantum
chemistry compilation pipeline.  The $Q^{3.17}$ composed scaling
(Sec.~\ref{sec:composed}) is a \emph{floor} that existing reduction
techniques can only improve upon.

\subsection{Fault-tolerant resource implications}
\label{sec:ft_resources}

For fault-tolerant quantum simulation, the total cost is controlled
by the product of gate count per Trotter step and the number of
steps~$r$.  We emphasize two structural advantages of the GeoVac
encoding in this regime.

First, the sub-quadratic 1-norm scaling ($\lambda \sim Q^{1.77}$)
means that the GeoVac advantage \emph{grows faster} for fault-tolerant
algorithms than for VQE.  The Trotter step count
$r \propto \lambda / \sqrt{\epsilon}$ depends on 1-norm, not term
count; the composed 1-norms (Table~\ref{tab:composed_onenorm}) show
that LiH at $Q = 84$ requires $r \approx 4{,}500$ steps at chemical
accuracy ($\epsilon = 10^{-3}$), while BeH$_2$ at $Q = 140$ requires
$r \approx 16{,}500$---modest counts for early fault-tolerant devices.

Second, the commutator-based bound (Table~\ref{tab:comm_bound})
provides a further $4$--$7\times$ step reduction that also grows
with system size ($r_{\mathrm{comm}} \sim Q^{1.47}$ vs.\
$r_{\lambda^2} \sim Q^{1.77}$).  This reflects the physical
mechanism: selection-rule sparsity ensures that an increasing fraction
of Pauli term pairs commute, reducing the effective error per step.

For qubitization~\cite{childs2021,lee2021}, the query complexity
scales as $O(\lambda t / \epsilon)$ with no quadratic dependence.
The $Q^{1.77}$ 1-norm scaling translates directly into a
$Q^{0.3\text{--}0.5}$ query-count advantage over encodings with
$\lambda \sim Q^{2+}$.  At $Q = 110$, this corresponds to a
$\sim 3\times$ reduction in queries; at $Q = 1000$ (extrapolated),
the advantage would grow to $\sim 10\times$.

These advantages compose with the Pauli term count reduction: fewer
terms means fewer gates per Trotter step, and lower 1-norm means
fewer steps.  The product $N_{\mathrm{Pauli}} \times r$ determines
the total gate budget, and both factors are simultaneously reduced
by the GeoVac encoding.

\subsection{Context: fault-tolerant Gaussian compression methods}
\label{sec:ft_gaussian}

The fault-tolerant QPE literature employs block-encoding methods---double
factorization (DF)~\cite{motta2021}, tensor hypercontraction
(THC)~\cite{lee2021}, and symmetry-compressed double factorization
(SCDF)~\cite{rocca2024}---that do not produce Pauli decompositions.
These methods target the block-encoding 1-norm~$\lambda$ and Toffoli
gate count rather than Pauli term count, and are implemented as
sequences of rotations within a linear combination of unitaries (LCU)
circuit.  Published benchmarks focus on large active spaces: FeMoCo at
108--152~qubits~\cite{lee2021}, cytochrome~P450 at
96--116~qubits~\cite{caesura2025}, and similar production-scale systems,
with further reductions reported through improved tensor factorization
and active volume compilation~\cite{caesura2025}.
To our knowledge, no published DF, THC, or SCDF $\lambda$~values exist
for molecules at the 10--70~qubit operating scale of the present work
(LiH at $Q = 30$, H$_2$O at $Q = 70$), because these systems are too
small for the overhead of block-encoding to be beneficial.
The Pauli term advantage reported here ($\sim\!50\times$--$320\times$,
Sec.~\ref{sec:composed_gaussian}; the retired pair-diagonal rule gave
51$\times$--1,712$\times$) is therefore against raw
Jordan--Wigner encoding, the representation used by VQE and other
near-term variational algorithms.  The 1-norm comparison
(Sec.~\ref{sec:onenorm}) and its scaling ($\lambda \sim Q^{1.77}$)
are the metrics relevant for both VQE shot budgets and fault-tolerant
qubitization query complexity.

\paragraph{Computed 1-norm comparison.}
To quantify the 1-norm comparison directly, we compute Gaussian raw
Jordan--Wigner baselines from published molecular integrals.
For LiH~STO-3G (6~spatial orbitals, $Q = 12$, FCI $E = -7.881$~Ha),
the raw JW Hamiltonian has \gvq{lih_sto3g_raw}{907}~Pauli terms, $\lambda = 34.3$~Ha, and
273~QWC groups.  The composed GeoVac LiH ($Q = 30$, electronic-only
with PK partitioned classically) has \gvq{lih_composed_pauli}{838}~Pauli terms,
$\lambda = 34.0$~Ha, and 64~QWC groups (measured 2026-08-30 under the
exact rule; the retired pair-diagonal rule gave 334, $32.6$~Ha and 21).
At matched raw-JW
convention---not matched qubit count (GeoVac uses $2.5\times$ more
qubits, $Q=30$ vs.\ $12$) and not matched accuracy (the composed LiH
carries $\sim$5\% $R_\mathrm{eq}$ geometry error and does not bind
under qubit FCI)---GeoVac is at \emph{near parity} on Pauli terms
($1.08\times$ fewer, 838 vs.\ 907) with $4.3\times$ fewer QWC groups
and essentially identical 1-norm ($0.99\times$).  Under the retired rule
this read ``$2.7\times$ fewer Pauli terms and $13\times$ fewer QWC
groups''---a decisive-looking win that the correction removes.
The 1-norm parity reflects that $\lambda$ is dominated by one-body
terms (core energies, PK), which scale with nuclear charge rather than
orbital count.  The Pauli and QWC advantage are the VQE-relevant
metrics.  For the validated He equal-qubit comparison at $Q = 28$,
the GeoVac lattice has $7.1\times$ lower 1-norm than Gaussian
cc-pVTZ ($\lambda = 74.21$ vs.\ $530.47$~Ha), demonstrating that the
1-norm advantage grows with basis size for single-center systems.

\paragraph{BLISS-THC at the production-scale ceiling.}
For context on the FTQC operating scale, Caesura
\textit{et al.}~\cite{caesura2025} report the current best-of-class
BLISS-THC factorization for cytochrome P450 in a (63e, 58o) active
space ($Q = 116$ spin-orbitals): 999~memory qubits,
$\lambda_{\mathrm{BLISS\text{-}THC}} \approx 133$~Ha,
and $1.71 \times 10^9$ Toffoli gates after circuit improvements;
the prior-art THC of Goings \textit{et al.}~\cite{goings2022}
reported $\lambda_{\mathrm{THC}} = 388.9$~Ha and
$7.79 \times 10^9$ Toffolis at the same active space.  Caesura's
$\sim$233$\times$ runtime speedup over baseline THC combines a
factor of $\sim$25 from Active Volume photonic compilation, a
factor of $\sim$8 from BLISS within THC, and $\sim$1.1 from
block-encoding circuit improvements; the $\lambda$ improvement
alone is roughly $3\times$, bringing the BLISS-THC 1-norm to within
a factor of 2 of the symmetry-shifted lower bound 69.3~Ha for
P450~\cite{caesura2025}.  These resource estimates are reported
for a single P450 active space, the smallest molecule in the
Caesura benchmark; the GeoVac multi-center library
reaches CO and N$_2$ at $Q = 100$ with
$2{,}790$~Pauli terms under the composed encoding (Sec.~\ref{sec:multi_center}),
but Pauli term count and 1-norm/Toffoli count are not directly
comparable metrics --- Pauli count drives VQE measurement overhead
and Trotterization step depth, while the BLISS-THC pipeline
optimizes the qubitization 1-norm and Toffoli budget through a
non-Pauli (block-encoded LCU) circuit family.  The two operating
points target different algorithmic regimes; the appropriate
head-to-head comparison would require running BLISS-THC on a
GeoVac-scale molecule (LiH, H$_2$O), which is below the published
benchmark range, or running GeoVac at P450 scale, which is above
the current library ceiling.

\subsection{Limitations}
\label{sec:limitations}

Several caveats apply to the present results.

\emph{System size.}  The scaling exponents are extracted from four data
points at $Q = 10$--$110$ qubits.  While the ERI density decay
($\approx M^{-0.49}$ under the exact rule) has an analytical foundation
in the Gaunt integral selection rules, the
precise exponents ($\alpha_{\mathrm{terms}} = 3.77$,
$\alpha_{\lambda} = 1.77$) may shift slightly as larger systems (Li, Be,
molecular H$_2$ at higher $n_{\max}$) are included.  All $R^2$~values
exceed 0.997, indicating good power-law behavior over the range tested.

\emph{Gaussian baseline.}  The genuine Gaussian integrals
(Sec.~\ref{sec:genuine_gaussian}) now cover He at three basis sizes
(STO-3G, cc-pVDZ, cc-pVTZ), providing validated equal-qubit
comparisons at $Q = 2, 10, 28$.  At the two largest sizes the
1-norm advantage holds---$3.8\times$ and $7.1\times$ lower---but the
Pauli-term comparison does \emph{not} run one way: at $Q = 10$ GeoVac has
$1.8\times$ \emph{more} terms, recovering to $1.5\times$ fewer at $Q = 28$.
(Corrected 2026-08-30; the retired pair-diagonal rule gave $1.3\times$ and
$8.1\times$ fewer terms and $7.1\times$ lower 1-norm at $Q=28$.)
The 6-31G entry in Table~\ref{tab:pauli} still uses synthetic
integrals with representative sparsity.  The molecular Gaussian
baselines use the published Pauli counts of
Trenev~\textit{et al.}~\cite{trenev2025} (Table~5, Appendix~B).  The structural
conclusion---Gaussian ERIs are near-dense while lattice ERIs
are selection-rule-sparse---is robust and well
established~\cite{helgaker2000}.

\emph{1-norm bound.}  The Trotter step counts in Table~\ref{tab:onenorm}
use the standard $\lambda^2$ upper bound (Childs et al.~\cite{childs2021},
Eq.~48).  As shown in Table~\ref{tab:comm_bound}, commutator-based bounds
that exploit the detailed Hamiltonian structure yield $7\times$ fewer Trotter
steps at $Q = 60$, with scaling exponent $\alpha = 1.47$ versus~$1.77$
(the $1.69$ previously quoted here was the retired $\lambda$ exponent).
The commutator sweep covers $n_{\max} = 2$--$4$ (three data points);
the $n_{\max} = 5$ system ($N = \gvq{he_n5_pauli}{2{,}434{,}441}$ terms under the exact
rule; retired: $227{,}338$) requires
an $N \times N$ anticommutativity matrix of order $10^{12}$ entries
($\sim\!6$~TB as bytes) and was not computed.  The $\sim\!52$~GB figure
previously quoted here was $227{,}338^2$, i.e.\ computed from the retired
term count.  The scaling exponent may shift slightly with additional data
points.

\emph{Orbital ordering.}  Jordan--Wigner Pauli weights depend on the
spin-orbital ordering.  We use the natural lattice ordering
$(n, l, m, \sigma)$ sorted lexicographically.  Optimized orderings (e.g.,
interleaving by angular momentum) could further reduce the average Pauli
weight without changing the term count.

\emph{Beyond two-electron integrals.}  For systems with strong
relativistic effects or explicit correlation factors, additional integral
classes arise that are not captured by the present analysis.  The selection
rule argument applies specifically to the nonrelativistic Coulomb ERI
tensor.


%% ========================================================================
\section{Conclusion}
\label{sec:conclusion}
%% ========================================================================

We have demonstrated that qubit Hamiltonians constructed from spectral
graph theory lattices produce structurally sparse Pauli decompositions
across three operationally relevant cost metrics:
\begin{itemize}
\item \textbf{Pauli term count:} $O(Q^{3.8})$, a $Q^{0.1\text{--}0.5}$-fold
  improvement over the $O(Q^{3.9\text{--}4.3})$ scaling of published
  Gaussian-basis encodings~\cite{trenev2025} (corrected 2026-08-30; the
  retired exponent $3.15$ gave $Q^{1.0\text{--}1.5}$).  The load-bearing
  advantage is the exact $Q$-linearity across molecules, not this
  within-molecule margin.
\item \textbf{QWC measurement groups:} $O(Q^{4.01})$, slightly steeper
  than the term count, so the shot budget for variational algorithms is
  \emph{not} a source of advantage under the exact rule (re-measured
  2026-08-30; the retired pair-diagonal fit gave $O(Q^{3.36})$, below the
  term exponent).
\item \textbf{Pauli 1-norm:} $O(Q^{1.77})$ (re-measured 2026-08-30;
  retired: $O(Q^{1.69})$), the key result for
  fault-tolerant simulation.  Because the 1-norm controls both
  Trotter step counts and qubitization query complexity, this
  sub-quadratic scaling implies that the GeoVac lattice basis
  produces Hamiltonians that are not just sparse in structure but
  \emph{concentrated in weight}.
\end{itemize}
The advantage is structural: angular momentum selection rules in the
Gaunt integrals enforce a decaying ERI density ($\approx M^{-0.49}$
under the exact rule), and the surviving
integrals are dominated by a small number of large diagonal
contributions.  Equal-qubit comparisons against validated Gaussian
cc-pVDZ and cc-pVTZ integrals confirm the advantage at the 1-norm
level: $7.1\times$ lower $\lambda$ at $Q = 28$, translating directly
into fewer Trotter steps.  This sparsity is a \emph{basis property},
present before any qubit tapering, Pauli grouping, or tensor
factorization---and compatible with all of these downstream
optimizations (Sec.~\ref{sec:composability}).

For quantum chemistry on quantum hardware, where circuit depth and
measurement budget are the primary constraints, the choice of
second-quantized basis can be as consequential as the choice of quantum
algorithm.  The GeoVac lattice Hamiltonian, with its $O(N)$ classical
construction cost~\cite{loutey_paper7}, $\sim Q^{3.8}$ Pauli encoding,
and $O(Q^{1.77})$ simulation-cost scaling (re-measured 2026-08-30;
retired: $O(Q^{1.69})$), provides a
starting point for both NISQ variational methods and fault-tolerant
Hamiltonian simulation.  For composed multi-center systems (Sec.~\ref{sec:composed}), the
pseudopotential approximation yields block-diagonal ERIs whose Pauli
count factorizes exactly as $N_{\mathrm{Pauli}} = c(n_{\max})\,Q$, with a
single within-molecule exponent of $3.17$ for every system tested
(corrected 2026-08-29; the retired rule gave $2.50$--$2.52$).
Against Gaussian baselines~\cite{trenev2025}: at $Q = 196$, composed
H$_2$O requires $\sim 320\times$ fewer Pauli terms than the Gaussian
power-law interpolation.  The decisive comparison is the system-size
axis, where the exact linearity in $Q$ stands against the Gaussian
$Q^{3.9\text{--}4.3}$---a gap that widens with system size, which is what
makes this advantage structural and persistent.  The composition
architecture~\cite{loutey_paper17}---which decomposes each molecule
into a fiber bundle of single-center orbital blocks---provides a
concrete path to polyatomic systems: each additional atomic center
adds an independent ERI block without creating cross-block terms,
provided the pseudopotential approximation remains valid.

Software implementing the encoder
(\texttt{geovac.qubit\_encoding}), measurement grouping
(\texttt{geovac.measurement\_grouping}), and Trotter analysis
(\texttt{geovac.trotter\_bounds}), together with all benchmarks
reported here, are available as part of the open-source GeoVac package.

\paragraph{Note added v2.0.24.}
The full $N$-electron quantum encoding comparison (Track~AS,
v2.0.23) has confirmed that the composed architecture is categorically
sparser than direct grid encoding of the exact 4-electron angular
Hamiltonian: at $Q = 10$, the full Jordan--Wigner encoding produces
3{,}\gvq{he_n2_pauli}{288}~Pauli terms with a 1-norm of 739~Ha, compared to
838~terms and 37~Ha for the composed encoding at $Q = 30$.
\emph{Vintage note (2026-08-30):} the composed term count is the
exact-rule value (retired: 334), but the 1-norms on both sides of this
comparison were measured under the retired rule and have not been
re-derived---the full Jordan--Wigner side is not reproducible from the
shipping library.  The qualitative conclusion (the composed architecture
is categorically sparser than direct grid encoding, $838 \ll 3{,}288$ at
three times the qubit count) is unaffected.
The root cause is structural: the full encoding maps finite-difference
grid points in the 11-dimensional angular space (dense FD coupling),
while the composed encoding maps orbital occupations with
block-diagonal ERIs from Gaunt selection rules.  Cross-pair $V_{ee}$
costs more in Pauli terms than the PK/$Z_{\rm eff}$/exchange overhead
it replaces.  This validates the composed architecture as the quantum
computing approach for multi-electron molecules.

\paragraph{Note added v2.0.44.}
Paper~20 presents the GeoVac Hamiltonian library as a practical tool
for the quantum computing community, including FCI accuracy benchmarks
across the LiH potential energy surface, matched-qubit (not
matched-accuracy) resource comparisons against Gaussian baselines,
and ecosystem export to
Qiskit, OpenFermion, and PennyLane.  The balanced coupled architecture
(Paper~19) achieves 0.20\% energy error at FCI, demonstrating that
the composed partition Hamiltonian is accurate when solved exactly.

\paragraph{Note added v2.4.0.}
The molecule library now spans 37 systems (35 composed molecules plus
the He and H$_2$ reference atoms), including
the full first transition series of hydrides (ScH--ZnH; $Z = 21$--30,
[Ar] frozen cores).  All ten use isolated $d$-orbital
valence blocks ($l_{\min} = l_{\max} = 2$) with ERI density of 4.0\%,
half that of $s$+$p$ blocks (8.9\%).  The resulting Pauli/$Q$
ratio of~9.23 falls below the universal coefficient of~11.10 \emph{under
the retired pair-diagonal ERI rule} (Sec.~\ref{sec:eri_rule})---an artifact
of dropping more $m$-swap integrals at higher angular momentum, not a
genuine saving (under the exact global-$M_L$ rule the $d$-block is denser).
Transition metals
are reclassified from ``out of scope'' to ``scoping demonstrated.''

\paragraph{Note added v2.5.0.}
A systematic investigation of nested (non-composed) hyperspherical
encodings (Track~DF, 6~sprints) established two results:
(i)~H-set angular momentum recoupling introduces additional ERI zeros
beyond bare Gaunt selection rules, reducing ERI density by 18--26\%
relative to uncoupled bases (Sec.~\ref{sec:recoupling_sparsity});
(ii)~the composed block-diagonal factorization is structurally necessary
for heteronuclear molecular PES---single-center and mixed-exponent
nested approaches produce correct binding energies but fail to reproduce
equilibrium geometries (33.7\% and 17.1\% $R_{\rm eq}$ error for LiH).
The 6$j$ recoupling sparsity is basis-agnostic and applies to any
encoding that uses angular momentum pair coupling.

%% ========================================================================
\begin{acknowledgments}
This work was supported by computational resources provided by the GeoVac
development infrastructure.
\end{acknowledgments}
%% ========================================================================

\begin{thebibliography}{30}

%% --- Quantum computing / quantum chemistry reviews ---

\bibitem{mcardle2020}
S.~McArdle, S.~Endo, A.~Aspuru-Guzik, S.~C.~Benjamin, and X.~Yuan,
``Quantum computational chemistry,''
Rev.\ Mod.\ Phys.\ \textbf{92}, 015003 (2020).

\bibitem{cao2019}
Y.~Cao, J.~Romero, J.~P.~Olson, \textit{et al.},
``Quantum chemistry in the age of quantum computing,''
Chem.\ Rev.\ \textbf{119}, 10856--10915 (2019).

%% --- Qubit reduction and Pauli grouping ---

\bibitem{bravyi2017}
S.~Bravyi, J.~M.~Gambetta, A.~Mezzacapo, and K.~Temme,
``Tapering off qubits to simulate fermionic Hamiltonians,''
arXiv:1701.08213 (2017).

\bibitem{verteletskyi2020}
V.~Verteletskyi, T.-C.~Yen, and A.~F.~Izmaylov,
``Measurement optimization in the variational quantum eigensolver using a
minimum clique cover,''
J.\ Chem.\ Phys.\ \textbf{152}, 124114 (2020).

\bibitem{yen2020}
T.-C.~Yen, V.~Verteletskyi, and A.~F.~Izmaylov,
``Measuring all compatible operators in one series of single-qubit
measurements using unitary transformations,''
J.\ Chem.\ Theory Comput.\ \textbf{16}, 2400--2409 (2020).

%% --- Tensor factorization ---

\bibitem{motta2021}
M.~Motta, E.~Ye, J.~R.~McClean, \textit{et al.},
``Low rank representations for quantum simulation of electronic structure,''
npj Quantum Inf.\ \textbf{7}, 83 (2021).

\bibitem{lee2021}
J.~Lee, D.~W.~Berry, C.~Gidney, \textit{et al.},
``Even more efficient quantum computations of chemistry through tensor
hypercontraction,''
PRX Quantum \textbf{2}, 030305 (2021).

\bibitem{goings2022}
J.~J.~Goings, A.~White, J.~Lee, C.~S.~Tautermann, M.~Degroote,
C.~Gidney, T.~Shiozaki, R.~Babbush, and N.~C.~Rubin,
``Reliably assessing the electronic structure of cytochrome P450 on
today's classical computers and tomorrow's quantum computers,''
Proc.\ Natl.\ Acad.\ Sci.\ \textbf{119}, e2203533119 (2022);
arXiv:2202.01244.

%% --- Jordan--Wigner and fermion-to-qubit mappings ---

\bibitem{jordan1928}
P.~Jordan and E.~Wigner,
``\"Uber das Paulische \"Aquivalenzverbot,''
Z.\ Phys.\ \textbf{47}, 631--651 (1928).

\bibitem{bravyi2002}
S.~B.~Bravyi and A.~Yu.~Kitaev,
``Fermionic quantum computation,''
Ann.\ Phys.\ (N.Y.) \textbf{298}, 210--226 (2002).

%% --- OpenFermion ---

\bibitem{mcclean2020}
J.~R.~McClean, N.~C.~Rubin, K.~J.~Sung, \textit{et al.},
``OpenFermion: The electronic structure package for quantum computers,''
Quantum Sci.\ Technol.\ \textbf{5}, 034014 (2020).

%% --- Trotter and simulation ---

\bibitem{childs2021}
A.~M.~Childs, Y.~Su, M.~C.~Tran, N.~Wiebe, and S.~Zhu,
``Theory of Trotter error with commutator scaling,''
Phys.\ Rev.\ X \textbf{11}, 011020 (2021).

\bibitem{campbell2019}
E.~Campbell,
``Random compiler for fast Hamiltonian simulation,''
Phys.\ Rev.\ Lett.\ \textbf{123}, 070503 (2019).

%% --- Hardware-efficient and adaptive ---

\bibitem{tang2021}
H.~L.~Tang, V.~O.~Shkolnikov, G.~S.~Barron, \textit{et al.},
``qubit-ADAPT-VQE: An adaptive algorithm for constructing hardware-efficient
ans\"atze on a quantum processor,''
PRX Quantum \textbf{2}, 020310 (2021).

\bibitem{kandala2017}
A.~Kandala, A.~Mezzacapo, K.~Temme, \textit{et al.},
``Hardware-efficient variational quantum eigensolver for small molecules and
quantum magnets,''
Nature \textbf{549}, 242--246 (2017).

%% --- Standard quantum chemistry references ---

\bibitem{szabo1996}
A.~Szabo and N.~S.~Ostlund,
\textit{Modern Quantum Chemistry: Introduction to Advanced Electronic
Structure Theory}
(Dover, Mineola, NY, 1996).

\bibitem{helgaker2000}
T.~Helgaker, P.~J{\o}rgensen, and J.~Olsen,
\textit{Molecular Electronic-Structure Theory}
(Wiley, Chichester, 2000).

%% --- Gaussian basis sets ---

\bibitem{hehre1969}
W.~J.~Hehre, R.~F.~Stewart, and J.~A.~Pople,
``Self-consistent molecular-orbital methods.\ I.\ Use of Gaussian
expansions of Slater-type atomic orbitals,''
J.\ Chem.\ Phys.\ \textbf{51}, 2657--2664 (1969).

%% --- Gaunt integrals ---

\bibitem{gaunt1929}
J.~A.~Gaunt,
``The triplets of helium,''
Philos.\ Trans.\ R.\ Soc.\ A \textbf{228}, 151--196 (1929).

%% --- NIST ---

\bibitem{nist_he}
A.~Kramida, Yu.~Ralchenko, J.~Reader, and NIST ASD Team,
\textit{NIST Atomic Spectra Database} (ver.~5.11),
\url{https://physics.nist.gov/asd} (2024).

%% --- GeoVac internal references ---

\bibitem{loutey_paper7}
J.~Loutey,
``The Dimensionless Vacuum: Recovering the Schr\"odinger Equation from Scale-Invariant Graph Topology,''
GeoVac Paper~7 (2026).

\bibitem{loutey_paper13}
J.~Loutey,
``The Hyperspherical Lattice: Two-Electron Atoms as Coupled Channel Graphs,''
GeoVac Paper~13 (2026).

\bibitem{loutey_paper17}
J.~Loutey,
``Composed Natural Geometries for Core-Valence Diatomics: LiH and BeH$^+$ from Discrete Graph Topology,''
GeoVac Paper~17 (2026).

\bibitem{loutey_paper25}
J.~Loutey,
``The Hopf Graph as a Lattice Gauge Structure: A Framework Observation,''
GeoVac Paper~25 (2026).

\bibitem{loutey_paper29}
J.~Loutey,
``Finite-Size Graph-RH for the GeoVac Hopf Graphs: Ihara Zeta, Bound-Crossing, and a Scope Boundary on Selberg-on-Hydrogen,''
GeoVac Paper~29 (2026).

\bibitem{burkat_fitzpatrick2025}
J.~Burkat and N.~Fitzpatrick,
``The Quantum Paldus Transform: Efficient Circuits with Applications,''
arXiv:2506.09151 (2025).

\bibitem{rocca2024}
D.~Rocca, C.~L.~Cortes, J.~F.~Gonthier, P.~J.~Ollitrault, and R.~M.~Parrish,
``Reducing the runtime of fault-tolerant quantum simulations in chemistry
through symmetry-compressed double factorization,''
J.\ Chem.\ Theory Comput.\ \textbf{20}, 4639--4653 (2024).

\bibitem{caesura2025}
A.~E.~Caesura, C.~L.~Cortes, W.~Pol, S.~Sim, M.~Steudtner,
G.-L.~R.~Anselmetti, M.~Degroote, N.~Moll, R.~Santagati, M.~Streif,
and C.~S.~Tautermann,
``Faster quantum chemistry simulations on a quantum computer with
improved tensor factorization and active volume compilation,''
PRX Quantum \textbf{6}, 030337 (2025); arXiv:2501.06165.

\bibitem{trenev2025}
D.~Trenev, P.~J.~Ollitrault, S.~M.~Harwood, T.~P.~Gujarati,
S.~Raman, A.~Mezzacapo, and S.~Mostame,
``Refining resource estimation for the quantum computation of
vibrational molecular spectra through Trotter error analysis,''
Quantum \textbf{9}, 1630 (2025); arXiv:2311.03719.
[Cited for the electronic-structure Gaussian-basis Pauli counts in
Table~5 (Appendix~B; Jordan--Wigner with 2-qubit reduction), used as
the Gaussian baseline in Table~\ref{tab:gaussian_published}; the
$Q^{3.92}$/$Q^{4.25}$ LiH/H$_2$O exponents and the
$Q^{3.9}$--$Q^{4.3}$ span they bracket are our own log-log fit of
those published counts.]

\bibitem{loutey_paper19}
J.~Loutey,
``Toward PK-Free Molecular Hamiltonians: Two-Center Sturmian
Integrals and Coupled Composition,''
GeoVac Paper~19 (2026).

% ----------------------------------------------------------------------
% References added 2026-06-04 (C3 follow-on closure); citations resolved
% 2026-06-28 by literature search:
%   ChildsBerry          -> Berry, Childs & Kothari, FOCS 2015 (arXiv:1501.01715)
%   MartinezYRomero2004  -> arXiv:physics/0402061 (2004)
%   Sunaga2025           -> RE-KEYED to Chawla2024 (arXiv:2406.04992 =
%                           Phys. Rev. A 111, 022817 (2025)); the RaH-18q
%                           benchmark was mis-attributed to Sunaga, then
%                           mis-keyed to Swain2022 (arXiv:2211.06907, a
%                           different 12q-active-space dipole paper with no
%                           RaH-18q cell). Corrected to Chawla et al. 2026-06-28.
%   BJL                  -> WITHDRAWN (unverifiable; cite removed, prose reworded)
% Dyall confirmed from Paper 31; Reiher 2017 / Szmytkowski 2007 well-known.
% ----------------------------------------------------------------------

\bibitem{Dyall}
K.~G.~Dyall and K.~F\ae gri,
\textit{Introduction to Relativistic Quantum Chemistry}
(Oxford University Press, Oxford, 2007).

\bibitem{reiher2017}
M.~Reiher, N.~Wiebe, K.~M.~Svore, D.~Wecker, and M.~Troyer,
``Elucidating reaction mechanisms on quantum computers,''
Proc.\ Natl.\ Acad.\ Sci.\ USA \textbf{114}, 7555 (2017);
arXiv:1605.03590.

\bibitem{Szmytkowski2007}
R.~Szmytkowski,
``Recurrence and differential relations for spherical spinors,''
J.~Math.\ Chem.\ \textbf{42}, 397 (2007).

\bibitem{BreitPauli}
H.~A.~Bethe and E.~E.~Salpeter,
\textit{Quantum Mechanics of One- and Two-Electron Atoms}
(Springer, Berlin, 1957); the Breit--Pauli Hamiltonian and the
spin-orbit, spin-spin, and spin-other-orbit operators are treated
canonically in Chs.~38--39.

\bibitem{ChildsBerry}
D.~W.~Berry, A.~M.~Childs, and R.~Kothari,
``Hamiltonian simulation with nearly optimal dependence on all
parameters,'' in Proc.\ 56th IEEE FOCS (2015), pp.~792--809;
arXiv:1501.01715.

\bibitem{MartinezYRomero2004}
R.~P.~Mart\'inez-y-Romero, H.~N.~N\'u\~nez-Y\'epez, and
A.~L.~Salas-Brito,
``Relativistic recursion relations for transition matrix elements,''
arXiv:physics/0402061 (2004).

\bibitem{Chawla2024}
P.~Chawla, V.~S.~Prasannaa, K.~R.~Swain, K.~Sugisaki, and B.~P.~Das,
``Relativistic VQE calculations of molecular electric dipole moments on
trapped ion quantum hardware,'' Phys.\ Rev.\ A \textbf{111}, 022817 (2025);
arXiv:2406.04992.
[RaH 18-qubit active space: \gvq{chawla_rah}{12{,}556} relativistic / 2{,}740
non-relativistic Pauli terms (47{,}099 / 4{,}249 two-electron integrals).]

\bibitem{christlmaier2023xtc}
E.~M.~Christlmaier, T.~Schraivogel, P.~L\'opez R\'ios, A.~Alavi, and
D.~Kats, ``xTC: An efficient treatment of three-body interactions in
transcorrelated methods,'' J.\ Chem.\ Phys.\ \textbf{159}, 014113 (2023).

\bibitem{lowsu2024qeve}
G.~H.~Low and Y.~Su, ``Quantum eigenvalue processing,'' SIAM
J.\ Comput.\ \textbf{55}, 135 (2026); arXiv:2401.06240 (2024).

\bibitem{uvarov2025qeve_tc}
A.~Uvarov and A.~F.~Izmaylov, ``Accuracy and resource advantages of
quantum eigenvalue estimation with non-Hermitian transcorrelated
electronic Hamiltonians,'' arXiv:2511.21867 (2025).

\end{thebibliography}

\end{document}
